\providecommand{\safeop}[1]{\mathsf{#1}}
\clearpage
\section*{Integrated Delta Layer: Robust Local Chart}
\addcontentsline{toc}{section}{Integrated Delta Layer: Robust Local Chart}
% --- local macro refresh for r50 ---
\providecommand{\HC}{}
\renewcommand{\HC}{\mathrm{HC}}
\providecommand{\MM}{}
\renewcommand{\MM}{\mathrm{MM}}
\providecommand{\OP}{}
\renewcommand{\OP}{\mathrm{OP}}
\providecommand{\rout}{}
\renewcommand{\rout}{\mathrm{rout}}
\providecommand{\unr}{}
\renewcommand{\unr}{\mathrm{unr}}
\providecommand{\near}{}
\renewcommand{\near}{\mathrm{near}}
\providecommand{\asym}{}
\renewcommand{\asym}{\mathrm{asym}}
\providecommand{\certred}{}
\renewcommand{\certred}{\mathrm{cert\_red}_{\gamma,\asym}}
\providecommand{\closecand}{}
\renewcommand{\closecand}{\mathrm{close\_cand}_{\gamma,\asym}}
\providecommand{\resproj}{}
\renewcommand{\resproj}{\mathrm{res\_proj}_{\gamma,\asym}}
\providecommand{\compmargin}{}
\renewcommand{\compmargin}{\mathrm{comp\_margin}_{\gamma,\asym}}
\providecommand{\debtrout}{}
\renewcommand{\debtrout}{\mathrm{debt\_rout}_{\gamma,\asym}}
% --- begin integrated body for r50 ---



\begin{abstract}
This delta layer continues the local proof chain photon-frame \(\HC\)-collar contraction line. The idealized non-ideal sample already produced a local photon-frame near-kernel closure candidate. The present layer tests whether that candidate survives a denser asymmetric chart, non-uniform curvature/readout weights, and explicitly nonzero routed \(\HC/\MM\) absorption. The result is still conservative: it is a local OP5 closure-candidate certificate for the declared chart only, not a global OP5 closure theorem.
\end{abstract}

\section{Position of the probe}

The chart certified a finite near-kernel photon-frame sample by checking
\[
 \widehat\lambda_{\gamma,K}^{\mathrm{loc},(2)}+\eta_{\mathrm{loc}}<1,
 \qquad L_{\unr}^{\mathrm{loc}}=0.
\]
That test was useful but still too symmetric: the sample table used a short four-point chart and encoded the remaining perturbation through a relatively coarse margin. The test strengthens the diagnostic in three ways:
\begin{enumerate}
 \item the photon-frame chart is denser and intentionally asymmetric;
 \item each sample point carries a non-uniform curvature/readout weight;
 \item the \(\HC/\MM\) absorption register is allowed to be nonzero, but only as routed debt-free bookkeeping.
\end{enumerate}
The intended question is therefore not whether OP5 is globally closed, but whether the local \(\HC\)-cone reading remains contraction-compatible after the first robustification of the sample.

\section{Asymmetric weighted photon-frame chart}

\begin{definition}[asymmetric photon-frame HC-absorption chart]
An asymmetric photon-frame chart is a finite packet
\[
\mathcal S_{\gamma,K}^{\asym,(2)}
=
\{(b_i,q_i,a_i,e_i,w_i,L_i^{\rout},L_i^{\unr})\}_{i=1}^{8},
\]
where:
\begin{itemize}
 \item \(b_i>0\) is the unshifted photon-frame radial readout size;
 \item \(q_i\) is the normalized second-shift radius ratio before explicit absorption;
 \item \(a_i\) is the normalized routed \(\HC/\MM\)-absorption contribution;
 \item \(e_i\) is the normalized finite chart error/tolerance contribution;
 \item \(w_i>0\) is a declared non-uniform readout weight;
 \item \(L_i^{\rout}\) records routed leakage/absorption bookkeeping;
 \item \(L_i^{\unr}\) records leakage that is not routed through an admissible \(\HC/\MM\) register.
\end{itemize}
The pointwise weighted contraction ratio is bounded by
\[
\lambda_i^{\mathrm{loc},(2)}
\le q_i+a_i+e_i.
\]
The chart estimator and routed/unrouted guards are
\[
\widehat\lambda_{\gamma,K}^{\mathrm{loc},(2)}
:=\max_{1\le i\le 8}\lambda_i^{\mathrm{loc},(2)},
\qquad
L_{\unr}^{\mathrm{loc}}:=\sum_{i=1}^{8}L_i^{\unr},
\qquad
B_{\HC/\MM}^{\mathrm{loc}}:=\mathrm{unpaid}(L^{\rout}).
\]
Here \(B_{\HC/\MM}^{\mathrm{loc}}=0\) means that nonzero absorption has been declared and routed without leaving residual absorption debt.
\end{definition}

\begin{definition}[local closure-candidate classifier]
The asymmetric photon-frame branch is classified by
\[
\mathrm{Class}_{\gamma,K}^{\mathrm{loc}}
=
\begin{cases}
\certred, & \widehat\lambda_{\gamma,K}^{\mathrm{loc},(2)}+\eta_{\mathrm{loc}}<1,
 \ L_{\unr}^{\mathrm{loc}}=0,
 \ B_{\HC/\MM}^{\mathrm{loc}}=0,\\[2mm]
\compmargin, & \widehat\lambda_{\gamma,K}^{\mathrm{loc},(2)}+\eta_{\mathrm{loc}}<1,
 \ L_{\unr}^{\mathrm{loc}}=0,
 \ B_{\HC/\MM}^{\mathrm{loc}}>0,\\[2mm]
\debtrout, & \widehat\lambda_{\gamma,K}^{\mathrm{loc},(2)}+\eta_{\mathrm{loc}}<1,
 \ L_{\unr}^{\mathrm{loc}}>0\text{ but routeable},\\[2mm]
\resproj, & \widehat\lambda_{\gamma,K}^{\mathrm{loc},(2)}+\eta_{\mathrm{loc}}\ge 1
 \text{ or leakage is not routeable.}
\end{cases}
\]
The certificate \(\certred\) is local to the declared chart and is read as \(\closecand\), not as a global OP5 theorem.
\end{definition}

\section{Explicit sample table}

The following sample is deliberately not a uniform halving table. The base radii are uneven, the weights are non-uniform, and the routed \(\HC/\MM\)-absorption contribution is nonzero in every row.

\begin{longtable}{>{\centering\arraybackslash}p{0.7cm} >{\centering\arraybackslash}p{1.6cm} >{\centering\arraybackslash}p{1.5cm} >{\centering\arraybackslash}p{1.5cm} >{\centering\arraybackslash}p{1.5cm} >{\centering\arraybackslash}p{1.5cm} >{\centering\arraybackslash}p{1.8cm} >{\centering\arraybackslash}p{1.4cm}}
\toprule
\(i\) & \(b_i\) & \(w_i\) & \(q_i\) & \(a_i\) & \(e_i\) & \(q_i+a_i+e_i\) & \(L_i^{\unr}\) \\
\midrule
\endfirsthead
\toprule
\(i\) & \(b_i\) & \(w_i\) & \(q_i\) & \(a_i\) & \(e_i\) & \(q_i+a_i+e_i\) & \(L_i^{\unr}\) \\
\midrule
\endhead
1 & \(1\) & \(1\) & \(1/2\) & \(1/64\) & \(1/64\) & \(17/32\) & \(0\) \\
2 & \(29/32\) & \(33/32\) & \(17/32\) & \(1/64\) & \(1/32\) & \(37/64\) & \(0\) \\
3 & \(23/32\) & \(17/16\) & \(9/16\) & \(1/32\) & \(1/64\) & \(39/64\) & \(0\) \\
4 & \(17/32\) & \(19/16\) & \(19/32\) & \(1/32\) & \(1/32\) & \(21/32\) & \(0\) \\
5 & \(13/32\) & \(5/4\) & \(5/8\) & \(1/64\) & \(1/64\) & \(21/32\) & \(0\) \\
6 & \(3/8\) & \(21/16\) & \(21/32\) & \(1/64\) & \(1/64\) & \(11/16\) & \(0\) \\
7 & \(9/32\) & \(11/8\) & \(5/8\) & \(1/32\) & \(1/32\) & \(11/16\) & \(0\) \\
8 & \(5/32\) & \(3/2\) & \(19/32\) & \(1/64\) & \(1/32\) & \(41/64\) & \(0\) \\
\bottomrule
\end{longtable}

All rows carry routed absorption \(a_i>0\). We declare the routed register as paid,
\[
L_{\rout}^{\mathrm{loc}}>0,
\qquad
B_{\HC/\MM}^{\mathrm{loc}}=0,
\qquad
L_{\unr}^{\mathrm{loc}}=0.
\]
The maximum row value is
\[
\widehat\lambda_{\gamma,K}^{\mathrm{loc},(2)}
=\max_i(q_i+a_i+e_i)
=\frac{11}{16}.
\]
With the asymmetric chart margin
\[
\eta_{\mathrm{loc}}=\frac{9}{64},
\]
one obtains
\[
\widehat\lambda_{\gamma,K}^{\mathrm{loc},(2)}+\eta_{\mathrm{loc}}
=\frac{11}{16}+\frac{9}{64}
=\frac{53}{64}<1.
\]

\begin{proposition}[asymmetric HC-absorption robustness certificate]
For the explicit table above, the asymmetric photon-frame branch receives the local classifier
\[
\mathrm{Class}_{\gamma,K}^{\mathrm{loc}}=\certred=\closecand.
\]
Equivalently, the local closure-candidate status survives the first denser asymmetric weighted chart with nonzero routed \(\HC/\MM\) absorption.
\end{proposition}

\begin{proof}
The estimator is the maximum of the row bounds in the explicit table. The largest row value is \(11/16\), attained in rows 6 and 7. The declared margin is \(9/64\), hence
\[
\widehat\lambda_{\gamma,K}^{\mathrm{loc},(2)}+\eta_{\mathrm{loc}}=53/64<1.
\]
The table has zero unrouted leakage in every row, so \(L_{\unr}^{\mathrm{loc}}=0\). The absorption register is nonzero, because \(a_i>0\) in every row, but it is explicitly routed and declared debt-free, so \(B_{\HC/\MM}^{\mathrm{loc}}=0\). The first branch of the classifier therefore applies. This proves the stated local certificate.
\end{proof}

\section{Field-equation bookkeeping update}

The branch adds the asymmetric photon-frame contraction register
\[
\mathcal R_{\gamma,\asym}^{(2)}
:=
\left(
\widehat\lambda_{\gamma,K}^{\mathrm{loc},(2)},
\eta_{\mathrm{loc}},
L_{\rout}^{\mathrm{loc}},
L_{\unr}^{\mathrm{loc}},
B_{\HC/\MM}^{\mathrm{loc}},
\{w_i\}_{i=1}^{8}
\right)
\]
to the OP5-facing source-return bookkeeping. The readout-side implication is
\[
\mathcal R_{\gamma,\asym}^{(2)}
\vdash
\left[
\widehat\lambda+\eta<1,
\ L_{\unr}=0,
\ B_{\HC/\MM}=0
\right]
\Longrightarrow
\closecand.
\]
Thus the \(\HC\)-cone interpretation is strengthened from an ideal or near-ideal photon-frame sample to a first robust asymmetric weighted sample. The correct structural reading is that OP5 remains more plausibly an \(\HC\)-cone/pass-count debt problem than an independent terminal endpoint problem, provided the remaining chart families keep satisfying the same contraction and routing guards.

\section{Conservativity guards}

The following guards are part of the layer.
\begin{enumerate}
 \item No global OP5 closure is claimed.
 \item The certificate is local to the declared finite photon-frame chart.
 \item Nonzero \(\HC/\MM\) absorption is not deletion of residuals; it is paid routed bookkeeping.
 \item The chart does not assert a literal Euclidean photon trajectory.
 \item The result does not close OP1, OP4, QCD/confinement, CKM/PMNS, or the quark-mass spectrum.
 \item Promotion beyond \(\closecand\) still requires either a covering argument over admissible photon-frame charts or a finite atlas theorem with stable contraction margins.
\end{enumerate}

\section{Next recommended step}

The highest-leverage next step is an atlas-style finite cover test: define a small family of admissible asymmetric photon-frame charts
\[
\mathfrak A_{\gamma}^{(2)}=
\{\mathcal S_{\gamma,K,j}^{\asym,(2)}\}_{j=1}^{M}
\]
and require a uniform bound
\[
\sup_{j\le M}\left(
\widehat\lambda_{\gamma,K,j}^{(2)}+\eta_j\right)<1,
\qquad
\sum_{j\le M}L_{\unr,j}=0,
\qquad
\sum_{j\le M}B_{\HC/\MM,j}=0.
\]
If such a finite-atlas certificate is obtained, the OP5=\(\HC\)-cone hypothesis can be promoted from a single-chart local closure candidate to a finite-atlas closure candidate.


% --- end integrated body for r50 ---

\clearpage
\section*{Integrated Delta Layer: Finite Atlas}
\addcontentsline{toc}{section}{Integrated Delta Layer: Finite Atlas}
% --- local macro refresh for r51 ---
\providecommand{\HC}{}
\renewcommand{\HC}{\mathrm{HC}}
\providecommand{\MM}{}
\renewcommand{\MM}{\mathrm{MM}}
\providecommand{\OP}{}
\renewcommand{\OP}{\mathrm{OP}}
\providecommand{\rout}{}
\renewcommand{\rout}{\mathrm{rout}}
\providecommand{\unr}{}
\renewcommand{\unr}{\mathrm{unr}}
\providecommand{\asym}{}
\renewcommand{\asym}{\mathrm{asym}}
\providecommand{\Atlas}{}
\renewcommand{\Atlas}{\mathfrak A}
\providecommand{\closecand}{}
\renewcommand{\closecand}{\mathrm{close\_cand}}
\providecommand{\atlascand}{}
\renewcommand{\atlascand}{\mathrm{atlas\_close\_cand}}
\providecommand{\certred}{}
\renewcommand{\certred}{\mathrm{cert\_red}}
\providecommand{\certcomp}{}
\renewcommand{\certcomp}{\mathrm{cert\_comp}}
\providecommand{\resproj}{}
\renewcommand{\resproj}{\mathrm{res\_proj}}
\providecommand{\ConeObs}{}
\renewcommand{\ConeObs}{D_{\OP5=\HC\text{-}\mathrm{cone}}}
% --- begin integrated body for r51 ---



\begin{abstract}
This delta layer upgrades the local proof chain photon-frame \(\HC\)-collar contraction probes from single-chart diagnostics to a finite-atlas closure-candidate test. The layer also states, explicitly and conservatively, what would still have to be proved in order to close OP5 under the hypothesis that OP5 is the \(\HC\)-cone/reflector surface of the source-return route rather than an independent endpoint obstruction. The result is an atlas-level closure candidate, not a final OP5 theorem.
\end{abstract}

\section{Position after}

The layer showed that the local photon-frame candidate survives a denser asymmetric chart with nonzero but debt-free routed \(\HC/\MM\) absorption:
\[
\widehat\lambda_{\gamma,K}^{\mathrm{loc},(2)}=\frac{11}{16},
\qquad
\eta_{\mathrm{loc}}=\frac{9}{64},
\qquad
\widehat\lambda_{\gamma,K}^{\mathrm{loc},(2)}+\eta_{\mathrm{loc}}=\frac{53}{64}<1.
\]
This is strong enough to justify an atlas test but not strong enough to close OP5. A single chart only says that one declared readout neighbourhood is compatible with the \(\HC\)-cone contraction picture. OP5 closure requires that the admissible OP5 branch is exhausted by such neighbourhoods, and that the transition logic between them does not reintroduce an endpoint obstruction.

\section{Finite photon-frame atlas}

\begin{definition}[finite photon-frame HC-cone atlas]
An finite photon-frame atlas for the OP5-facing branch is a finite family
\[
\Atlas_{\gamma}^{(2)}=
\{\mathcal S_{\gamma,K,j}^{(2)}\}_{j=1}^{M}
\]
of admissible photon-frame \(\HC\)-collar charts. Each chart carries the data
\[
\mathcal S_{\gamma,K,j}^{(2)}
\leadsto
\left(
\widehat\lambda_j,
\eta_j,
L_{\unr,j},
B_{\HC/\MM,j},
\tau_j
\right),
\]
where \(\widehat\lambda_j\) is the finite contraction estimator, \(\eta_j\) is the chart margin, \(L_{\unr,j}\) is the unrouted leakage register, \(B_{\HC/\MM,j}\) is the unpaid \(\HC/\MM\)-absorption debt, and \(\tau_j\) records the transition-map debt between overlapping charts.
\end{definition}

\begin{definition}[Atlas classifier]
The atlas receives the classifier by the following finite rule.
\[
\mathrm{Class}(\Atlas_{\gamma}^{(2)})=\atlascand
\]
exactly if
\[
\sup_{j\le M}(\widehat\lambda_j+\eta_j)<1,\qquad
\sum_{j\le M}L_{\unr,j}=0,\qquad
\sum_{j\le M}B_{\HC/\MM,j}=0,\qquad
\sum_{j\le M}\tau_j=0.
\]
It receives the compatible-open classifier \(\certcomp\) if the uniform contraction inequality holds but some leakage, absorption, or transition register is nonzero while still routeable. It receives the residual-projection classifier \(\resproj\) if the uniform contraction bound fails or if any register is not routeable. Thus an atlas is not merely a list of successful sample tables. It also has to verify that no obstruction appears on overlaps or transitions.
\end{definition}

\section{Explicit finite-atlas sample}

The following five-chart atlas is a compact first test. It collects the already motivated ideal, near-kernel, non-ideal, asymmetric, and tilted/edge regimes into one finite certificate table.

{\small\setlength{\tabcolsep}{3pt}
\begin{longtable}{>{\raggedright\arraybackslash}p{3.0cm} >{\centering\arraybackslash}p{1.75cm} >{\centering\arraybackslash}p{1.55cm} >{\centering\arraybackslash}p{1.95cm} >{\centering\arraybackslash}p{1.55cm} >{\centering\arraybackslash}p{2.0cm}}
\toprule
Chart & \(\widehat\lambda_j\) & \(\eta_j\) & \(\widehat\lambda_j+\eta_j\) & \(L_{\unr,j}\) & \(B_{\HC/\MM,j}+\tau_j\) \\
\midrule
\endfirsthead
\toprule
Chart & \(\widehat\lambda_j\) & \(\eta_j\) & \(\widehat\lambda_j+\eta_j\) & \(L_{\unr,j}\) & \(B_{\HC/\MM,j}+\tau_j\) \\
\midrule
\endhead
central ideal collar & \(1/2\) & \(1/8\) & \(5/8\) & \(0\) & \(0\) \\
near-kernel symmetric & \(21/32\) & \(1/8\) & \(25/32\) & \(0\) & \(0\) \\
non-ideal sample & \(21/32\) & \(1/8\) & \(25/32\) & \(0\) & \(0\) \\
asymmetric HC-absorption & \(11/16\) & \(9/64\) & \(53/64\) & \(0\) & \(0\) \\
tilted edge collar & \(45/64\) & \(7/64\) & \(13/16\) & \(0\) & \(0\) \\
\bottomrule
\end{longtable}
}

The atlas maximum is therefore
\[
\Lambda_{\Atlas}^{\mathrm{loc}}
:=
\sup_{j\le 5}(\widehat\lambda_j+\eta_j)
=
\frac{53}{64}<1.
\]
All explicit leakage, absorption-debt, and transition-debt registers are declared zero at the atlas level:
\[
\sum_{j\le 5}L_{\unr,j}=0,
\qquad
\sum_{j\le 5}B_{\HC/\MM,j}=0,
\qquad
\sum_{j\le 5}\tau_j=0.
\]

\begin{proposition}[finite-atlas closure-candidate certificate]
The explicit five-chart photon-frame atlas above receives the classifier
\[
\mathrm{Class}(\Atlas_{\gamma}^{(2)})=
\atlascand.
\]
It promotes the single-chart result to a finite-atlas OP5 closure candidate.
\end{proposition}

\begin{proof}
The largest value in the third numerical column is \(53/64\), which is strictly smaller than one. The declared atlas registers satisfy
\[
\sum_{j\le5}L_{\unr,j}=0,
\qquad
\sum_{j\le5}B_{\HC/\MM,j}=0,
\qquad
\sum_{j\le5}\tau_j=0.
\]
The first branch of the atlas classifier applies. Hence the finite atlas receives the atlas closure-candidate status.
\end{proof}

\section{What would be needed to close OP5}

The atlas is still a candidate because it assumes, rather than proves, that the five declared chart types cover all admissible OP5-facing photon-frame readouts. Under the working hypothesis
\[
\OP5 = \HC\text{-cone/reflector surface of the source-return route},
\]
OP5 would be closed by proving the following theorem.

\begin{theorem}[Conditional OP5 closure theorem to be proved]
Assume that the OP5-facing source-return route factorizes through the \(\HC\)-cone reflector,
\[
R_{\OP}^{\HC\text{-}\mathrm{cone}}
= T_{\Gamma}^{\vee}\circ J_{\HC}^{\mathrm{cone}}\circ T_{\Gamma},
\]
and that its cone-factorization obstruction vanishes globally,
\[
\ConeObs^{\chi,(1)}=0.
\]
Then OP5 closes if there exists a finite photon-frame atlas \(\Atlas_{\gamma}^{(2)}\) satisfying all of the following conditions.

\begin{enumerate}
 \item \textbf{Coverage.} Every admissible OP5-facing readout state is contained in at least one chart of \(\Atlas_{\gamma}^{(2)}\).
 \item \textbf{Uniform contraction.} There is a constant \(\Lambda<1\) such that
 \[
 \widehat\lambda_j+\eta_j\le\Lambda<1
 \qquad\text{for every chart }j.
 \]
 \item \textbf{Transition stability.} On every nonempty chart overlap, the transition map preserves the contraction class and creates no unpaid transition debt:
 \[
 \tau_{jk}=0
 \qquad\text{for all admissible overlaps }j\cap k.
 \]
 \item \textbf{No unrouted leakage.} Every leakage term is either zero or routed through an admissible \(\HC/\MM\) register, and the total unrouted leakage vanishes:
 \[
 L_{\unr}^{\Atlas}=0.
 \]
 \item \textbf{No unpaid absorption debt.} Routed \(\HC/\MM\) absorption may be nonzero, but all such absorption must be paid inside the route bookkeeping:
 \[
 B_{\HC/\MM}^{\Atlas}=0.
 \]
 \item \textbf{Periapsis/no-through-collapse guard.} The active \(\HC\)-collar forces every admissible source-return branch to have a collar-side periapsis rather than a through-collapse endpoint.
 \item \textbf{Pass-count consistency.} The OP5 pass count is exactly the source-return pass count seen through the second Hilbert-shift/\(\HC\)-cone frame, not an additional independent terminal transport pass.
 \item \textbf{Statement identification.} The formal statement of OP5 is shown to be equivalent to the absence of the above cone-factorization, leakage, debt, and transition obstructions.
\end{enumerate}

If all eight items hold, then OP5 is not an independent residual endpoint problem. It is discharged as the \(\HC\)-cone source-return contraction theorem.
\end{theorem}

\begin{remark}[Why does not yet close OP5]
The table verifies uniform contraction for a finite declared atlas and records zero debt in the declared registers. It does not yet prove the coverage item, the global vanishing of \(\ConeObs^{\chi,(1)}\), or the equivalence between the formal OP5 statement and the \(\HC\)-cone source-return theorem. Those are the remaining proof burdens.
\end{remark}

\section{Practical proof route after}

The shortest conservative route from to a genuine OP5 closure proof is:
\begin{enumerate}
 \item write the admissible OP5 readout space \(\mathcal X_{\OP5}^{\mathrm{adm}}\);
 \item prove that the five chart types form a finite cover of \(\mathcal X_{\OP5}^{\mathrm{adm}}\);
 \item prove that all chart overlaps have zero transition debt;
 \item prove that \(\ConeObs^{\chi,(1)}=0\) under the \(\HC\)-cone factorization hypothesis;
 \item prove that no admissible branch crosses the collapse core under an active \(\HC\)-collar, i.e. every branch has a periapsis/readout-turning point;
 \item identify OP5's original obstruction statement with the atlas obstruction package.
\end{enumerate}
If these six tasks are completed, the numerical contraction certificates become supporting evidence for a structural theorem rather than isolated probes.

\section{Conservativity guards}

The layer uses the following guards.
\begin{enumerate}
 \item The result is a finite-atlas closure candidate, not a final OP5 closure.
 \item The five-chart cover is declared as a test atlas; its coverage of all admissible OP5 states remains to be proved.
 \item The \(\HC\)-cone interpretation remains conditional on vanishing cone-factorization obstruction.
 \item Nonzero \(\HC/\MM\) absorption is permitted only as explicit routed bookkeeping, never as residual deletion.
 \item The result does not imply OP1, OP4, QCD/confinement, CKM/PMNS, or quark-mass spectrum closure.
\end{enumerate}

\section{Next recommended step}

The highest-leverage next step is to define the admissible OP5 readout space
\[
\mathcal X_{\OP5}^{\mathrm{adm}}
\]
and prove a finite-cover lemma
\[
\mathcal X_{\OP5}^{\mathrm{adm}}
\subseteq
\bigcup_{j=1}^{5}\mathcal S_{\gamma,K,j}^{(2)}.
\]
That lemma is the first real bridge from atlas diagnostics to OP5 closure.


% --- end integrated body for r51 ---

\clearpage
\section*{Integrated Delta Layer: Tunnel Menu}
\addcontentsline{toc}{section}{Integrated Delta Layer: Tunnel Menu}
% --- local macro refresh for r52 ---
\providecommand{\HC}{}
\renewcommand{\HC}{\mathrm{HC}}
\providecommand{\MM}{}
\renewcommand{\MM}{\mathrm{MM}}
\providecommand{\OP}{}
\renewcommand{\OP}{\mathrm{OP}}
\providecommand{\adm}{}
\renewcommand{\adm}{\mathrm{adm}}
\providecommand{\tun}{}
\renewcommand{\tun}{\mathrm{tun}}
\providecommand{\cand}{}
\renewcommand{\cand}{\mathrm{cand}}
\providecommand{\unr}{}
\renewcommand{\unr}{\mathrm{unr}}
\providecommand{\rout}{}
\renewcommand{\rout}{\mathrm{rout}}
\providecommand{\Halo}{}
\renewcommand{\Halo}{\mathsf{Halo}}
\providecommand{\Exit}{}
\renewcommand{\Exit}{\mathsf{Exit}}
\providecommand{\Trap}{}
\renewcommand{\Trap}{\mathsf{Trap}}
\providecommand{\Gate}{}
\renewcommand{\Gate}{\mathsf{Gate}}
\providecommand{\Atlas}{}
\renewcommand{\Atlas}{\mathfrak A}
\providecommand{\epsr}{}
\renewcommand{\epsr}{\varepsilon_r}
\providecommand{\epsh}{}
\renewcommand{\epsh}{\varepsilon_h}
\providecommand{\Xop}{}
\renewcommand{\Xop}{\mathcal X_{\OP5}^{\adm}}
\providecommand{\Ttun}{}
\renewcommand{\Ttun}{\mathcal T_{\tun}}
\providecommand{\Dhc}{}
\renewcommand{\Dhc}{D_{\HC\text{-}\mathrm{cone}}}
\providecommand{\ConeObs}{}
\renewcommand{\ConeObs}{D_{\OP5=\HC\text{-}\mathrm{cone}}}
% --- begin integrated body for r52 ---



\begin{abstract}
This delta layer inserts the nonstandard-halo step between the robust photon-frame contraction probe and the finite-atlas closure-candidate logic. Its purpose is to reduce the apparent combinatorial overload of OP5: instead of testing all admissible source-return routes as free trajectories, the OP5-facing branch is first trapped in a double-epsilon halo around the collapse kernel. Exit from this halo is then possible only through explicitly gated tunnel points. The result is conservative: OP5 is not closed here. The proof burden is reduced from an all-routes problem to a finite or finitely classifiable tunnel-point problem, provided the stated halo-capture, halo-invariance, and tunnel-gating hypotheses are later proved.
\end{abstract}

\section{Version position}

The immediate pre-line contains two relevant ingredients. First, produced a robust asymmetric photon-frame local certificate with
\[
\widehat\lambda_{\gamma,K}^{\mathrm{loc},(2)}=\frac{11}{16},
\qquad
\eta_{\mathrm{loc}}=\frac{9}{64},
\qquad
\widehat\lambda_{\gamma,K}^{\mathrm{loc},(2)}+\eta_{\mathrm{loc}}=\frac{53}{64}<1.
\]
Second, an finite-atlas draft upgrades this local logic to a declared finite atlas, but such an atlas still needs an independent explanation of why the admissible OP5 search does not branch into uncontrolled route families. The present layer supplies that missing structural filter: a double-epsilon halo channel with tunnel-point reduction.

\section{Double-epsilon halo algebra}

\begin{definition}[Two-scale halo readout algebra]
Let \(\mathcal A\) be the OP5-facing readout algebra and let \(K\) denote the collapse kernel. A two-scale halo readout is a triple
\[
(x;\rho_r(x),\rho_h(x))\in \mathcal A\times\mathbb R_{\ge0}\times\mathbb R_{\ge0},
\]
where \(\rho_r(x)\) is the radial kernel-distance readout and \(\rho_h(x)\) is the \(\HC\)-cone residual readout. The two infinitesimal windows are denoted
\[
\epsr>0,
\qquad
\epsh>0.
\]
The corresponding double-epsilon halo is
\[
\Halo_{\epsr,\epsh}(K)
:=
\{x\in\Xop:\rho_r(x)\le\epsr\ \text{and}\ \rho_h(x)\le\epsh\}.
\]
The first inequality is radial closeness to the collapse kernel. The second inequality is closeness to the \(\HC\)-cone/reflector normal form.
\end{definition}

\begin{remark}[Why two epsilon channels are needed]
A one-parameter halo around \(K\) only records radial nearness. OP5 is not merely a radial-distance problem; the relevant question is whether the source-return route is already captured by the active \(\HC\)-cone collar. Therefore the halo must carry two channels: one for kernel nearness and one for \(\HC\)-normality. In this sense \(\epsr\) controls ``near the kernel'', while \(\epsh\) controls ``near the \(\HC\)-cone route''.
\end{remark}

\begin{definition}[Halo boundary and exit side]
The halo boundary is the union of the radial and \(\HC\)-residual boundary faces:
\[
\partial\Halo_{\epsr,\epsh}(K)
=
\{\rho_r=\epsr,\rho_h\le\epsh\}
\cup
\{\rho_r\le\epsr,\rho_h=\epsh\}.
\]
An exit candidate is a point
\[
p\in\partial\Halo_{\epsr,\epsh}(K)
\]
for which the source-return update points outward in at least one of the two halo channels before applying the tunnel selector.
\end{definition}

\section{Tunnel points}

\begin{definition}[Tunnel candidate]
A point \(p\in\partial\Halo_{\epsr,\epsh}(K)\) is a tunnel candidate if it satisfies the four tests
\[
\begin{aligned}
&\text{(T1)}\quad p\in\partial\Halo_{\epsr,\epsh}(K),\\
&\text{(T2)}\quad \Dhc(p)\le\epsh,\\
&\text{(T3)}\quad \chi_{\tun}(p)=1,\\
&\text{(T4)}\quad \operatorname{rank}J_{\Exit}(p)<\operatorname{rank}_{\max}.
\end{aligned}
\]
Here \(\chi_{\tun}\) is the tunnel selector and \(J_{\Exit}\) is the exit Jacobian for the paired radial/\(\HC\)-residual channel. The resulting candidate set is
\[
\mathcal T_{\tun}^{\cand}(\epsr,\epsh)
:=
\{p\in\partial\Halo_{\epsr,\epsh}(K):\text{(T1)--(T4) hold}\}.
\]
\end{definition}

\begin{definition}[Resolved tunnel point]
A tunnel candidate \(p\in\mathcal T_{\tun}^{\cand}\) is resolved if exactly one of the following outcomes is certified:
\[
\Gate(p)
\in
\{\mathrm{blocked},\mathrm{HC/MM\ absorbed},\mathrm{atlas\ routed},\mathrm{residual\ projected}\}.
\]
The first three outcomes are controlled. The fourth outcome records a remaining obstruction and cannot be used to close OP5.
\end{definition}

\begin{remark}[Tunnel points as the replacement for route explosion]
The intended reduction is not that particles literally escape through geometric holes. The tunnel language is a readout-side compression: a halo-trapped branch can only leave the halo where the radial boundary, the \(\HC\)-cone residual boundary, and the selector/rank conditions simultaneously permit an exit. Thus an apparent continuum of route choices is replaced by a finite or finitely classifiable family of boundary obstructions.
\end{remark}

\section{Capture, invariance, and tunnel reduction}

\begin{definition}[Halo capture]
The OP5-facing route satisfies halo capture if every admissible readout enters the double halo after finitely many \(\HC\)-normalized source-return updates:
\[
\forall x\in\Xop\ \exists n<\infty:
R_{\HC}^{n}(x)\in\Halo_{\epsr,\epsh}(K).
\]
\end{definition}

\begin{definition}[Halo invariance off tunnels]
The double halo is invariant off tunnel candidates if
\[
x\in\Halo_{\epsr,\epsh}(K)\setminus\mathcal T_{\tun}^{\cand}
\quad\Longrightarrow\quad
R_{\HC}(x)\in\Halo_{\epsr,\epsh}(K).
\]
Equivalently, the active \(\HC\)-collar produces a periapsis/turning-point readout before a through-collapse exit can occur, unless the state lies on a declared tunnel candidate.
\end{definition}

\begin{proposition}[tunnel-point reduction]
Assume halo capture and halo invariance off tunnels. Then every admissible OP5-facing branch is either permanently trapped in the double-epsilon halo or reaches a point of \(\mathcal T_{\tun}^{\cand}(\epsr,\epsh)\). Consequently, the OP5 route problem reduces to resolving the tunnel candidate set.
\end{proposition}

\begin{proof}
Let \(x\in\Xop\). By halo capture, there is a finite \(n\) such that \(x_n:=R_{\HC}^{n}(x)\in\Halo_{\epsr,\epsh}(K)\). If no later iterate lies in \(\mathcal T_{\tun}^{\cand}\), halo invariance implies by induction that every later iterate remains in \(\Halo_{\epsr,\epsh}(K)\). If a later iterate fails to remain in the halo, the contrapositive of halo invariance says that the failure can occur only at a tunnel candidate. Hence every branch is trapped or reaches the tunnel candidate set. This proves the reduction.
\end{proof}

\section{Finite tunnel menu}

The first conservative tunnel inventory is not a numerical discovery claim. It is a finite classification scheme for the exits that must be tested.

\begin{definition}[First OP5 tunnel menu]
The first tunnel menu is
\[
\mathfrak T_{\OP5}^{\mathrm{loc}}
=
\{T_r,T_h,T_{rh},T_\chi,T_{\mathrm{rk}},T_{\mathrm{ov}}\},
\]
where:
\begin{center}
\begin{tabular}{>{\raggedright\arraybackslash}p{2.4cm} >{\raggedright\arraybackslash}p{9.4cm}}
\toprule
Class & Meaning \\
\midrule
\(T_r\) & radial-face exit: \(\rho_r=\epsr\), \(\rho_h<\epsh\); \\
\(T_h\) & \(\HC\)-residual-face exit: \(\rho_r<\epsr\), \(\rho_h=\epsh\); \\
\(T_{rh}\) & corner exit: \(\rho_r=\epsr\), \(\rho_h=\epsh\); \\
\(T_\chi\) & selector-switch exit caused by a change in \(\chi_{\tun}\); \\
\(T_{\mathrm{rk}}\) & rank-drop exit caused by \(\operatorname{rank}J_{\Exit}<\operatorname{rank}_{\max}\); \\
\(T_{\mathrm{ov}}\) & atlas-overlap exit caused by transition debt between two photon-frame charts. \\
\bottomrule
\end{tabular}
\end{center}
A tunnel candidate must be assigned to at least one class in \(\mathfrak T_{\OP5}^{\mathrm{loc}}\).
\end{definition}

\begin{proposition}[Finite tunnel classification target]
If the candidate tunnel set admits a finite decomposition
\[
\mathcal T_{\tun}^{\cand}(\epsr,\epsh)
=
\bigsqcup_{\alpha\in\mathfrak T_{\OP5}^{\mathrm{loc}}}
\mathcal T_\alpha
\]
and every nonempty \(\mathcal T_\alpha\) has one of the controlled outcomes
\[
\Gate(\mathcal T_\alpha)
\in
\{\mathrm{blocked},\mathrm{HC/MM\ absorbed},\mathrm{atlas\ routed}\},
\]
then OP5 has no uncontrolled tunnel escape inside the declared double-epsilon halo.
\end{proposition}

\begin{proof}
By the preceding tunnel-point reduction, every non-trapped branch can only exit through \(\mathcal T_{\tun}^{\cand}\). The finite decomposition assigns each such candidate to one of the declared tunnel classes. If every class is blocked, absorbed without unpaid debt, or routed into an atlas chart, then no class contributes an uncontrolled exit. Therefore the double halo has no uncontrolled tunnel escape.
\end{proof}

\section{Connection to the finite atlas}

Once a tunnel point is atlas-routed, it must land in a contraction chart rather than in a free branch. The bridge to the finite-atlas line is therefore:
\[
\mathcal T_{\tun}^{\mathrm{atlas}}
\subseteq
\bigcup_{j=1}^{M}\mathcal S_{\gamma,K,j}^{(2)},
\qquad
\sup_{j\le M}(\widehat\lambda_j+\eta_j)<1,
\]
with the same zero-debt guards
\[
L_{\unr}^{\Atlas}=0,
\qquad
B_{\HC/\MM}^{\Atlas}=0,
\qquad
\tau^{\Atlas}=0.
\]
For typographic clarity, the intended symbol in the preceding display is the chart margin \(\eta_j\). Thus the atlas-routed tunnel branch inherits the atlas criterion exactly when all tunnel exits are covered by charts with a uniform strict contraction margin.

\section{Conditional OP5 closure statement after the halo step}

\begin{theorem}[Conditional OP5 closure via double-epsilon halo]
Assume the working identification
\[
\OP5=\HC\text{-cone source-return reflector problem}
\]
and assume that the cone-factorization obstruction vanishes,
\[
\ConeObs^{\chi,(1)}=0.
\]
Then OP5 closes if the following six conditions are proved.
\begin{enumerate}
 \item \textbf{Halo capture:} every admissible OP5 readout enters \(\Halo_{\epsr,\epsh}(K)\) after finitely many \(\HC\)-normalized source-return updates.
 \item \textbf{Halo invariance off tunnels:} the double halo is invariant outside \(\mathcal T_{\tun}^{\cand}\).
 \item \textbf{Finite tunnel classification:} \(\mathcal T_{\tun}^{\cand}\) is empty, finite, or decomposes into finitely many tunnel classes.
 \item \textbf{Tunnel gating:} each nonempty tunnel class is blocked, debt-free \(\HC/\MM\)-absorbed, or routed into a contraction atlas.
 \item \textbf{No hidden leakage or debt:} all tunnel and atlas registers satisfy
 \[
 L_{\unr}=0,
 \qquad
 B_{\HC/\MM}=0,
 \qquad
 \tau=0.
 \]
 \item \textbf{OP5 statement identification:} the original OP5 obstruction is equivalent to the cone-factorization/tunnel/atlas obstruction package above.
\end{enumerate}
Under these hypotheses the OP5 endpoint problem is exhausted by controlled halo trapping and controlled tunnel resolution. Hence OP5 is discharged as a double-epsilon \(\HC\)-cone halo theorem rather than as an independent terminal transport obstruction.
\end{theorem}

\begin{remark}[What proves and what remains]
The layer proves the formal reduction: capture plus invariance off tunnels reduces OP5 to tunnel resolution. It does not yet prove capture, invariance, finite tunnel classification, or global vanishing of \(\ConeObs^{\chi,(1)}\). Those become the next concrete proof obligations.
\end{remark}

\section{Immediate next proof tasks}

The next executable line should be:
\begin{enumerate}
 \item instantiate \(\rho_r\) and \(\rho_h\) from the current photon-frame \(\HC\)-collar readout;
 \item compute the first symbolic boundary faces \(\rho_r=\epsr\) and \(\rho_h=\epsh\);
 \item define the exit Jacobian \(J_{\Exit}\) in the two-channel coordinates;
 \item test whether the first tunnel menu \(\mathfrak T_{\OP5}^{\mathrm{loc}}\) is complete for the declared halo grammar;
 \item route each nonempty tunnel class to blocked, \(\HC/\MM\)-absorbed, atlas-routed, or residual-projected status.
\end{enumerate}
If these tasks succeed, the finite atlas is no longer trying to cover arbitrary OP5 routes. It only has to cover the gated tunnel exits of the double-epsilon halo.

\section{Conservativity guards}

\begin{enumerate}
 \item does not claim global OP5 closure.
 \item The halo is a readout-side structural filter, not a claim that physical particles literally sit inside a geometric shell.
 \item Tunnel points are selector/rank/obstruction points, not literal wormholes.
 \item \(\HC/\MM\)-absorption is allowed only as explicit debt-free bookkeeping, never as residual deletion.
 \item The finite tunnel menu is a first classification target and must still be tested for completeness.
 \item OP1, OP4, QCD/confinement, CKM/PMNS, and the quark-mass spectrum remain untouched by this layer.
\end{enumerate}


% --- end integrated body for r52 ---

\clearpage
\section*{Integrated Delta Layer: Halo Channel and Response Alphabet}
\addcontentsline{toc}{section}{Integrated Delta Layer: Halo Channel and Response Alphabet}
% --- local macro refresh for r53 ---
\providecommand{\HC}{}
\renewcommand{\HC}{\mathrm{HC}}
\providecommand{\MM}{}
\renewcommand{\MM}{\mathrm{MM}}
\providecommand{\OP}{}
\renewcommand{\OP}{\mathrm{OP}}
\providecommand{\adm}{}
\renewcommand{\adm}{\mathrm{adm}}
\providecommand{\tun}{}
\renewcommand{\tun}{\mathrm{tun}}
\providecommand{\cand}{}
\renewcommand{\cand}{\mathrm{cand}}
\providecommand{\unr}{}
\renewcommand{\unr}{\mathrm{unr}}
\providecommand{\rout}{}
\renewcommand{\rout}{\mathrm{rout}}
\providecommand{\Cone}{}
\renewcommand{\Cone}{\mathrm{cone}}
\providecommand{\rad}{}
\renewcommand{\rad}{\mathrm{rad}}
\providecommand{\proj}{}
\renewcommand{\proj}{\mathrm{proj}}
\providecommand{\coll}{}
\renewcommand{\coll}{\mathrm{coll}}
\providecommand{\rk}{}
\renewcommand{\rk}{\mathrm{rk}}
\providecommand{\ov}{}
\renewcommand{\ov}{\mathrm{ov}}
\providecommand{\Exit}{}
\renewcommand{\Exit}{\mathsf{Exit}}
\providecommand{\Halo}{}
\renewcommand{\Halo}{\mathsf{Halo}}
\providecommand{\Trap}{}
\renewcommand{\Trap}{\mathsf{Trap}}
\providecommand{\Gate}{}
\renewcommand{\Gate}{\mathsf{Gate}}
\providecommand{\Atlas}{}
\renewcommand{\Atlas}{\mathfrak A}
\providecommand{\epsr}{}
\renewcommand{\epsr}{\varepsilon_r}
\providecommand{\epsh}{}
\renewcommand{\epsh}{\varepsilon_h}
\providecommand{\Xop}{}
\renewcommand{\Xop}{\mathcal X_{\OP5}^{\adm}}
\providecommand{\Ttun}{}
\renewcommand{\Ttun}{\mathcal T_{\tun}}
\providecommand{\Dhc}{}
\renewcommand{\Dhc}{D_{\HC\text{-}\mathrm{cone}}}
\providecommand{\ConeObs}{}
\renewcommand{\ConeObs}{D_{\OP5=\HC\text{-}\mathrm{cone}}}
\providecommand{\RHC}{}
\renewcommand{\RHC}{R_{\HC}}
\providecommand{\Upi}{}
\renewcommand{\Upi}{U_{\pi}^{(2)}}
\providecommand{\Pcr}{}
\renewcommand{\Pcr}{\Pi_{\Cone\to\rad}^{\HC,\gamma}}
\providecommand{\rhoR}{}
\renewcommand{\rhoR}{\rho_{r}^{\gamma}}
\providecommand{\rhoH}{}
\renewcommand{\rhoH}{\rho_{h}^{\gamma}}
\providecommand{\Eexit}{}
\renewcommand{\Eexit}{\mathcal E_{\Exit}^{\gamma,K}}
\providecommand{\Jexit}{}
\renewcommand{\Jexit}{J_{\Exit}^{\gamma,K}}
\providecommand{\Th}{}
\renewcommand{\Th}{\Theta_{\gamma,K}^{(2)}}
\providecommand{\rHC}{}
\renewcommand{\rHC}{r_{\HC}^{(1)}}
% --- begin integrated body for r53 ---



\begin{abstract}
This delta layer executes the next proof obligation left by the local pass. The double-epsilon halo symbols are made typographically stable, the two halo channels are instantiated in the photon-frame near-kernel collar, and the exit Jacobian is defined as the differential of the paired radial/HC-residual channel. The result is a first conservative tunnel-menu execution: within the declared photon-frame halo grammar every boundary exit is assigned to one of the six tunnel classes. The layer does not close OP5 globally. It converts the next obstruction into explicit tests: collar admissibility, channel regularity, selector gating, rank-drop isolation, atlas-overlap compatibility, and zero hidden leakage/debt.
\end{abstract}

\section{Macro hygiene and version position}

The layer introduced the double-epsilon halo but the visual notation for the halo and exit operators must be kept stable in downstream TeX. In this layer the intended printed symbols are
\[
\Halo,\qquad \Exit,\qquad J_{\Exit},\qquad \mathcal E_{\Exit}.
\]
For source-level robustness the lower-case spellings are aliased to the same mathematical symbols, but the body of this layer uses the capitalized operators. Thus the rendered document should not display raw control words for the halo or exit operators.

The reduction says that the OP5 route problem is reduced to tunnel resolution once halo capture and halo invariance off tunnels are proved. The present layer executes the first concrete part of that program: it instantiates the two halo channels in the photon-frame HC-collar chart and defines the exit Jacobian which detects rank-drop tunnel candidates.

\section{Photon-frame collar chart}

\begin{definition}[Photon-frame HC-collar chart]
Let \(U_{\gamma,K}\subset \Xop\) be an admissible near-kernel photon-frame chart. The chart carries:
\[
 d_K^G(x),\qquad r_{\gamma}^{(2)}(x),\qquad h_{\gamma}^{(2)}(x),\qquad q_{\gamma}^{(2)}(x):=(h_{\gamma}^{(2)}(x))^2-(r_{\gamma}^{(2)}(x))^2,
\]
where \(d_K^G\) is the metric distance to the collapse kernel, \(r_{\gamma}^{(2)}\) and \(h_{\gamma}^{(2)}\) are the second-shift radial and HC-height coordinates, and \(q_{\gamma}^{(2)}\) is the photon-frame cone residual. The active first-shift HC collar radius is denoted \(\rHC(x)>0\). The admissible collar guard is
\[
 d_K^G(x)\ge \rHC(x)>0.
\]
This guard prevents interpreting a projected line through \(K\) as a literal through-collapse trajectory.
\end{definition}

\begin{definition}[Second-shift tilt operator]
The second-shift cone-to-radial tilt is
\[
\Th
:=
\Pcr\circ \Upi-\Upi\circ\Pcr.
\]
It records the failure of radial projection and the second Hilbert shift to commute. A projected path may visually pass through \(K\) after the second shift; in this layer that appearance is read first as tilt debt rather than as immediate OP5 failure.
\end{definition}

\section{Concrete halo-channel instantiation}

\begin{definition}[Radial halo channel]
The radial halo channel is
\[
\rhoR(x)
:=
 d_K^G\bigl(\Pcr(\Upi x),K\bigr).
\]
The radial halo face is
\[
\rhoR(x)=\epsr.
\]
The photon-frame near-kernel collar is admissible only on the guarded subregion
\[
\rHC(x)\le \rhoR(x)\le \epsr.
\]
Thus \(\rhoR\le \epsr\) means that the projected second-shift readout is inside the declared near-kernel radial halo, while \(\rhoR\ge\rHC\) records the HC periapsis/collar exclusion of literal through-collapse.
\end{definition}

\begin{definition}[HC-residual halo channel]
Fix a harmless normalization floor \(\delta_0>0\) and a tilt weight \(\mu_{\Theta}\ge0\). The HC-residual halo channel is
\[
\rhoH(x)
:=
\frac{|q_{\gamma}^{(2)}(x)|}{(h_{\gamma}^{(2)}(x))^2+(r_{\gamma}^{(2)}(x))^2+\delta_0}
+
\mu_{\Theta}\,\|\Th(x)\|_{G}.
\]
The HC-residual halo face is
\[
\rhoH(x)=\epsh.
\]
The first term measures cone defect. The second term measures the second-shift projection tilt. Therefore \(\rhoH\le\epsh\) means that the readout is simultaneously close to the ideal photon-frame HC cone and free of uncontrolled projection tilt up to the declared tolerance.
\end{definition}

\begin{definition}[Executed double-epsilon halo]
The executed photon-frame halo is
\[
\Halo_{\epsr,\epsh}^{\gamma,K}
:=
\{x\in U_{\gamma,K}\cap \Xop:
\rHC(x)\le \rhoR(x)\le\epsr,
\ \rhoH(x)\le\epsh\}.
\]
The extra lower bound \(\rHC\le\rhoR\) is not an additional escape channel. It is the previously declared HC-collar guard: the admissible branch approaches a periapsis collar rather than crossing \(K\) literally.
\end{definition}

\section{Exit map and exit Jacobian}

\begin{definition}[Two-channel exit map]
The exit map is
\[
\Eexit:U_{\gamma,K}\longrightarrow \mathbb R^2,
\qquad
\Eexit(x):=\bigl(\rhoR(x)-\epsr,\ \rhoH(x)-\epsh\bigr).
\]
The halo boundary is \((\Eexit)^{-1}(0\times(-\infty,0])\cup(\Eexit)^{-1}((-
\infty,0]\times0)\) together with the collar face \(\rhoR=\rHC\). The collar face is not treated as an outward OP5 exit; it is the no-through-collapse reflecting side of the HC collar.
\end{definition}

\begin{definition}[Exit Jacobian]
At a boundary point \(p\), the exit Jacobian is
\[
\Jexit(p)
:=
D\Eexit(p)
=
\begin{pmatrix}
D\rhoR(p)\\[2mm]
D\rhoH(p)
\end{pmatrix}.
\]
The maximal two-channel rank is \(\operatorname{rank}_{\max}=2\). The rank-drop tunnel condition is therefore
\[
\operatorname{rank}\Jexit(p)<2,
\]
which is equivalently
\[
D\rhoR(p)\wedge D\rhoH(p)=0
\]
whenever the two channel differentials are defined.
\end{definition}

\begin{definition}[Outward exit condition]
A boundary point \(p\in\partial\Halo_{\epsr,\epsh}^{\gamma,K}\) is outward in the radial channel if there exists an admissible update vector \(v\) with
\[
D\rhoR(p)[v]>0
\quad\text{on}\quad \rhoR(p)=\epsr.
\]
It is outward in the HC-residual channel if
\[
D\rhoH(p)[v]>0
\quad\text{on}\quad \rhoH(p)=\epsh.
\]
A genuine tunnel candidate must be outward in at least one of these two channels and must pass the selector/rank tests.
\end{definition}

\section{First tunnel-menu execution}

\begin{proposition}[Boundary exhaustion by the tunnel menu]
Within the photon-frame halo grammar, every outward boundary exit candidate belongs to at least one of the six tunnel classes
\[
\mathfrak T_{\OP5}^{\mathrm{loc}}
=
\{T_r,T_h,T_{rh},T_{\chi},T_{\rk},T_{\ov}\}.
\]
\end{proposition}

\begin{proof}
Let \(p\in\partial\Halo_{\epsr,\epsh}^{\gamma,K}\) be an outward exit candidate. If \(p\) lies on the radial face with \(\rhoR(p)=\epsr\) and \(\rhoH(p)<\epsh\), then \(p\in T_r\). If \(p\) lies on the HC-residual face with \(\rhoR(p)<\epsr\) and \(\rhoH(p)=\epsh\), then \(p\in T_h\). If both equalities hold, then \(p\in T_{rh}\). If the tunnel selector changes value or must be evaluated at a selector discontinuity, then \(p\in T_{\chi}\). If \(\operatorname{rank}\Jexit(p)<2\), then \(p\in T_{\rk}\). Finally, if the exit occurs while changing from one photon-frame chart to another, then the point lies in an atlas-overlap region and belongs to \(T_{\ov}\). These cases exhaust the declared boundary, selector, rank, and overlap mechanisms of the grammar.
\end{proof}

\begin{definition}[response alphabet]
For each tunnel class \(T\), the first execution response is drawn from
\[
\mathrm{Resp}_{\mathrm{loc}}(T)
\in
\{\mathrm{blocked},\mathrm{absorbed}_{\HC/\MM},\mathrm{atlas\ routed},\mathrm{rank\ isolated},\mathrm{residual\ open}\}.
\]
Only the first three responses are closing responses. Rank-isolated is a reduction response: it converts the question into a finite critical-set test. Residual-open is a non-closing response.
\end{definition}

\begin{center}
\begin{tabular}{>{\raggedright\arraybackslash}p{2.4cm} >{\raggedright\arraybackslash}p{5.2cm} >{\raggedright\arraybackslash}p{5.2cm}}
\toprule
Tunnel class & test & conservative response \\
\midrule
\(T_r\) & \(\rhoR=\epsr,\ \rhoH<\epsh\), radial outward derivative & atlas routed if covered by a strict contraction chart; otherwise residual open \\
\(T_h\) & \(\rhoH=\epsh,\ \rhoR<\epsr\), HC-residual outward derivative & absorbed if HC/MM absorption is explicit and debt-free; otherwise residual open \\
\(T_{rh}\) & \(\rhoR=\epsr,\ \rhoH=\epsh\) & split into \(T_r\) and \(T_h\) if rank two; rank isolated if rank drops \\
\(T_{\chi}\) & tunnel selector changes or is discontinuous & blocked when selector returns zero; otherwise routed to atlas-overlap or residual open \\
\(T_{\rk}\) & \(D\rhoR\wedge D\rhoH=0\) & rank isolated as a critical set of the two-channel exit map \\
\(T_{\ov}\) & chart transition or finite-atlas overlap debt & atlas routed if transition maps preserve strict contraction and zero debt \\
\bottomrule
\end{tabular}
\end{center}

\section{Rank-drop isolation and finite critical set target}

\begin{definition}[Rank-drop critical register]
The rank-drop register is
\[
\mathcal C_{\rk}^{\mathrm{loc}}
:=
\{p\in\partial\Halo_{\epsr,\epsh}^{\gamma,K}:
D\rhoR(p)\wedge D\rhoH(p)=0,
\ \chi_{\tun}(p)=1\}.
\]
It is the only place where the two halo faces can fail to give independent boundary directions.
\end{definition}

\begin{proposition}[Rank-drop reduction]
Assume \(\rhoR\) and \(\rhoH\) are piecewise analytic on a finite photon-frame atlas. Then the rank-drop tunnel class is reduced to the piecewise analytic critical-set problem
\[
D\rhoR\wedge D\rhoH=0
\quad\text{on}\quad
\partial\Halo_{\epsr,\epsh}^{\gamma,K}.
\]
If this critical set is empty, blocked by \(\chi_{\tun}=0\), or covered by finitely many strict contraction charts with zero hidden debt, then \(T_{\rk}\) contributes no uncontrolled OP5 exit.
\end{proposition}

\begin{proof}
By definition, rank-drop occurs exactly when the differential of the exit map has rank less than two. For a two-channel map this is exactly the vanishing of the wedge \(D\rhoR\wedge D\rhoH\), away from nondifferentiable selector points which are assigned to \(T_{\chi}\). Piecewise analyticity transfers the problem to a finite critical-set calculation on each chart. The stated empty, blocked, or chart-routed alternatives are precisely the controlled outcomes in the response alphabet.
\end{proof}

\section{First executed closure implication}

\begin{theorem}[executed halo/tunnel reduction]
Assume the OP5=HC-cone working identification and the photon-frame channel definitions above. Suppose the following four conditions hold on the declared finite photon-frame atlas:
\begin{enumerate}
 \item \textbf{Channel regularity:} \(\rhoR\) and \(\rhoH\) are piecewise analytic and compatible with the HC-collar guard \(\rHC\le\rhoR\).
 \item \textbf{Selector completeness:} every admissible outward boundary exit is evaluated by \(\chi_{\tun}\).
 \item \textbf{Tunnel response closure:} each nonempty class in \(\{T_r,T_h,T_{rh},T_{\chi},T_{\rk},T_{\ov}\}\) receives one of the three closing responses: blocked, debt-free HC/MM absorbed, or atlas routed through strict contraction charts.
 \item \textbf{No hidden debt:}
 \[
 L_{\unr}=0,
 \qquad
 B_{\HC/\MM}=0,
 \qquad
 \tau=0.
 \]
\end{enumerate}
Then the double-epsilon photon-frame halo has no uncontrolled OP5 exit. Consequently, if halo capture and halo invariance off tunnels are also proved globally for the admissible OP5 readout class, OP5 reduces to a closed HC-cone halo theorem.
\end{theorem}

\begin{proof}
Boundary exhaustion assigns every outward exit candidate to the finite tunnel menu. The response-closure assumption assigns every nonempty class to a controlled outcome. Blocked classes do not exit; debt-free HC/MM-absorbed classes have no unpaid residual; atlas-routed classes return to strict contraction charts; and the no-hidden-debt condition rules out suppressed leakage. Hence no uncontrolled exit remains inside the executed halo grammar. Combining this with capture and invariance would trap every admissible OP5 branch or route every tunnel exit through controlled charts. This is exactly the HC-cone halo theorem formulation of OP5.
\end{proof}

\section{What still has to be shown for an actual OP5 close}

The layer executes the definitions and the first tunnel-menu reduction. It does not yet prove the global OP5 close. The remaining non-negotiable proof obligations are:
\begin{enumerate}
 \item prove finite-atlas coverage of the admissible photon-frame near-kernel collar;
 \item prove halo capture for every admissible OP5 readout;
 \item prove halo invariance away from the finite tunnel menu;
 \item compute or bound \(\mathcal C_{\rk}^{\mathrm{loc}}\);
 \item prove that selector-switch points are blocked or chart-routed;
 \item prove that all atlas overlaps preserve the strict margin \(\sup_j(\widehat\lambda_j+\eta_j)<1\);
 \item prove zero hidden leakage and zero hidden HC/MM debt.
\end{enumerate}

If these obligations are discharged, OP5 is no longer an all-routes problem. It is exhausted by the double-epsilon halo, the finite tunnel menu, and strict contraction routing at the tunnel exits.

\section{Conservativity guards}

\begin{enumerate}
 \item does not claim global OP5 closure.
 \item The executed halo is a readout-side HC-collar filter, not a literal particle shell.
 \item The collar face \(\rhoR=\rHC\) is a no-through-collapse guard and not an escape face.
 \item The second-shift projected through-K appearance remains tilt debt unless the tunnel tests certify otherwise.
 \item HC/MM absorption must remain explicit and debt-free.
 \item OP1, OP4, QCD/confinement, CKM/PMNS, and the quark-mass spectrum are not changed by this layer.
\end{enumerate}


% --- end integrated body for r53 ---

\clearpage
\section*{Integrated Delta Layer: Normal-Collar Rank Test}
\addcontentsline{toc}{section}{Integrated Delta Layer: Normal-Collar Rank Test}
% --- local macro refresh for r54 ---
\providecommand{\HC}{}
\renewcommand{\HC}{\mathrm{HC}}
\providecommand{\MM}{}
\renewcommand{\MM}{\mathrm{MM}}
\providecommand{\OP}{}
\renewcommand{\OP}{\mathrm{OP}}
\providecommand{\adm}{}
\renewcommand{\adm}{\mathrm{adm}}
\providecommand{\tun}{}
\renewcommand{\tun}{\mathrm{tun}}
\providecommand{\unr}{}
\renewcommand{\unr}{\mathrm{unr}}
\providecommand{\rk}{}
\renewcommand{\rk}{\mathrm{rk}}
\providecommand{\ov}{}
\renewcommand{\ov}{\mathrm{ov}}
\providecommand{\HaloOp}{}
\renewcommand{\HaloOp}{\mathsf{Halo}}
\providecommand{\ExitOp}{}
\renewcommand{\ExitOp}{\mathsf{Exit}}
\providecommand{\Atlas}{}
\renewcommand{\Atlas}{\mathfrak A}
\providecommand{\epsr}{}
\renewcommand{\epsr}{\varepsilon_r}
\providecommand{\epsh}{}
\renewcommand{\epsh}{\varepsilon_h}
\providecommand{\Xop}{}
\renewcommand{\Xop}{\mathcal X_{\OP 5}^{\adm}}
\providecommand{\rhoR}{}
\renewcommand{\rhoR}{\rho_{r}^{\gamma}}
\providecommand{\rhoH}{}
\renewcommand{\rhoH}{\rho_{h}^{\gamma}}
\providecommand{\Eexit}{}
\renewcommand{\Eexit}{\mathcal E_{\ExitOp}^{\gamma,K}}
\providecommand{\Jexit}{}
\renewcommand{\Jexit}{J_{\ExitOp}^{\gamma,K}}
\providecommand{\rHC}{}
\renewcommand{\rHC}{r_{\HC}^{(1)}}
\providecommand{\Crk}{}
\renewcommand{\Crk}{\mathcal C_{\rk}^{\mathrm{loc}}}
\providecommand{\Cchi}{}
\renewcommand{\Cchi}{\mathcal C_{\chi}^{\mathrm{loc}}}
\providecommand{\Wreg}{}
\renewcommand{\Wreg}{W_j^{\mathrm{loc}}}
\providecommand{\Lunr}{}
\renewcommand{\Lunr}{L_{\unr}}
\providecommand{\Bhcmm}{}
\renewcommand{\Bhcmm}{B_{\HC/\MM}}
\providecommand{\marg}{}
\renewcommand{\marg}{\mathfrak m}
% --- begin integrated body for r54 ---



\begin{abstract}
This delta layer executes the next obstruction left by the local pass. The exit map is decomposed chartwise on a finite photon-frame halo atlas. The rank-drop tunnel class is bounded by an explicit determinant register, and the selector-switch tunnel class is bounded by a finite switch skeleton. The result is not a global OP5 close. It reduces the uncontrolled tunnel search to two finite residual registers plus atlas-overlap margin checks.
\end{abstract}

\section{Purpose and notation hygiene, declared local OP1/OP4 closure, release-final}

The local proof chain halo channel reduced the OP5 near-kernel route problem to a two-channel boundary problem. The printed operators in this layer are
\[
\HaloOp,\qquad \ExitOp,\qquad \Jexit,\qquad \Eexit .
\]
No lowercase display macro is used. The task is to execute the instruction
\[
\operatorname{rank}\Jexit<2
\quad\Longleftrightarrow\quad
D\rhoR\wedge D\rhoH=0,
\]
and to combine it with the selector-switch test.

\section{Finite photon-frame halo atlas}

\begin{definition}[finite halo atlas]
An photon-frame halo atlas is a finite family
\[
\Atlas_{\gamma,K}^{\mathrm{loc}}=\{(U_j,\varphi_j)\}_{j=1}^{M}
\]
covering the active halo boundary
\[
\partial \HaloOp_{\epsr,\epsh}^{\gamma,K}
\cap \{\rHC\le \rhoR\le \epsr,\ \rhoH\le\epsh\}.
\]
On each chart we write
\[
r_j:=\rhoR\circ\varphi_j^{-1},
\qquad
h_j:=\rhoH\circ\varphi_j^{-1}.
\]
The chart is regular for when these two functions are piecewise analytic and all nonsmooth selector boundaries are explicitly assigned to the switch register.
\end{definition}

\begin{definition}[Exit determinant]
On a smooth two-channel chart piece with active coordinates \((u,v)\), define
\[
\Wreg(u,v)
=
\det
\begin{pmatrix}
\partial_u r_j & \partial_v r_j\\[1mm]
\partial_u h_j & \partial_v h_j
\end{pmatrix}
=
\partial_u r_j\,\partial_v h_j-
\partial_v r_j\,\partial_u h_j.
\]
The rank-drop register is
\[
\Crk=
\bigcup_j
\{p\in U_j\cap\partial \HaloOp_{\epsr,\epsh}^{\gamma,K}:
W_j^{\mathrm{loc}}(p)=0,\ \chi_{\tun}(p)=1\}.
\]
Thus the rank-drop tunnel class is not an arbitrary route family; it is a finite union of explicitly declared chartwise critical loci.
\end{definition}

\begin{lemma}[Rank-drop localization]
Every rank-drop tunnel point belongs either to \(\Crk\) or to a selector-switch boundary.
\end{lemma}

\begin{proof}
On a smooth chart piece, gives rank drop exactly when the two channel differentials are dependent. In active coordinates this is exactly \(W_j^{\mathrm{loc}}=0\). If the point is not smooth because the selector changes there, it is assigned to the selector-switch register. Since the atlas is finite, the localization is finite.
\end{proof}

\section{Normal collar charts}

\begin{definition}[Normal collar form]
A smooth halo chart piece is in normal collar form when it admits coordinates \((a,b,z)\) such that
\[
r_j(a,b,z)=r_{\HC,j}^{(1)}(z)+\epsr a,
\qquad
h_j(a,b,z)=\epsh b+\Delta_j(a,b,z),
\]
with
\[
0\le a\le1,
\qquad
0\le b\le1,
\qquad
|\partial_b\Delta_j(a,b,z)|\le \theta_j\epsh,
\qquad
0\le\theta_j<1.
\]
The number \(1-\theta_j\) is the local HC-residual transversality margin.
\end{definition}

\begin{proposition}[Normal collar charts have no rank drop]
On a normal collar chart piece,
\[
W_j^{\mathrm{loc}}=\epsr\bigl(\epsh+\partial_b\Delta_j\bigr)
\]
in the active \((a,b)\)-plane. Hence
\[
|W_j^{\mathrm{loc}}|\ge \epsr\epsh(1-\theta_j)>0.
\]
Consequently \(\Crk\) is empty on that chart piece.
\end{proposition}

\begin{proof}
The active derivatives are \(\partial_a r_j=\epsr\), \(\partial_b r_j=0\), and \(\partial_b h_j=\epsh+\partial_b\Delta_j\). Therefore the determinant is \(\epsr(\epsh+\partial_b\Delta_j)\). The assumed bound on \(\partial_b\Delta_j\) keeps this determinant strictly away from zero.
\end{proof}

\section{Selector-switch register}

\begin{definition}[Finite selector skeleton]
The tunnel selector is a finite Boolean gate
\[
\chi_{\tun}:\partial \HaloOp_{\epsr,\epsh}^{\gamma,K}\longrightarrow\{0,1\}.
\]
It is constant on cells cut out by finitely many piecewise analytic switch functions
\[
\sigma_{j,\ell}:U_j\longrightarrow\mathbb R.
\]
The selector-switch register is
\[
\Cchi=
\bigcup_{j,\ell}
\{p\in U_j\cap\partial \HaloOp_{\epsr,\epsh}^{\gamma,K}:
\sigma_{j,\ell}(p)=0\}.
\]
A switch point is closing only if the outward-side selector is zero or if a strict atlas-transition chart routes it with zero hidden debt.
\end{definition}

\begin{lemma}[Selector localization]
Every selector ambiguity or selector discontinuity of an outward boundary candidate belongs to \(\Cchi\).
\end{lemma}

\begin{proof}
Away from the switch skeleton all switch signs are fixed, so the selector is locally constant. Hence the selector can change only where at least one \(\sigma_{j,\ell}\) vanishes.
\end{proof}

\section{Tunnel response execution}

The six tunnel classes receive the following responses:
\[
\begin{array}{c|l}
\text{class} & \text{response}\\
\hline
T_r & \text{atlas routed if strict contraction covers the radial face; otherwise open or blocked}\\
T_h & \text{debt-free HC/MM absorbed if explicit absorption certificate exists; otherwise open}\\
T_{rh} & \text{split into }T_r,T_h\text{ if }W_j^{\mathrm{loc}}\ne0;\text{ rank isolated if }W_j^{\mathrm{loc}}=0\\
T_{\chi} & \text{blocked if outward selector is zero; otherwise switch isolated or atlas routed}\\
T_{\rk} & \text{rank isolated inside }\Crk\text{ until empty/block/absorb/route is proved}\\
T_{\ov} & \text{atlas routed if overlap margins remain strictly below one and debt is zero}
\end{array}
\]

\begin{definition}[Overlap margin]
For a chart transition \(\Gamma_{jk}:U_j\to U_k\), set
\[
\marg_{jk}
:=
\sup_{x\in U_j\cap U_k}
\bigl(\widehat\lambda_k(\Gamma_{jk}x)+\eta_k(\Gamma_{jk}x)\bigr).
\]
The overlap is contraction-preserving if
\[
\marg_{jk}<1,
\qquad
\Lunr=0,
\qquad
\Bhcmm=0
\]
on the overlap.
\end{definition}

\begin{theorem}[rank/switch register reduction]
Assume the OP5=HC-cone working identification and the photon-frame halo-channel definitions. Suppose the finite halo atlas satisfies:
\begin{enumerate}
\item each smooth chart piece is piecewise analytic;
\item selector changes occur only on \(\Cchi\);
\item every normal collar chart piece has \(\theta_j<1\);
\item all active atlas overlaps are contraction-preserving;
\item hidden leakage and hidden HC/MM debt vanish on the executed halo boundary:
\[
\Lunr=0,
\qquad
\Bhcmm=0.
\]
\end{enumerate}
Then every outward tunnel candidate is either blocked, debt-free HC/MM absorbed, atlas routed, or contained in
\[
\mathcal R_{\mathrm{res}}^{\mathrm{loc}}:=\Crk\cup\Cchi.
\]
If, in addition, \(\Crk\) is empty, blocked, absorbed, or routed and \(\Cchi\) is blocked or routed, then the tunnel menu has no uncontrolled OP5 exit.
\end{theorem}

\begin{proof}
The boundary exhaustion assigns every outward exit to the finite tunnel menu. The determinant localization sends rank-drop points to \(\Crk\), while normal collar transversality eliminates rank drop on regular pieces. The selector localization sends selector ambiguity to \(\Cchi\). Radial, HC-residual, corner, rank, selector, and overlap tunnel classes are then handled by the response table. The only non-closing objects left are the two residual registers. If these are also empty, blocked, absorbed, or routed, no uncontrolled tunnel remains.
\end{proof}

\section{First executable normal-form sample}

For the simplest ideal photon-frame collar sample, take active coordinates \((a,b)\) and set
\[
r(a,b)=\rHC+\epsr a,
\qquad
h(a,b)=\epsh b+\delta ab,
\qquad
|\delta|<\epsh.
\]
Then
\[
W^{\mathrm{loc}}(a,b)
=
\det
\begin{pmatrix}
\epsr & 0\\
\delta b & \epsh+\delta a
\end{pmatrix}
=
\epsr(\epsh+\delta a)>0.
\]
Thus the rank-drop register is empty in the ideal normal-form sample. Any remaining tunnel issue in that sample must come from a selector switch, an atlas overlap, or an explicitly declared HC/MM absorption boundary.

\section{Conservativity guards}

This layer proves a structural reduction, not a final physical closure of OP5. The double-epsilon halo remains a readout-side structural filter, not a literal particle shell. Tunnel points are selector, rank, or obstruction points, not literal spacetime wormholes. The ideal normal-form sample is a diagnostic model, not a measurement claim. HC/MM absorption must remain explicit and debt-free. OP1, OP4, QCD/confinement, CKM/PMNS, and the quark-mass spectrum are not changed by this layer.


% --- end integrated body for r54 ---

\clearpage
\section*{Integrated Delta Layer: Selector and Overlap Test}
\addcontentsline{toc}{section}{Integrated Delta Layer: Selector and Overlap Test}
% --- local macro refresh for r55 ---
\providecommand{\HC}{}
\renewcommand{\HC}{\mathrm{HC}}
\providecommand{\MM}{}
\renewcommand{\MM}{\mathrm{MM}}
\providecommand{\OP}{}
\renewcommand{\OP}{\mathrm{OP}}
\providecommand{\adm}{}
\renewcommand{\adm}{\mathrm{adm}}
\providecommand{\tun}{}
\renewcommand{\tun}{\mathrm{tun}}
\providecommand{\unr}{}
\renewcommand{\unr}{\mathrm{unr}}
\providecommand{\rk}{}
\renewcommand{\rk}{\mathrm{rk}}
\providecommand{\ov}{}
\renewcommand{\ov}{\mathrm{ov}}
\providecommand{\HaloOp}{}
\renewcommand{\HaloOp}{\mathsf{Halo}}
\providecommand{\ExitOp}{}
\renewcommand{\ExitOp}{\mathsf{Exit}}
\providecommand{\Atlas}{}
\renewcommand{\Atlas}{\mathfrak A}
\providecommand{\epsr}{}
\renewcommand{\epsr}{\varepsilon_r}
\providecommand{\epsh}{}
\renewcommand{\epsh}{\varepsilon_h}
\providecommand{\Xop}{}
\renewcommand{\Xop}{\mathcal X_{\OP 5}^{\adm}}
\providecommand{\rhoR}{}
\renewcommand{\rhoR}{\rho_{r}^{\gamma}}
\providecommand{\rhoH}{}
\renewcommand{\rhoH}{\rho_{h}^{\gamma}}
\providecommand{\Eexit}{}
\renewcommand{\Eexit}{\mathcal E_{\ExitOp}^{\gamma,K}}
\providecommand{\Jexit}{}
\renewcommand{\Jexit}{J_{\ExitOp}^{\gamma,K}}
\providecommand{\rHC}{}
\renewcommand{\rHC}{r_{\HC}^{(1)}}
\providecommand{\Crk}{}
\renewcommand{\Crk}{\mathcal C_{\rk}^{\mathrm{loc}}}
\providecommand{\Cchi}{}
\renewcommand{\Cchi}{\mathcal C_{\chi}^{\mathrm{loc}}}
\providecommand{\Wreg}{}
\renewcommand{\Wreg}{W_j^{\mathrm{loc}}}
\providecommand{\Lunr}{}
\renewcommand{\Lunr}{L_{\unr}}
\providecommand{\Bhcmm}{}
\renewcommand{\Bhcmm}{B_{\HC/\MM}}
\providecommand{\marg}{}
\renewcommand{\marg}{\mathfrak m}
\providecommand{\thetaMax}{}
\renewcommand{\thetaMax}{\theta_{\max}^{\mathrm{loc}}}
\providecommand{\Ncollar}{}
\renewcommand{\Ncollar}{\mathsf N_{\HC}^{\gamma,K}}
% --- begin integrated body for r55 ---



\begin{abstract}
This delta layer executes the follow-up left by the local pass. It gives a sufficient, finite, chartwise criterion which places the relevant photon-frame halo charts into normal-collar form with a uniform margin
\(
|\partial_b\Delta_j|\le \theta_j\varepsilon_h,
\theta_j<1.
\)
Under this certificate the rank-drop register is empty on all certified collar pieces. The layer does not close OP5 globally. It removes the rank-drop tunnel channel from the certified photon-frame normal-collar atlas and leaves only selector-switch, atlas-overlap, leakage/debt, or non-certified residual pieces.
\end{abstract}

\section{Purpose}

The layer reduced uncontrolled halo exits to a rank-drop register, a selector-switch register, and atlas-overlap margin checks. The key unresolved certificate was whether the photon-frame halo charts may actually be placed in normal-collar form. This layer gives the required local criterion and the first finite symbolic uniformization certificate.

The printed operators remain renderer-safe:
\[
\HaloOp,\qquad \ExitOp,\qquad \Jexit,\qquad \Eexit .
\]
No raw lowercase display macro is used.

\section{Photon-frame collar coordinates}

\begin{definition}[Certified photon-frame collar piece]
A smooth halo chart piece \(U_j\subset \partial\HaloOp_{\epsr,\epsh}^{\gamma,K}\) is the local pass-certified when it admits coordinates
\[
(a,b,z)\in [0,1]\times[0,1]\times Z_j
\]
with the following three properties:
\begin{enumerate}
\item \(a\) is the normalized radial collar coordinate after the second photon-frame shift,
\[
\rhoR(a,b,z)=r_{\HC,j}^{(1)}(z)+\epsr a;
\]
\item \(b\) is the normalized HC-cone height coordinate;
\item the HC-residual channel splits as
\[
\rhoH(a,b,z)=\epsh b+\Delta_j(a,b,z)
\]
with a declared tilt-debt derivative bound
\[
|\partial_b\Delta_j(a,b,z)|\le \theta_j\epsh,
\qquad 0\le \theta_j<1.
\]
\end{enumerate}
The number \(1-\theta_j\) is the normal-collar safety margin.
\end{definition}

\begin{remark}
The coordinate \(b\) is not defined to be \(\rhoH/\epsh\). It is the HC-cone height coordinate inherited from the photon-frame collar. This avoids a circular proof of transversality. The nontrivial check is exactly the derivative bound on \(\Delta_j\).
\end{remark}

\begin{definition}[Uniform certificate]
A finite photon-frame halo atlas
\[
\Atlas_{\gamma,K}^{\mathrm{loc}}=\{U_j\}_{j=1}^M
\]
is uniformly normal-collar certified when every smooth active piece is the local pass-certified and
\[
\thetaMax:=\max_{1\le j\le M}\theta_j<1.
\]
The certified rank-drop register is then
\[
\Crk=
\bigcup_j
\{p\in U_j\cap\partial\HaloOp_{\epsr,\epsh}^{\gamma,K}: W_j^{\mathrm{loc}}(p)=0,
\chi_{\tun}(p)=1\}.
\]
\end{definition}

\section{Rank-drop elimination on certified pieces}

\begin{lemma}[determinant lower bound]
On an the local pass-certified collar piece,
\[
\Wreg(a,b,z)=\epsr\bigl(\epsh+\partial_b\Delta_j(a,b,z)\bigr).
\]
Consequently
\[
|\Wreg(a,b,z)|
\ge \epsr\epsh(1-\theta_j)>0.
\]
\end{lemma}

\begin{proof}
The active derivatives in the \((a,b)\)-plane are
\[
\partial_a\rhoR=\epsr,
\qquad
\partial_b\rhoR=0,
\qquad
\partial_b\rhoH=\epsh+\partial_b\Delta_j.
\]
Therefore the exit determinant is
\[
\det
\begin{pmatrix}
\partial_a\rhoR&\partial_b\rhoR\\[1mm]
\partial_a\rhoH&\partial_b\rhoH
\end{pmatrix}
=\epsr(\epsh+\partial_b\Delta_j).
\]
The derivative bound gives \(\epsh+\partial_b\Delta_j\ge \epsh(1-\theta_j)>0\) after orienting the collar height positively. Hence the determinant is bounded away from zero.
\end{proof}

\begin{proposition}[Certified collar pieces contain no rank-drop tunnels]
If \(U_j\) is the local pass-certified, then
\[
U_j\cap\Crk=\varnothing .
\]
If the atlas is uniformly the local pass-certified, then the rank-drop tunnel channel is empty on the certified smooth halo boundary:
\[
\Crk\cap\bigcup_j U_j^{\mathrm{cert}}=\varnothing .
\]
\end{proposition}

\begin{proof}
Rank drop is equivalent to \(W_j^{\mathrm{loc}}=0\) on a smooth chart piece. The determinant lower bound excludes this equality on each certified piece. Taking the finite union gives the atlas statement.
\end{proof}

\section{Finite symbolic certificate}

The first executable certificate uses three representative collar pieces. These are not physical measurements; they are the minimal symbolic chart types required by the current halo grammar.

\[
\begin{array}{c|c|c|c|c}
\text{piece} & \text{role} & \Delta_j(a,b,z)/\epsh & \theta_j & \text{rank response}\\
\hline
U_1 & \text{radial face} & \frac14 ab & \frac14 & \text{empty}\\[1mm]
U_2 & \text{HC-residual face} & \frac38 b(1-a) & \frac38 & \text{empty}\\[1mm]
U_3 & \text{corner/overlap collar} & \frac12 ab(1-a) & \frac12 & \text{empty}
\end{array}
\]

Thus
\[
\thetaMax=\frac12<1
\]
and hence
\[
|W_j^{\mathrm{loc}}|\ge \frac12\epsr\epsh
\]
on the three certified representative pieces.

\begin{lemma}[First finite symbolic normal-collar sample]
For the three-piece certificate above, the rank-drop register is empty on all certified pieces.
\end{lemma}

\begin{proof}
For every listed piece, \(\theta_j\le 1/2\). The determinant lower bound gives strict positivity of \(W_j^{\mathrm{loc}}\). Hence no listed piece contributes to \(\Crk\).
\end{proof}

\section{Remaining tunnel channels}

After the certificate, the local proof chain tunnel menu is reduced as follows:
\[
\begin{array}{c|l}
\text{tunnel class} & \text{status}\\
\hline
T_r & \text{covered by normal collar if radial face is in the certified atlas}\\
T_h & \text{covered by normal collar or debt-free HC/MM absorption}\\
T_{rh} & \text{rank split eliminated on certified pieces; overlap still checked}\\
T_\chi & \text{unchanged: selector-switch register must be blocked or routed}\\
T_{\rk} & \text{empty on certified collar pieces}\\
T_{\ov} & \text{unchanged: transition margins must remain strictly below one}
\end{array}
\]

Thus does not remove every OP5 obstruction. It removes the rank-drop channel exactly where the photon-frame halo atlas has the certified normal-collar form.

\begin{theorem}[normal-collar rank elimination]
Assume the OP5=HC-cone working identification, the halo-channel execution, and the rank/switch register reduction. Suppose a finite photon-frame halo atlas is uniformly the local pass-certified with \(\thetaMax<1\). Then no rank-drop tunnel occurs on the certified smooth halo boundary. Every remaining uncontrolled candidate must lie in one of the following residual classes:
\[
\mathcal R_{\mathrm{res}}^{\mathrm{loc}}
\subseteq
\mathcal C_{\chi}^{\mathrm{loc}}
\cup
\mathcal C_{\ov}^{\mathrm{loc}}
\cup
\mathcal C_{\mathrm{leak/debt}}^{\mathrm{loc}}
\cup
\mathcal C_{\mathrm{uncert}}^{\mathrm{loc}},
\]
where \(\mathcal C_{\mathrm{uncert}}^{\mathrm{loc}}\) denotes non-certified pieces not yet placed in normal-collar form.
If the selector-switch register is blocked or routed, all atlas-overlap margins are contraction-preserving, hidden leakage and HC/MM debt vanish, and \(\mathcal C_{\mathrm{uncert}}^{\mathrm{loc}}=\varnothing\), then the halo tunnel menu has no uncontrolled exit on the finite certified atlas.
\end{theorem}

\begin{proof}
The certified determinant lower bound eliminates \(T_{\rk}\) on every smooth certified chart piece. The reduction already states that all remaining outward candidates are rank, selector, overlap, leakage/debt, or residual non-certified pieces. Removing rank on certified pieces leaves exactly the displayed residual classes. If these classes are blocked, routed, or empty under the stated hypotheses, no uncontrolled exit remains on the finite certified atlas.
\end{proof}

\section{Meaning for OP5}

The result is a strong local structural advance. It says that, provided the photon-frame halo atlas is uniformly normal-collar certified, OP5 is no longer obstructed by a generic rank-collapse phenomenon near the HC collar. The hard work is pushed to four explicit checks:
\[
\text{selector switches},\qquad
\text{atlas overlaps},\qquad
\text{hidden leakage/debt},\qquad
\text{uncertified pieces}.
\]
A global OP5 close would require these four checks to vanish, block, or route over the full finite atlas.

\section{Conservativity guards}

This layer does not claim a global OP5 closure. It proves a conditional rank-drop elimination theorem for certified photon-frame normal-collar charts. The halo remains a readout-side structural filter, not a literal particle shell. Tunnel points remain selector, rank, overlap, leakage, or obstruction points. HC/MM absorption must remain explicit and debt-free. OP1, OP4, QCD/confinement, CKM/PMNS, and the quark-mass spectrum are untouched.


% --- end integrated body for r55 ---

\clearpage
\section*{Integrated Delta Layer: Finite Capture}
\addcontentsline{toc}{section}{Integrated Delta Layer: Finite Capture}
% --- local macro refresh for r56 ---
\providecommand{\HC}{}
\renewcommand{\HC}{\mathrm{HC}}
\providecommand{\MM}{}
\renewcommand{\MM}{\mathrm{MM}}
\providecommand{\OP}{}
\renewcommand{\OP}{\mathrm{OP}}
\providecommand{\adm}{}
\renewcommand{\adm}{\mathrm{adm}}
\providecommand{\tun}{}
\renewcommand{\tun}{\mathrm{tun}}
\providecommand{\unr}{}
\renewcommand{\unr}{\mathrm{unr}}
\providecommand{\ov}{}
\renewcommand{\ov}{\mathrm{ov}}
\providecommand{\rk}{}
\renewcommand{\rk}{\mathrm{rk}}
\providecommand{\HaloOp}{}
\renewcommand{\HaloOp}{\mathsf{Halo}}
\providecommand{\ExitOp}{}
\renewcommand{\ExitOp}{\mathsf{Exit}}
\providecommand{\Atlas}{}
\renewcommand{\Atlas}{\mathfrak A}
\providecommand{\epsr}{}
\renewcommand{\epsr}{\varepsilon_r}
\providecommand{\epsh}{}
\renewcommand{\epsh}{\varepsilon_h}
\providecommand{\Xop}{}
\renewcommand{\Xop}{\mathcal X_{\OP 5}^{\adm}}
\providecommand{\rhoR}{}
\renewcommand{\rhoR}{\rho_{r}^{\gamma}}
\providecommand{\rhoH}{}
\renewcommand{\rhoH}{\rho_{h}^{\gamma}}
\providecommand{\Jexit}{}
\renewcommand{\Jexit}{J_{\ExitOp}^{\gamma,K}}
\providecommand{\Cchi}{}
\renewcommand{\Cchi}{\mathcal C_{\chi}^{\mathrm{loc}}}
\providecommand{\Cov}{}
\renewcommand{\Cov}{\mathcal C_{\ov}^{\mathrm{loc}}}
\providecommand{\Crk}{}
\renewcommand{\Crk}{\mathcal C_{\rk}^{\mathrm{loc}}}
\providecommand{\Lunr}{}
\renewcommand{\Lunr}{L_{\unr}}
\providecommand{\Bhcmm}{}
\renewcommand{\Bhcmm}{B_{\HC/\MM}}
\providecommand{\cert}{}
\renewcommand{\cert}{\mathrm{cert}}
\providecommand{\route}{}
\renewcommand{\route}{\mathrm{route}}
\providecommand{\block}{}
\renewcommand{\block}{\mathrm{block}}
\providecommand{\abs}{}
\renewcommand{\abs}{\mathrm{abs}}
\providecommand{\res}{}
\renewcommand{\res}{\mathrm{res}}
\providecommand{\marg}{}
\renewcommand{\marg}{\mathfrak m}
\providecommand{\Ssel}{}
\renewcommand{\Ssel}{\mathfrak S_{\chi}^{\mathrm{loc}}}
\providecommand{\Gtrans}{}
\renewcommand{\Gtrans}{\mathfrak G_{\ov}^{\mathrm{loc}}}
\providecommand{\Ncollar}{}
\renewcommand{\Ncollar}{\mathsf N_{\HC}^{\gamma,K}}
% --- begin integrated body for r56 ---



\begin{abstract}
This delta layer executes the two residual channels left by on certified photon-frame normal-collar pieces: selector switches and atlas overlaps. It introduces a finite selector-switch stratum, a finite overlap transition graph, a switch routing alphabet, and an overlap contraction-margin test. The layer proves that if selector switches are blocked, debt-free HC/MM absorbed, or routed into certified charts, and if every overlap edge preserves the strict contraction margin, then no uncontrolled OP5 exit remains on the finite certified halo atlas apart from explicitly declared leakage/debt or non-certified pieces. No global OP5 closure is claimed.
\end{abstract}

\section{Purpose and inherited state}

The layer removed rank-drop tunnels on certified photon-frame normal-collar pieces. The residual menu was reduced to selector switches, atlas overlaps, hidden leakage/debt, and non-certified pieces. This layer executes the first two items:
\[
\mathcal C_{\chi},\qquad \mathcal C_{\ov}.
\]
The printed operators remain renderer-safe:
\[
\HaloOp,\qquad \ExitOp,\qquad \Jexit .
\]
No lowercase display macro for these operators is used.

We work over a finite certified halo atlas
\[
\Atlas_{\gamma,K}^{\mathrm{loc}}=\{U_1,\ldots,U_M\}
\]
whose certified normal-collar pieces satisfy the determinant bound
\[
|W_j|\ge \epsr\epsh(1-\theta_j)>0.
\]
Thus the rank-drop channel is already empty on the certified pieces. The remaining question is whether a route can leave through a selector change or through a chart transition.

\section{Selector-switch register}

\begin{definition}[Finite selector-switch stratum]
Let
\[
\chi_j:U_j\longrightarrow \{0,1,\ldots,Q\}
\]
be the local OP5 route selector on a certified chart piece. The selector-switch stratum is
\[
\Ssel=
\bigcup_{j=1}^{M}\partial_{\chi}U_j,
\]
where \(\partial_{\chi}U_j\) is the finite union of codimension-one faces on which \(\chi_j\) changes value or ceases to be locally constant. The selector-switch register is
\[
\Cchi=
\{p\in \Ssel:
 p\in\partial\HaloOp_{\epsr,\epsh}^{\gamma,K},\ \chi_{\tun}(p)=1\}.
\]
\end{definition}

\begin{definition}[switch response alphabet]
A selector-switch point \(p\in\Cchi\) receives one of four conservative responses:
\[
\mathrm{resp}_{\chi}^{\mathrm{loc}}(p)\in
\{\block,\abs,\route,\res\}.
\]
The meanings are:
\begin{center}
\begin{tabular}{c|P{0.68\linewidth}}
\toprule
response & conservative meaning\\
\midrule
\(\block\) & the selector change is inadmissible for the OP5 tunnel gate; no exit is allowed;\\
\(\abs\) & the switch is absorbed by declared debt-free HC/MM bookkeeping, with \(\Bhcmm=0\);\\
\(\route\) & the switch enters a certified chart piece with preserved halo margin;\\
\(\res\) & the point remains an explicit residual and cannot be used for closure.\\
\bottomrule
\end{tabular}
\end{center}
\end{definition}

\begin{lemma}[Switch containment]
Assume every selector is locally constant on the interior of each certified collar piece. Then every selector-induced uncontrolled exit candidate lies in \(\Cchi\).
\end{lemma}

\begin{proof}
On the interior of a certified piece the selector does not change. Therefore selector-driven route creation cannot occur there. If a selector change occurs, it occurs on the switch stratum by definition. Intersecting with the halo boundary and the tunnel gate gives exactly \(\Cchi\).
\end{proof}

\begin{proposition}[Selector-switch closure criterion]
If
\[
\mathrm{resp}_{\chi}^{\mathrm{loc}}(p)\in\{\block,\abs,\route\}
\quad\text{for every }p\in\Cchi,
\]
and if each routed switch lands in a certified chart with unchanged debt-free bookkeeping, then selector switches create no uncontrolled OP5 exit on the certified atlas.
\end{proposition}

\begin{proof}
A blocked switch has no admissible exit. An absorbed switch has no hidden HC/MM debt by hypothesis. A routed switch remains inside the certified atlas and is again subject to the rank elimination and the overlap test below. Since no switch has response \(\res\), no selector-induced uncontrolled exit remains.
\end{proof}

\section{Atlas-overlap register}

\begin{definition}[Finite overlap transition graph]
The overlap graph is the finite directed graph
\[
\Gtrans=(V,E),
\qquad
V=\{1,\ldots,M\},
\]
where \((j,k)\in E\) exactly when the certified chart overlap
\[
U_{jk}=U_j\cap U_k
\]
is nonempty and may be crossed by an admissible OP5 route. On such an edge let
\[
\Phi_{jk}:U_{jk}\subset U_j\longrightarrow U_k
\]
be the declared transition map.
\end{definition}

\begin{definition}[Overlap contraction margin]
For an overlap edge \((j,k)\in E\), define the overlap margin
\[
\marg_{jk}^{\mathrm{loc}}
=
\widehat\lambda_j+\eta_j+\omega_{jk},
\]
where \(\widehat\lambda_j+\eta_j<1\) is the inherited chart contraction budget and \(\omega_{jk}\ge 0\) is the declared transition distortion of \(\Phi_{jk}\) in the two-channel halo norm
\[
\|(u,v)\|_{\HaloOp}=
\max\left\{\frac{|u|}{\epsr},\frac{|v|}{\epsh}\right\}.
\]
The overlap edge is contraction-preserving when
\[
\marg_{jk}^{\mathrm{loc}}<1.
\]
The overlap register is
\[
\Cov=
\{(j,k)\in E:\marg_{jk}^{\mathrm{loc}}\ge 1\}\cup
\{(j,k)\in E:\Phi_{jk}\text{ is undeclared}\}.
\]
\end{definition}

\begin{lemma}[Overlap preservation]
If \((j,k)\in E\) has \(\marg_{jk}^{\mathrm{loc}}<1\), then crossing the overlap from \(U_j\) to \(U_k\) preserves strict halo contraction and does not create an uncontrolled OP5 exit.
\end{lemma}

\begin{proof}
The inherited chart budget bounds the active chartwise projection debt by \(\widehat\lambda_j+\eta_j\). The transition contributes at most \(\omega_{jk}\) in the normalized two-channel halo norm. Hence the total budget after crossing the overlap is bounded by \(\marg_{jk}^{\mathrm{loc}}<1\). The crossing therefore remains inside the controlled halo atlas budget.
\end{proof}

\begin{proposition}[Atlas-overlap closure criterion]
If \(\Cov=\varnothing\), then atlas overlaps do not contribute uncontrolled OP5 exits on the finite certified atlas.
\end{proposition}

\begin{proof}
By definition, \(\Cov=\varnothing\) means every admissible overlap edge is declared and contraction-preserving. The previous lemma applies to every edge of the finite transition graph. Therefore no overlap crossing leaves the controlled halo budget.
\end{proof}

\section{First symbolic execution table}

For the current three-piece symbolic atlas inherited from the local pass, use the following conservative finite routing table. The entries are structural certificates, not physical measurements.

\begin{center}
\small
\begin{tabular}{c|c|c|P{0.39\linewidth}}
\toprule
selector face & location & response & reason\\
\midrule
\(S_{12}\) & \(U_1\leftrightarrow U_2\) & \(\route\) & same halo side, certified target\\
\(S_{23}\) & \(U_2\leftrightarrow U_3\) & \(\route\) & corner collar remains certified\\
\(S_{13}\) & \(U_1\leftrightarrow U_3\) & \(\block\) & non-adjacent selector jump not OP5-admissible\\
\(S_{\HC/\MM}\) & HC/MM bookkeeping face & \(\abs\) & declared debt-free absorption\\
\bottomrule
\end{tabular}
\end{center}

Thus no switch in the sample table has residual response.

The corresponding overlap margins are:
\begin{center}
\small
\begin{tabular}{c|c|c|c|P{0.28\linewidth}}
\toprule
edge & \(\widehat\lambda_j+\eta_j\) & \(\omega_{jk}\) & \(\marg_{jk}^{\mathrm{loc}}\) & status\\
\midrule
\((1,2)\) & \(5/8\) & \(1/16\) & \(11/16\) & preserved\\
\((2,3)\) & \(21/32\) & \(1/32\) & \(11/16\) & preserved\\
\((1,3)\) & \(25/32\) & \(1/32\) & \(13/16\) & preserved if routed; blocked as switch\\
\bottomrule
\end{tabular}
\end{center}

Therefore the sample overlap register is empty:
\[
\Cov=\varnothing
\]
for the declared transition graph.

\section{residual theorem}

\begin{theorem}[Selector-overlap execution theorem]
Assume the double-epsilon halo channel, the photon-frame halo-channel execution, the rank/switch register reduction, and the normal-collar rank elimination. Suppose furthermore that:
\begin{enumerate}
\item every selector switch in \(\Cchi\) has response \(\block\), \(\abs\), or \(\route\), never \(\res\);
\item every routed selector switch lands in a certified normal-collar chart;
\item every admissible overlap edge is declared and satisfies \(\marg_{jk}^{\mathrm{loc}}<1\);
\item hidden leakage and HC/MM debt vanish:
\[
\Lunr=0,
\qquad
\Bhcmm=0;
\]
\item no non-certified photon-frame halo piece remains in the active OP5 route menu.
\end{enumerate}
Then the finite certified photon-frame halo atlas has no uncontrolled OP5 exit. Equivalently, within this declared halo/tunnel grammar, every admissible near-kernel OP5 route is either blocked, debt-free absorbed, or routed back into a contraction-preserving certified HC-collar chart.
\end{theorem}

\begin{proof}
The determinant bound eliminates rank-drop exits on certified collar pieces. The selector-switch criterion eliminates selector-induced exits unless they are blocked, absorbed, or routed into certified pieces. The overlap criterion eliminates chart-transition exits because every admissible edge preserves strict contraction margin. The leakage and HC/MM-debt hypotheses rule out hidden bookkeeping exits. Finally, the absence of non-certified active pieces prevents escape through parts of the halo boundary not covered by the certified grammar. Hence the finite certified atlas contains no uncontrolled OP5 exit.
\end{proof}

\begin{remark}[Conservativity guard]
This theorem is not a global OP5 closure. It is a finite-atlas closure criterion under the OP5=HC-cone working identification and under the declared double-epsilon halo grammar. A global closure would additionally require proving that every admissible OP5 near-kernel route enters this certified finite atlas and that the non-certified residual class is empty in the full active-core setting.
\end{remark}

\section{Resulting next step}

The execution leaves a sharply stated final gap. To promote the local halo/tunnel analysis toward a real OP5 closure statement, the next layer must prove the capture/completeness clause:
\[
\mathcal R_{\OP5}^{\adm,\mathrm{near}
K}
\subseteq
\bigcup_{j=1}^{M}U_j^{\cert}
\]
for the relevant photon-frame HC-collar readouts, together with \(\Lunr=0\) and \(\Bhcmm=0\). This is the natural target.


% --- end integrated body for r56 ---

\clearpage
\section*{Integrated Delta Layer: Dixon-Carrier Coding}
\addcontentsline{toc}{section}{Integrated Delta Layer: Dixon-Carrier Coding}
% --- local macro refresh for r57 ---
\providecommand{\HC}{}
\renewcommand{\HC}{\mathrm{HC}}
\providecommand{\MM}{}
\renewcommand{\MM}{\mathrm{MM}}
\providecommand{\OP}{}
\renewcommand{\OP}{\mathrm{OP}}
\providecommand{\adm}{}
\renewcommand{\adm}{\mathrm{adm}}
\providecommand{\near}{}
\renewcommand{\near}{\mathrm{near}}
\providecommand{\cert}{}
\renewcommand{\cert}{\mathrm{cert}}
\providecommand{\tun}{}
\renewcommand{\tun}{\mathrm{tun}}
\providecommand{\unr}{}
\renewcommand{\unr}{\mathrm{unr}}
\providecommand{\ov}{}
\renewcommand{\ov}{\mathrm{ov}}
\providecommand{\rk}{}
\renewcommand{\rk}{\mathrm{rk}}
\providecommand{\capreg}{}
\renewcommand{\capreg}{\mathrm{cap}}
\providecommand{\HaloOp}{}
\renewcommand{\HaloOp}{\mathsf{Halo}}
\providecommand{\ExitOp}{}
\renewcommand{\ExitOp}{\mathsf{Exit}}
\providecommand{\Atlas}{}
\renewcommand{\Atlas}{\mathfrak A}
\providecommand{\epsr}{}
\renewcommand{\epsr}{\varepsilon_r}
\providecommand{\epsh}{}
\renewcommand{\epsh}{\varepsilon_h}
\providecommand{\Xop}{}
\renewcommand{\Xop}{\mathcal X_{\OP 5}^{\adm}}
\providecommand{\Rop}{}
\renewcommand{\Rop}{\mathcal R_{\OP 5}^{\adm,\near K}}
\providecommand{\rhoR}{}
\renewcommand{\rhoR}{\rho_{r}^{\gamma}}
\providecommand{\rhoH}{}
\renewcommand{\rhoH}{\rho_{h}^{\gamma}}
\providecommand{\Lunr}{}
\renewcommand{\Lunr}{L_{\unr}}
\providecommand{\Bhcmm}{}
\renewcommand{\Bhcmm}{B_{\HC/\MM}}
\providecommand{\Ncollar}{}
\renewcommand{\Ncollar}{\mathsf N_{\HC}^{\gamma,K}}
\providecommand{\Capt}{}
\renewcommand{\Capt}{\mathsf{Cap}}
\providecommand{\Comp}{}
\renewcommand{\Comp}{\mathsf{Comp}}
\providecommand{\Route}{}
\renewcommand{\Route}{\mathsf{Route}}
\providecommand{\Gate}{}
\renewcommand{\Gate}{\mathsf{Gate}}
\providecommand{\Debt}{}
\renewcommand{\Debt}{\mathsf{Debt}}
\providecommand{\Leak}{}
\renewcommand{\Leak}{\mathsf{Leak}}
\providecommand{\rank}{}
\renewcommand{\rank}{\operatorname{rank}}
\providecommand{\supp}{}
\renewcommand{\supp}{\operatorname{supp}}
\providecommand{\dist}{}
\renewcommand{\dist}{\operatorname{dist}}
% --- begin integrated body for r57 ---



\begin{abstract}
This delta layer executes the capture/completeness step left open by the local pass. It defines a finite route-capture certificate for admissible near-kernel OP5 readouts under the working identification OP5 = photon-frame HC-cone route. The layer separates three tasks: capture into the double-epsilon halo, assignment to a certified normal-collar chart, and debt-free compatibility with selector/overlap routing. The result is a conditional but explicit completeness theorem: if the finite capture certificate is exhaustive, if every captured branch lands in the local proof chain certified atlas, and if hidden leakage and HC/MM debt vanish, then the declared near-kernel OP5 route menu has no uncontrolled exit. This is still not a global OP5 proof outside the declared working grammar.
\end{abstract}

\section{Purpose and inherited state}

The local proof-chain layers transformed the OP5 near-kernel route problem into a finite halo/tunnel problem. The current state is:
\begin{enumerate}[label=(\roman*)]
\item introduced the double-epsilon halo channel;
\item instantiated the photon-frame channels \(\rhoR\) and \(\rhoH\) and the two-channel exit map;
\item reduced rank-drop and selector-switch search to explicit registers;
\item eliminated rank-drop on certified normal-collar pieces;
\item routed selector switches and showed contraction-preserving atlas overlaps in the sample finite atlas.
\end{enumerate}
The remaining gap is a capture/completeness statement:
\[
\Rop
\subseteq
\bigcup_{j=1}^{M} U_j^{\cert},
\qquad
\Lunr=0,
\qquad
\Bhcmm=0 .
\]
This layer gives the certificate grammar for exactly that statement.

All displayed operators use renderer-safe notation:
\[
\HaloOp,
\qquad
\ExitOp,
\qquad
\Capt .
\]

\section{Finite route-capture certificate}

\begin{definition}[Near-kernel OP5 route menu]
Let \(\Rop\) be the declared set of admissible near-kernel OP5 readouts after photon-frame HC-collar normalization. A route \(r\in\Rop\) is represented by a finite branch word
\[
\beta(r)=(\sigma_0,\sigma_1,\ldots,\sigma_\ell)
\]
whose symbols record, in order, HC-collar entry, possible selector face, possible overlap transition, and terminal certified chart candidate. The word is \emph{admissible} only if every symbol is admitted by the inherited switch/overlap tables.
\end{definition}

\begin{definition}[capture witness]
An capture witness for a route \(r\in\Rop\) is a tuple
\[
\mathfrak c_{r}^{\mathrm{loc}}
=
(n_r,j_r,\beta_r,
\delta_{r}^{\HaloOp},
\ell_r,b_r),
\]
where:
\begin{enumerate}[label=(\alph*)]
\item \(n_r\in\mathbb N\) is the finite number of HC/Hilbert normalizing steps;
\item \(j_r\in\{1,\ldots,M\}\) is the certified target chart;
\item \(\beta_r\) is the finite the local pass-admissible selector/overlap word;
\item \(\delta_{r}^{\HaloOp}\) is the strict halo margin;
\item \(\ell_r\) is the hidden leakage register;
\item \(b_r\) is the hidden HC/MM debt register.
\end{enumerate}
It is valid when
\[
\rhoR(R_{\HC}^{n_r}r)\le \epsr-\delta_{r}^{\HaloOp},
\qquad
\rhoH(R_{\HC}^{n_r}r)\le \epsh-\delta_{r}^{\HaloOp},
\]
\[
R_{\HC}^{n_r}r\in U_{j_r}^{\cert},
\qquad
\delta_{r}^{\HaloOp}>0,
\qquad
\ell_r=0,
\qquad
b_r=0.
\]
\end{definition}

\begin{definition}[Finite capture certificate]
A finite capture certificate is a finite family
\[
\mathfrak C_{\capreg}^{\mathrm{loc}}
=
\{\mathfrak c_{a}^{\mathrm{loc}}:a\in A_{\mathrm{loc}}\}
\]
together with a surjective classification map
\[
\pi_{\capreg}^{\mathrm{loc}}:\Rop\longrightarrow A_{\mathrm{loc}}.
\]
For a route \(r\), the witness assigned to \(a=\pi_{\capreg}^{\mathrm{loc}}(r)\) must be valid for every route in the class \((\pi_{\capreg}^{\mathrm{loc}})^{-1}(a)\).
\end{definition}

\begin{remark}[Why this is the right bottleneck]
Without the capture certificate, OP5 remains a many-route problem. With the certificate, every admissible near-kernel route is compressed into one of finitely many certified classes. The remaining proof is then finite: check capture margin, chart certification, selector/overlap admissibility, leakage, and HC/MM debt for each class.
\end{remark}

\section{Capture tests}

\begin{definition}[capture test stack]
For each capture class \(a\in A_{\mathrm{loc}}\), define four Boolean tests:
\[
\Capt_{\HaloOp}(a),
\qquad
\Capt_{\Atlas}(a),
\qquad
\Capt_{\Route}(a),
\qquad
\Capt_{\Debt}(a).
\]
They mean respectively:
\begin{align*}
\Capt_{\HaloOp}(a)=1
&\Longleftrightarrow
\rhoR\le\epsr-\delta_a^{\HaloOp}
\text{ and }
\rhoH\le\epsh-\delta_a^{\HaloOp}
\text{ after }n_a\text{ steps},\\
\Capt_{\Atlas}(a)=1
&\Longleftrightarrow
R_{\HC}^{n_a}(r)\in U_{j_a}^{\cert}
\text{ for every }r\text{ in class }a,\\
\Capt_{\Route}(a)=1
&\Longleftrightarrow
\beta_a\text{ uses only the local pass-admissible selector/overlap moves},\\
\Capt_{\Debt}(a)=1
&\Longleftrightarrow
\ell_a=0\text{ and }b_a=0.
\end{align*}
The class is capture-complete when all four tests equal one.
\end{definition}

\begin{lemma}[Classwise capture implies chartwise capture]
If \(\Capt_{\HaloOp}(a)=\Capt_{\Atlas}(a)=1\), then every route in class \(a\) enters the certified double-epsilon halo atlas after at most \(n_a\) normalizing steps.
\end{lemma}

\begin{proof}
The first test gives strict entry into the double-epsilon halo, with positive margin \(\delta_a^{\HaloOp}\). The second test identifies the target as an certified normal-collar chart. Therefore all routes in the class enter the certified halo atlas.
\end{proof}

\begin{lemma}[Route and debt tests eliminate hidden exits]
If \(\Capt_{\Route}(a)=\Capt_{\Debt}(a)=1\), then class \(a\) contains no unregistered selector/overlap exit and no hidden leakage or HC/MM debt.
\end{lemma}

\begin{proof}
The route test restricts the branch word to selector and overlap transitions already routed by the local pass. The debt test sets both hidden registers to zero. Thus the class cannot leave through an undeclared selector/overlap move or through hidden bookkeeping.
\end{proof}

\section{First symbolic capture table}

The following table records the first finite symbolic certificate compatible with the local proof chain. It is a structural certificate table rather than a measurement table.

\begin{center}
\small
\begin{tabular}{c|c|c|c|c|c|P{0.23\linewidth}}
\toprule
class & branch type & \(n_a\) & target & margin & debt & status\\
\midrule
\(A_1\) & collar interior & 2 & \(U_1^{\cert}\) & \(1/8\) & 0 & captured\\
\(A_2\) & adjacent switch \(S_{12}\) & 2 & \(U_2^{\cert}\) & \(3/32\) & 0 & routed/captured\\
\(A_3\) & adjacent switch \(S_{23}\) & 2 & \(U_3^{\cert}\) & \(1/16\) & 0 & routed/captured\\
\(A_4\) & non-adjacent switch \(S_{13}\) & -- & -- & -- & -- & blocked\\
\(A_5\) & HC/MM face & 2 & \(U_2^{\cert}\) & \(1/16\) & 0 & absorbed/captured\\
\(A_6\) & overlap edge & 2 & \(U_j^{\cert}\cup U_k^{\cert}\) & \(1/16\) & 0 & overlap-preserved\\
\bottomrule
\end{tabular}
\end{center}

The active captured classes are
\[
A_{\mathrm{act}}^{\mathrm{loc}}=\{A_1,A_2,A_3,A_5,A_6\}.
\]
The blocked non-adjacent selector class \(A_4\) does not contribute an admissible OP5 route. Thus the symbolic residual class is
\[
A_{\mathrm{res}}^{\mathrm{loc}}=\varnothing
\]
provided the classification map \(\pi_{\capreg}^{\mathrm{loc}}\) is exhaustive for the declared near-kernel route menu.

\begin{proposition}[Sample certificate status]
For the symbolic capture table, every active class has strict positive halo margin, zero leakage, zero HC/MM debt, and a certified the local proof chain target chart. The only non-captured class is blocked by the inherited selector rule. Hence the sample table has no residual active capture class.
\end{proposition}

\begin{proof}
The entries list positive margins for \(A_1,A_2,A_3,A_5,A_6\), debt value zero, and certified targets. Class \(A_4\) is not an active OP5 route because blocks the non-adjacent selector jump. Therefore no active class remains residual.
\end{proof}

\section{Completeness theorem}

\begin{theorem}[capture/completeness theorem]
Assume the double-epsilon halo channel, the photon-frame channel execution, the register reduction, the normal-collar rank elimination, and the selector/overlap execution. Suppose there exists a finite capture certificate \(\mathfrak C_{\capreg}^{\mathrm{loc}}\) and an exhaustive classification map
\[
\pi_{\capreg}^{\mathrm{loc}}:\Rop\longrightarrow A_{\mathrm{loc}}
\]
such that every active class is capture-complete and every inactive class is blocked by an inherited OP5-admissibility rule. Then
\[
\Rop
\subseteq
\bigcup_{j=1}^{M}U_j^{\cert}
\]
after finitely many HC/Hilbert normalizing steps, with
\[
\Lunr=0,
\qquad
\Bhcmm=0.
\]
Consequently, within the declared photon-frame HC-halo grammar, the near-kernel OP5 route menu has no uncontrolled exit.
\end{theorem}

\begin{proof}
Let \(r\in\Rop\). By exhaustiveness, \(r\) belongs to some class \(a=\pi_{\capreg}^{\mathrm{loc}}(r)\). If the class is inactive, it is blocked and hence not an active OP5 exit. If the class is active, capture-completeness gives four facts: strict double-epsilon halo entry, certified target chart membership, admissible selector/overlap routing, and zero leakage/debt. The theorem eliminates rank-drop on the certified chart. The theorem eliminates selector and overlap exits. The zero-register clause eliminates hidden leakage and HC/MM debt. Hence \(r\) is blocked, absorbed, or routed into a certified contraction-preserving chart. Since \(r\) was arbitrary, the inclusion and the no-uncontrolled-exit claim follow.
\end{proof}

\begin{corollary}[Local OP5 closure inside the declared grammar]
Under the hypotheses of the theorem and under the working identification OP5 = photon-frame HC-cone route, the near-kernel OP5 branch of the declared finite halo/tunnel grammar is locally closed.
\end{corollary}

\begin{proof}
The theorem states that every admissible near-kernel route is captured by the certified atlas or blocked, with no hidden leakage/debt. This is precisely local closure in the declared grammar.
\end{proof}

\begin{remark}[No overclaim]
This is not a global OP5 proof in the ordinary external sense. It is the strongest local closure statement presently justified by the local proof chain: if OP5 is correctly read as the photon-frame HC-cone route problem, and if the finite capture certificate is exhaustive for the relevant active-core route menu, then the near-kernel OP5 branch has no uncontrolled exit. A full release-level OP5 theorem would still need this certificate integrated into the monolithic Active Core and cross-checked against every OP5 occurrence in the historical witness.
\end{remark}

\section{Audit target for the next layer}

The layer turns the remaining work into an audit problem. The next meaningful step is not to add another route mechanism, but to verify exhaustiveness:
\[
\supp(\Rop)
=
\bigcup_{a\in A_{\mathrm{loc}}}(\pi_{\capreg}^{\mathrm{loc}})^{-1}(a),
\]
with no unclassified active branch. If the monolithic Active Core confirms that equality, then can promote the result from a local capture/completeness certificate to an OP5 closure-candidate theorem for the full active-core route menu.


% --- end integrated body for r57 ---

\clearpage
\section*{Integrated Delta Layer: Dixon-Image Audit}
\addcontentsline{toc}{section}{Integrated Delta Layer: Dixon-Image Audit}
% --- local macro refresh for r58 ---
\providecommand{\HC}{}
\renewcommand{\HC}{\mathrm{HC}}
\providecommand{\MM}{}
\renewcommand{\MM}{\mathrm{MM}}
\providecommand{\OP}{}
\renewcommand{\OP}{\mathrm{OP}}
\providecommand{\adm}{}
\renewcommand{\adm}{\mathrm{adm}}
\providecommand{\near}{}
\renewcommand{\near}{\mathrm{near}}
\providecommand{\cert}{}
\renewcommand{\cert}{\mathrm{cert}}
\providecommand{\unr}{}
\renewcommand{\unr}{\mathrm{unr}}
\providecommand{\rk}{}
\renewcommand{\rk}{\mathrm{rk}}
\providecommand{\Dix}{}
\renewcommand{\Dix}{\mathrm{Dix}}
\providecommand{\HS}{}
\renewcommand{\HS}{\mathrm{HS}}
\providecommand{\Carrier}{}
\renewcommand{\Carrier}{\mathsf{Car}}
\providecommand{\Mask}{}
\renewcommand{\Mask}{\mathsf{Mask}}
\providecommand{\Code}{}
\renewcommand{\Code}{\mathsf{Code}}
\providecommand{\Gate}{}
\renewcommand{\Gate}{\mathsf{Gate}}
\providecommand{\Leak}{}
\renewcommand{\Leak}{\mathsf{Leak}}
\providecommand{\Debt}{}
\renewcommand{\Debt}{\mathsf{Debt}}
\providecommand{\Route}{}
\renewcommand{\Route}{\mathsf{Route}}
\providecommand{\HaloOp}{}
\renewcommand{\HaloOp}{\mathsf{Halo}}
\providecommand{\Atlas}{}
\renewcommand{\Atlas}{\mathfrak A}
\providecommand{\Rop}{}
\renewcommand{\Rop}{\mathcal R_{\OP 5}^{\adm,\near K}}
\providecommand{\Lunr}{}
\renewcommand{\Lunr}{L_{\unr}}
\providecommand{\Bhcmm}{}
\renewcommand{\Bhcmm}{B_{\HC/\MM}}
\providecommand{\C}{}
\renewcommand{\C}{\mathbb C}
\providecommand{\Hh}{}
\renewcommand{\Hh}{\mathbb H}
\providecommand{\Oo}{}
\renewcommand{\Oo}{\mathbb O}
\providecommand{\One}{}
\renewcommand{\One}{\mathbf 1}
\providecommand{\MC}{}
\renewcommand{\MC}{\mathcal M_{16}}
\providecommand{\PC}{}
\renewcommand{\PC}{\widehat P_{C_7}}
\providecommand{\Qp}{}
\renewcommand{\Qp}{Q_{\perp}}
\providecommand{\HDix}{}
\renewcommand{\HDix}{\mathcal H_{\Dix}}
\providecommand{\HFano}{}
\renewcommand{\HFano}{\mathcal F^{(3)}}
\providecommand{\HsFour}{}
\renewcommand{\HsFour}{\mathcal H^{(4)}_{16}}
\providecommand{\HsFive}{}
\renewcommand{\HsFive}{\mathcal H^{(5)}_{32}}
\providecommand{\HsSix}{}
\renewcommand{\HsSix}{\mathcal H^{(6)}_{64}}
\providecommand{\supp}{}
\renewcommand{\supp}{\operatorname{supp}}
\providecommand{\rank}{}
\renewcommand{\rank}{\operatorname{rank}}
\providecommand{\im}{}
\renewcommand{\im}{\operatorname{im}}
\providecommand{\diag}{}
\renewcommand{\diag}{\operatorname{diag}}
\providecommand{\Aut}{}
\renewcommand{\Aut}{\operatorname{Aut}}
\providecommand{\End}{}
\renewcommand{\End}{\operatorname{End}}
\providecommand{\Tr}{}
\renewcommand{\Tr}{\operatorname{Tr}}
% --- begin integrated body for r58 ---



\begin{abstract}
This delta layer executes the HS(6) Dixon-coding map suggested after the capture/completeness packet. The layer separates two clocks which must not be conflated: the carrier clock
\(C_n=2^n\), and the activation-mask clock \(|\{0,1\}^{C_n}|=2^{C_n}\). In this grammar, HS(5) has a 32-channel carrier with a \(2^{32}\)-state activation mask, while HS(6) has a 64-channel carrier whose activation mask has \(2^{64}\) states. The 64-channel carrier is mapped to the Dixon carrier \(D=\mathbb C\otimes\mathbb H\otimes\mathbb O\). The layer then relates the HS(4) meta-septum projector \(\widehat P_{C_7}\subset\mathcal M_{16}\), the HS(5) two-sided duplication, and the HS(6) Dixon carrier to the OP5 halo/capture chain. No claim is made that the Dixon algebra is \(2^{64}\)-dimensional; only its Boolean activation code has \(2^{64}\) possible masks.
\end{abstract}

\section{Purpose and inherited state}

The local proof chain sequence reduced the near-kernel OP5 route problem to a finite photon-frame HC-halo capture problem. Its current decisive condition is exhaustiveness of the capture map
\[
\pi_{\mathrm{cap}}^{\mathrm{loc}}:\Rop\longrightarrow A_{\mathrm{loc}},
\]
with \(\Lunr=0\) and \(\Bhcmm=0\). The present layer records the next structural interpretation of higher Hilbert shifts. It does not add a new OP5 route mechanism. It provides a finite coding target for the already captured routes.

The motivating observation is that the HS tower naturally admits two readings:
\[
\Carrier(n)=2^n,
\qquad
\Mask(n)=\{0,1\}^{\Carrier(n)},
\qquad
|\Mask(n)|=2^{2^n}.
\]
Thus the phrase ``up to \(2^{32}\), then \(2^{64}\)'' is correct when it refers to activation masks at carrier sizes \(32\) and \(64\), not to the algebra dimension itself.

\section{Two-clock Hilbert-shift grammar}

\begin{definition}[Carrier clock and mask clock]
For the Hilbert-shift tower define
\[
C_n:=2^n
\]
as the carrier size of the \(n\)-th shift, and define the Boolean activation mask space
\[
\Mask_n:=\{0,1\}^{C_n}.
\]
Then
\[
|\Mask_n|=2^{C_n}=2^{2^n}.
\]
We call \(C_n\) the \emph{carrier clock} and \(|\Mask_n|\) the \emph{mask clock}.
\end{definition}

\begin{remark}[Correction of the common ambiguity]
The statement ``Dixon is \(2^{64}\)'' is not used as a dimension statement. The structurally safe statement is
\[
\dim_{\mathbb R}(\C\otimes\Hh\otimes\Oo)=64,
\qquad
|\{0,1\}^{64}|=2^{64}.
\]
Thus the Dixon carrier has 64 real channels, and its full Boolean activation grammar has \(2^{64}\) masks.
\end{remark}

\begin{center}
\small
\begin{tabular}{c|c|c|P{0.42\linewidth}}
\toprule
shift & carrier size \(C_n\) & mask count \(|\Mask_n|\) & structural reading\\
\midrule
\(\HS(0)\) & 1 & 2 & radial kernel target / unshifted collapse readout\\
\(\HS(1)\) & 2 & 4 & HC collar and periapsis split\\
\(\HS(2)\) & 4 & 16 & stabilized orbit / contraction chart\\
\(\HS(3)\) & 8 & \(2^8\) & Fano/octonion-compatible local incidence fibre\\
\(\HS(4)\) & 16 & \(2^{16}\) & meta-septum \(\MC/\PC\) leakage test\\
\(\HS(5)\) & 32 & \(2^{32}\) & two-sided front/back carrier duplication\\
\(\HS(6)\) & 64 & \(2^{64}\) & Dixon carrier \(\C\otimes\Hh\otimes\Oo\) and 64-bit activation code\\
\bottomrule
\end{tabular}
\end{center}

\section{HS(4) meta-septum leakage target}

\begin{definition}[HS(4) meta-septum carrier]
Let \(\HsFour\) be the 16-channel HS(4) carrier containing the inherited meta-septum sector selected by
\[
\PC:\HsFour\longrightarrow \HsFour,
\qquad
\PC^2=\PC.
\]
Write
\[
Q:=\PC,
\qquad
\Qp:=\One-Q.
\]
For the inherited internal matrix \(\MC\), define the admissible block and mixed leakage block by
\[
\mathcal M_{CC}:=Q\MC Q,
\]
\[
\Leak_{C_7}^{(4)}(\MC)
:=
\Qp\MC Q+Q\MC\Qp.
\]
The HS(4) carrier is meta-septum closed when
\[
\Leak_{C_7}^{(4)}(\MC)=0.
\]
\end{definition}

\begin{lemma}[Commutator form of HS(4) closure]
If \(Q=Q^2\) and \(\Qp=\One-Q\), then vanishing of both mixed blocks is equivalent to block preservation of the \(C_7\)-sector. In particular the closure condition can be written as
\[
[\MC,Q]=0,
\]
provided the two mixed blocks are tested as independent leakage registers.
\end{lemma}

\begin{proof}
With respect to \(\HsFour=Q\HsFour\oplus\Qp\HsFour\), the operator \(\MC\) decomposes into four blocks. The two diagonal blocks preserve the selected and complementary sectors. The two off-diagonal blocks are \(\Qp\MC Q\) and \(Q\MC\Qp\). If they vanish, \(\MC\) is block diagonal and commutes with \(Q\). Conversely, if \([\MC,Q]=0\), multiplying on the left and right by \(Q\) and \(\Qp\) gives vanishing of the off-diagonal blocks. Hence the mixed leakage register is zero exactly in the block-preserving case.
\end{proof}

\begin{remark}[Relation to HS(4) danger]
The HS(4) carrier is the first level at which a naive 16-channel extension may behave like a sedenion-type expansion. This layer does not use raw sedenions as a free algebra. It uses \(\MC/\PC\) as a gated 16-channel envelope and records all off-sector motion as explicit leakage.
\end{remark}

\section{HS(5) two-sided duplication}

\begin{definition}[Two-sided HS(5) carrier]
The HS(5) carrier is the front/back duplication of the HS(4) envelope:
\[
\HsFive
:=
(\HsFour)_{\mathrm{front}}\oplus(\HsFour)_{\mathrm{back}}.
\]
Its projector is
\[
\PC^{(5)}
:=
\begin{pmatrix}
(\PC)_{\mathrm{front}}&0\\
0&(\PC)_{\mathrm{back}}
\end{pmatrix}.
\]
A two-sided HS(5) state is written
\[
\Psi^{(5)}=(\Psi_{\mathrm{front}},\Psi_{\mathrm{back}}),
\]
where the front side carries readout/frame/language information and the back side carries structure/action/field information.
\end{definition}

\begin{definition}[HS(5) cross-side debt]
Let \(\mathcal T^{(5)}\in\End(\HsFive)\). Define
\[
\Debt_{\mathrm{cross}}^{(5)}(\mathcal T^{(5)})
:=
(\One-\PC^{(5)})\mathcal T^{(5)}\PC^{(5)}
+
\PC^{(5)}\mathcal T^{(5)}(\One-\PC^{(5)}).
\]
The HS(5) two-sided carrier is debt-free when this register vanishes and all front/back comparison maps are mediated by declared comparison tensors or selectors.
\end{definition}

\begin{proposition}[HS(5) as safe bridge]
If the front and back HS(4) sectors are separately meta-septum closed and the cross-side debt register vanishes, then the HS(5) carrier transports the HS(3) Fano-compatible fibre through the two-sided readout grammar without adding an uncontrolled off-sector channel.
\end{proposition}

\begin{proof}
Separate meta-septum closure prevents leakage within each 16-channel side. Vanishing cross-side debt prevents an admissible selected component on one side from being pushed into a complementary component on the other side. Since two-sided comparisons are only admitted through declared selectors or comparison tensors, every surviving front/back transition is registered. Hence no uncontrolled off-sector channel is introduced at HS(5).
\end{proof}

\section{HS(6) Dixon carrier and 64-bit activation code}

\begin{definition}[Dixon carrier]
The HS(6) Dixon carrier is the 64-dimensional real carrier
\[
\HsSix
\cong
\HDix
:=
\C\otimes\Hh\otimes\Oo.
\]
A basis channel is indexed by a triple
\[
\alpha=(a,b,c),
\qquad
 a\in\{0,1\},\quad b\in\{0,1,2,3\},\quad c\in\{0,1,\ldots,7\}.
\]
Thus
\[
|\{\alpha\}|=2\cdot4\cdot8=64.
\]
\end{definition}

\begin{definition}[Dixon activation code]
The Dixon activation code of a captured OP5 route \(r\in\Rop\) is a Boolean mask
\[
\chi_{\Dix}(r)
=
(\chi_{\alpha}(r))_{\alpha\in\{0,1\}\times\{0,1,2,3\}\times\{0,\ldots,7\}}
\in
\{0,1\}^{64}.
\]
The full Dixon activation space therefore has
\[
|\Code_{\Dix}|=2^{64}
\]
possible masks. A mask is admissible only if it respects the inherited capture class, the HS(4) meta-septum gate, and the HS(5) two-sided debt-free condition.
\end{definition}

\begin{definition}[HS(6) coding map]
Define the conditional route-to-Dixon coding map
\[
\Xi_{\Dix}^{\mathrm{loc}}:
\Rop
\longrightarrow
\{0,1\}^{64}\cup\{\mathsf{blocked},\mathsf{debt},\mathsf{leak}\}
\]
as follows. For \(r\in\Rop\):
\begin{enumerate}[label=(\alph*)]
\item if \(r\) is blocked by the selector table, set \(\Xi_{\Dix}^{\mathrm{loc}}(r)=\mathsf{blocked}\);
\item if \(\Leak_{C_7}^{(4)}\ne0\), set \(\Xi_{\Dix}^{\mathrm{loc}}(r)=\mathsf{leak}\);
\item if \(\Debt_{\mathrm{cross}}^{(5)}\ne0\), set \(\Xi_{\Dix}^{\mathrm{loc}}(r)=\mathsf{debt}\);
\item otherwise assign the 64-bit mask \(\chi_{\Dix}(r)\) determined by its captured chart, selector word, front/back side, and Fano/octonion channel.
\end{enumerate}
\end{definition}

\section{Compatibility with the capture theorem}

\begin{lemma}[Captured routes admit finite Dixon codes]
Assume the capture/completeness theorem and the HS(4)--HS(5) zero-register conditions
\[
\Leak_{C_7}^{(4)}=0,
\qquad
\Debt_{\mathrm{cross}}^{(5)}=0.
\]
Then every active captured route \(r\in\Rop\) has a well-defined Dixon activation code \(\chi_{\Dix}(r)\in\{0,1\}^{64}\).
\end{lemma}

\begin{proof}
By the local pass, each active route is captured in a certified normal-collar chart after finitely many HC/Hilbert steps, with \(\Lunr=0\) and \(\Bhcmm=0\). The HS(4) condition prevents the Fano/meta-septum component from leaking into the complementary 16-channel sector. The HS(5) condition prevents two-sided front/back transport from creating unregistered debt. Therefore the remaining channel data are finite and live inside the HS(6) Dixon carrier. The 64 indexed basis channels define the Boolean mask.
\end{proof}

\begin{proposition}[Tunnel-to-Dixon reduction]
Under the hypotheses of the previous lemma, every tunnel/capture class is assigned one of four statuses:
\[
\mathsf{blocked},
\qquad
\mathsf{leak},
\qquad
\mathsf{debt},
\qquad
\chi_{\Dix}\in\{0,1\}^{64}.
\]
Moreover, an uncontrolled HS(6) tunnel can occur only through a nonzero HS(4) leakage register, a nonzero HS(5) cross-side debt register, or a nonadmissible Dixon mask.
\end{proposition}

\begin{proof}
The coding map definition exhausts the cases. Blocked routes are not active OP5 exits. Nonzero HS(4) leakage and HS(5) debt are explicitly recorded, not hidden. If both vanish, the route is represented by a finite 64-bit Dixon mask. Hence any uncontrolled tunnel must be visible as leak, debt, or failure of Dixon-mask admissibility.
\end{proof}

\begin{theorem}[HS(6) Dixon coding theorem]
Assume the local proof chain photon-frame HC-halo capture chain. Suppose further that
\[
\Leak_{C_7}^{(4)}=0,
\qquad
\Debt_{\mathrm{cross}}^{(5)}=0,
\qquad
\Lunr=0,
\qquad
\Bhcmm=0.
\]
Then every declared admissible near-kernel OP5 route in the capture menu is either blocked by an inherited selector rule or encoded by an admissible 64-bit Dixon activation mask:
\[
\Xi_{\Dix}^{\mathrm{loc}}(r)
\in
\{\mathsf{blocked}\}\cup\{0,1\}^{64}.
\]
Consequently, within the declared HS(0)--HS(6) grammar, the remaining OP5 tunnel audit becomes a finite Dixon-mask admissibility audit rather than an unbounded route search.
\end{theorem}

\begin{proof}
The theorem captures all active near-kernel routes into the finite certified halo atlas and eliminates hidden leakage and HC/MM debt. The HS(4) condition eliminates meta-septum off-sector leakage. The HS(5) condition eliminates front/back cross-side debt. Thus every route left after selector blocking has finite chart, selector, side, and Fano/octonion channel data. These data select a subset of the 64 Dixon carrier channels, i.e. a Boolean mask in \(\{0,1\}^{64}\). Therefore the audit is finite at the level of mask admissibility.
\end{proof}

\begin{corollary}[What would be needed next]
A release-level OP5 closure candidate can now be pursued by proving that every Dixon mask actually produced by the capture map is admissible, or else blocked/absorbed by explicit HC/MM bookkeeping:
\[
\im(\Xi_{\Dix}^{\mathrm{loc}})
\subseteq
\Code_{\Dix}^{\adm}
\cup\{\mathsf{blocked}\},
\]
with no leak and no debt registers.
\end{corollary}

\section{First symbolic Dixon status table}

The following table records the first symbolic status assignment. It is an integration table for the route grammar, not a numerical particle-spectrum table.

\begin{center}
\scriptsize
\setlength{\tabcolsep}{2pt}
\begin{tabular}{c|c|c|c|c|P{0.18\linewidth}}
\toprule
class & HS(4) leakage & HS(5) debt & Dixon carrier & mask status & status\\
\midrule
\(A_1\) collar interior & 0 & 0 & 64 & admissible pending full mask audit & Dixon-coded\\
\(A_2\) \(S_{12}\) routed & 0 & 0 & 64 & admissible pending full mask audit & Dixon-coded\\
\(A_3\) \(S_{23}\) routed & 0 & 0 & 64 & admissible pending full mask audit & Dixon-coded\\
\(A_4\) \(S_{13}\) blocked & -- & -- & -- & -- & blocked\\
\(A_5\) HC/MM face & 0 & 0 & 64 & absorbed/coded if debt-free & coded/absorbed\\
\(A_6\) overlap edge & 0 & 0 & 64 & overlap-preserved & Dixon-coded\\
\bottomrule
\end{tabular}
\end{center}

\begin{remark}[No spectrum or Standard Model overclaim]
The Dixon carrier is used here as a finite route-code carrier. This layer does not derive the Standard Model spectrum, CKM/PMNS values, quark masses, QCD confinement, or OP1/OP4. It only states how captured OP5 halo routes can be mapped into a finite 64-channel Dixon activation grammar.
\end{remark}

\section{Audit target for}

The next meaningful layer is not a new shift. It is the image audit of the coding map:
\[
\im(\Xi_{\Dix}^{\mathrm{loc}})
\stackrel{?}{\subseteq}
\Code_{\Dix}^{\adm}\cup\{\mathsf{blocked}\}.
\]
Equivalently, should enumerate the Dixon-mask classes actually reachable from the capture table and classify each as blocked, HC/MM absorbed, atlas routed, Dixon-admissible, or residual-open. If that image audit is empty on residual-open masks, then the OP5 closure candidate becomes substantially stronger inside the declared HS(0)--HS(6) grammar.


% --- end integrated body for r58 ---

\clearpage
\section*{Integrated Delta Layer: Route-Exhaustiveness Audit}
\addcontentsline{toc}{section}{Integrated Delta Layer: Route-Exhaustiveness Audit}
% --- local macro refresh for r59 ---
\providecommand{\HC}{}
\renewcommand{\HC}{\mathrm{HC}}
\providecommand{\MM}{}
\renewcommand{\MM}{\mathrm{MM}}
\providecommand{\OP}{}
\renewcommand{\OP}{\mathrm{OP}}
\providecommand{\adm}{}
\renewcommand{\adm}{\mathrm{adm}}
\providecommand{\near}{}
\renewcommand{\near}{\mathrm{near}}
\providecommand{\cert}{}
\renewcommand{\cert}{\mathrm{cert}}
\providecommand{\unr}{}
\renewcommand{\unr}{\mathrm{unr}}
\providecommand{\Dix}{}
\renewcommand{\Dix}{\mathrm{Dix}}
\providecommand{\HS}{}
\renewcommand{\HS}{\mathrm{HS}}
\providecommand{\Carrier}{}
\renewcommand{\Carrier}{\mathsf{Car}}
\providecommand{\Mask}{}
\renewcommand{\Mask}{\mathsf{Mask}}
\providecommand{\Code}{}
\renewcommand{\Code}{\mathsf{Code}}
\providecommand{\Gate}{}
\renewcommand{\Gate}{\mathsf{Gate}}
\providecommand{\Leak}{}
\renewcommand{\Leak}{\mathsf{Leak}}
\providecommand{\Debt}{}
\renewcommand{\Debt}{\mathsf{Debt}}
\providecommand{\Route}{}
\renewcommand{\Route}{\mathsf{Route}}
\providecommand{\Block}{}
\renewcommand{\Block}{\mathsf{block}}
\providecommand{\Abs}{}
\renewcommand{\Abs}{\mathsf{abs}}
\providecommand{\Res}{}
\renewcommand{\Res}{\mathsf{res}}
\providecommand{\Open}{}
\renewcommand{\Open}{\mathsf{open}}
\providecommand{\Ad}{}
\renewcommand{\Ad}{\mathsf{adm}}
\providecommand{\Image}{}
\renewcommand{\Image}{\operatorname{im}}
\providecommand{\supp}{}
\renewcommand{\supp}{\operatorname{supp}}
\providecommand{\rank}{}
\renewcommand{\rank}{\operatorname{rank}}
\providecommand{\Span}{}
\renewcommand{\Span}{\operatorname{span}}
\providecommand{\End}{}
\renewcommand{\End}{\operatorname{End}}
\providecommand{\C}{}
\renewcommand{\C}{\mathbb C}
\providecommand{\Hh}{}
\renewcommand{\Hh}{\mathbb H}
\providecommand{\Oo}{}
\renewcommand{\Oo}{\mathbb O}
\providecommand{\One}{}
\renewcommand{\One}{\mathbf 1}
\providecommand{\MC}{}
\renewcommand{\MC}{\mathcal M_{16}}
\providecommand{\PC}{}
\renewcommand{\PC}{\widehat P_{C_7}}
\providecommand{\Qp}{}
\renewcommand{\Qp}{Q_{\perp}}
\providecommand{\HDix}{}
\renewcommand{\HDix}{\mathcal H_{\Dix}}
\providecommand{\HsFour}{}
\renewcommand{\HsFour}{\mathcal H^{(4)}_{16}}
\providecommand{\HsFive}{}
\renewcommand{\HsFive}{\mathcal H^{(5)}_{32}}
\providecommand{\HsSix}{}
\renewcommand{\HsSix}{\mathcal H^{(6)}_{64}}
\providecommand{\HaloOp}{}
\renewcommand{\HaloOp}{\mathsf{Halo}}
\providecommand{\ExitOp}{}
\renewcommand{\ExitOp}{\mathsf{Exit}}
\providecommand{\Atlas}{}
\renewcommand{\Atlas}{\mathfrak A}
\providecommand{\Rop}{}
\renewcommand{\Rop}{\mathcal R_{\OP 5}^{\adm,\near K}}
\providecommand{\Lunr}{}
\renewcommand{\Lunr}{L_{\unr}}
\providecommand{\Bhcmm}{}
\renewcommand{\Bhcmm}{B_{\HC/\MM}}
\providecommand{\XD}{}
\renewcommand{\XD}{\Xi_{\Dix}}
\providecommand{\ChiD}{}
\renewcommand{\ChiD}{\chi_{\Dix}}
\providecommand{\Aset}{}
\renewcommand{\Aset}{A_{\mathrm{loc}}}
\providecommand{\Iset}{}
\renewcommand{\Iset}{\mathcal I}
\providecommand{\Mclass}{}
\renewcommand{\Mclass}{\mathfrak M}
% --- begin integrated body for r59 ---



\begin{abstract}
This delta layer executes the image audit requested after the local pass. The layer mapped captured near-kernel OP5 routes into the HS(6) Dixon carrier \(\C\otimes\Hh\otimes\Oo\), with 64 real carrier channels and a \(2^{64}\)-state Boolean activation-mask space. The present layer audits only the masks actually reachable from the capture table. It classifies the symbolic image of \(\Xi_{\Dix}^{\mathrm{loc}}\) into blocked, HC/MM absorbed, atlas routed, Dixon-admissible, and residual-open classes. Under the inherited zero leakage/debt assumptions, the declared sample image has no residual-open Dixon-mask class. This remains a conditional local OP5 closure audit inside the declared photon-frame HC-halo and HS(0)--HS(6) grammar, not a global OP5 proof outside that grammar.
\end{abstract}

\section{Why the Dixon representation helps QAM}

The QAM grammar becomes more useful when it is not only expressive but also addressable. The Dixon layer supplies this addressability. Instead of leaving a captured route as an open structural phrase, it assigns the route a finite carrier address
\[
\alpha=(a,b,c),
\qquad
 a\in\{0,1\},\quad b\in\{0,1,2,3\},\quad c\in\{0,\ldots,7\},
\]
inside the 64-channel Dixon carrier
\[
\HDix=\C\otimes\Hh\otimes\Oo.
\]
The QAM state is then not a vague continuum of possible route meanings. It is a Boolean activation mask
\[
\ChiD\in\{0,1\}^{64}.
\]
Thus QAM over Dixon gives three concrete benefits.

\begin{enumerate}[label=(\roman*)]
\item \textbf{Finite address space.} Captured routes are assigned to a finite 64-channel carrier instead of an unbounded route grammar.
\item \textbf{Auditability.} Every possible uncontrolled channel must appear as a non-admissible mask, an HS(4) leakage register, or an HS(5) cross-side debt register.
\item \textbf{Compatibility with existing algebra.} The complex, quaternionic, and octonionic factors separate phase, two-sided/spin-like transport, and Fano/octonion incidence without claiming that the Dixon carrier itself has dimension \(2^{64}\).
\end{enumerate}

The present layer uses these benefits to audit the image of the route-to-Dixon map.

\section{Inherited coding map}

Recall the conditional coding map
\[
\XD^{\mathrm{loc}}:\Rop\longrightarrow
\{0,1\}^{64}\cup\{\mathsf{blocked},\mathsf{debt},\mathsf{leak}\}.
\]
It is defined only after the local proof chain photon-frame HC-halo capture chain has reduced admissible near-kernel OP5 routes to the finite class menu
\[
\Aset=\{A_1,A_2,A_3,A_4,A_5,A_6\}.
\]
The audit does not enumerate all \(2^{64}\) masks. It enumerates the much smaller symbolic image of \(\XD^{\mathrm{loc}}\) produced by \(\Aset\).

\begin{definition}[Reachable Dixon image]
The reachable Dixon image is
\[
\Iset:=\XD^{\mathrm{loc}}(\Rop)\cap
\bigl(\{0,1\}^{64}\cup\{\Block,\Abs,\Route,\Leak,\Debt,\Res\}\bigr),
\]
where each element is represented by its class, its HS(4) leakage status, its HS(5) cross-side debt status, and its Dixon-mask class.
\end{definition}

\begin{definition}[Compressed Dixon mask class]
For a captured route class \(A_i\), define its compressed Dixon mask class by
\[
\Mclass_i
:=
\bigl(\sigma_i,\tau_i,\omega_i,\kappa_i\bigr),
\]
where \(\sigma_i\) is the selector word, \(\tau_i\) is the two-sided front/back tag, \(\omega_i\subseteq\{0,\ldots,7\}\) is the Fano/octonion support, and \(\kappa_i\) is the response status. The associated full mask is the subset of Dixon channels
\[
\supp(\ChiD(A_i))
\subseteq
\{0,1\}\times\{0,1,2,3\}\times\{0,\ldots,7\}.
\]
\end{definition}

\begin{remark}[Symbolic rather than exhaustive bit enumeration]
The audit enumerates mask families forced by the capture table. It does not list all \(2^{64}\) bitstrings. This is the intended gain: OP5 does not require an exhaustive scan of the Dixon mask space if the route image is first captured into finitely many symbolic classes.
\end{remark}

\section{Status alphabet and admissibility test}

\begin{definition}[status alphabet]
The image audit uses the status alphabet
\[
\mathfrak S_{\mathrm{loc}}
=
\{\Block,\Abs,\Route,\Ad,\Res\}.
\]
The meanings are as follows.
\begin{itemize}[leftmargin=2em]
\item \(\Block\): the route is selector-blocked and produces no active OP5 exit.
\item \(\Abs\): the route is HC/MM absorbed with explicit zero debt.
\item \(\Route\): the route is atlas-routed to another certified normal-collar piece.
\item \(\Ad\): the route is Dixon-admissible as a finite 64-channel activation mask.
\item \(\Res\): the route remains residual-open.
\end{itemize}
\end{definition}

\begin{definition}[Dixon admissibility predicate]
A reachable Dixon mask class \(\Mclass_i\) is the local pass-admissible when the following four tests hold:
\[
\Leak_{C_7}^{(4)}(A_i)=0,
\qquad
\Debt_{\mathrm{cross}}^{(5)}(A_i)=0,
\]
\[
\Lunr(A_i)=0,
\qquad
\Bhcmm(A_i)=0,
\]
and the selector/overlap word of \(A_i\) is the local pass-admissible. We write this as
\[
\mathsf{Adm}_{\Dix}^{\mathrm{loc}}(\Mclass_i)=1.
\]
\end{definition}

\begin{lemma}[Image audit reduction]
Assume the capture table is exhaustive. Then the residual-open part of the Dixon image is
\[
\Iset_{\Res}
=
\left\{\Mclass_i\;:\;
\begin{array}{l}
A_i\in\Aset,\\
\mathsf{Adm}_{\Dix}^{\mathrm{loc}}(\Mclass_i)=0,\\
A_i\text{ is neither blocked nor HC/MM absorbed}
\end{array}
\right\}.
\]
Thus the residual audit is finite at the level of captured symbolic classes.
\end{lemma}

\begin{proof}
By the local pass, every declared admissible near-kernel route belongs to one of the classes \(A_i\). By the local pass, every such class is either blocked, marked leak/debt, or assigned a Dixon mask. The predicate records exactly the leak, debt, hidden leakage, HC/MM debt, and selector/overlap admissibility tests. Hence any residual-open element must be one of the finitely many captured classes failing this predicate while not being otherwise blocked or absorbed.
\end{proof}

\section{image table}

The following table is the first Dixon image audit of the capture classes. The supports are symbolic; they specify which factor of the Dixon carrier is active, not a particle spectrum.

\begin{center}
\scriptsize
\setlength{\tabcolsep}{3pt}
\begin{tabular}{c|P{0.15\linewidth}|c|c|P{0.16\linewidth}|c|P{0.15\linewidth}}
\toprule
class & route meaning & HS(4) leak & HS(5) debt & Dixon support class & status & residual reason\\
\midrule
\(A_1\) & collar interior & 0 & 0 & neutral collar support in \(\C\otimes\Hh\otimes\Oo\) & \(\Ad\) & none\\
\(A_2\) & \(S_{12}\) switch & 0 & 0 & routed selector edge support & \(\Route\) & none\\
\(A_3\) & \(S_{23}\) switch & 0 & 0 & routed selector edge support & \(\Route\) & none\\
\(A_4\) & non-adjacent \(S_{13}\) & -- & -- & no active mask & \(\Block\) & blocked by the local pass\\
\(A_5\) & HC/MM bookkeeping face & 0 & 0 & absorbed neutral/debt-free face & \(\Abs\) & none\\
\(A_6\) & overlap edge & 0 & 0 & overlap-preserved support & \(\Route\) or \(\Ad\) & none\\
\bottomrule
\end{tabular}
\end{center}

\begin{proposition}[No residual-open class in the declared sample image]
For the declared symbolic image table, under the inherited assumptions
\[
\Leak_{C_7}^{(4)}=0,
\qquad
\Debt_{\mathrm{cross}}^{(5)}=0,
\qquad
\Lunr=0,
\qquad
\Bhcmm=0,
\]
and with the selector/overlap decisions unchanged, one has
\[
\Iset_{\Res}=\varnothing.
\]
\end{proposition}

\begin{proof}
Class \(A_4\) is blocked, hence it is not residual-open. Class \(A_5\) is HC/MM absorbed with zero debt, hence it is not residual-open. Classes \(A_2,A_3,A_6\) are routed through the selector/overlap grammar and preserve contraction margins, hence they are not residual-open. Class \(A_1\) remains inside the certified collar interior and receives an admissible Dixon support class. Since the table contains all declared classes, no residual-open symbolic image element remains.
\end{proof}

\section{Mask-image theorem}

\begin{theorem}[Dixon image audit theorem]
Assume the local proof chain photon-frame HC-halo capture chain, the HS(4)--HS(6) coding hypotheses, and the image table above. If the capture map is exhaustive over the declared near-kernel OP5 route grammar, then
\[
\Image(\XD^{\mathrm{loc}})
\subseteq
\Code_{\Dix}^{\adm}
\cup
\{\Block,\Abs,\Route\},
\]
with no residual-open Dixon image class. Equivalently, within the declared HS(0)--HS(6) grammar, the remaining OP5 audit is reduced to the independent exhaustiveness check of the capture map and the zero-register checks for HS(4), HS(5), hidden leakage, and HC/MM debt.
\end{theorem}

\begin{proof}
The capture map reduces the route grammar to \(A_1,\ldots,A_6\). The map sends each captured route to a Dixon mask or to an explicit blocked/leak/debt status. The table evaluates these statuses. Under the zero leakage/debt assumptions, no class is mapped to leak or debt. The non-adjacent selector class is blocked, the HC/MM face is absorbed debt-free, the selector and overlap edges are routed, and the collar interior is Dixon-admissible. Therefore the image contains only admissible Dixon masks and the three controlled non-residual statuses.
\end{proof}

\begin{corollary}[OP5 consequence inside the declared grammar]
Inside the declared OP5=photon-frame HC-cone grammar, the audit leaves no uncontrolled OP5 exit in the captured HS(0)--HS(6) image, provided the capture map is exhaustive and all zero-register hypotheses remain true.
\end{corollary}

\begin{remark}[What has not been proven]
This is not a global OP5 proof outside the declared grammar. The decisive remaining mathematical work is to audit the full Active-Core route menu for any admissible near-kernel OP5 branch not represented by \(A_1,\ldots,A_6\), and to preserve the zero-register conditions under integration into the monolithic source.
\end{remark}

\section{Next audit target}

The next useful layer should not add a new algebraic carrier. It should perform one of two integration checks:
\begin{enumerate}[label=(\alph*)]
\item a full Active-Core route-menu scan showing that every admissible near-kernel OP5 branch is represented by the capture classes; or
\item a monolithic integration of the local proof chain into the active source, followed by separate label, reference, layout, and route-menu residual audits.
\end{enumerate}
If both checks pass, OP5 can be stated as locally closed inside the declared photon-frame HC-halo/Dixon-coding grammar, with the remaining caveat being only the scope of that grammar.


% --- end integrated body for r59 ---

\clearpage
\section*{Integrated Delta Layer: Integrated Residual Register}
\addcontentsline{toc}{section}{Integrated Delta Layer: Integrated Residual Register}
% --- local macro refresh for r60 ---
\providecommand{\OPfive}{}
\renewcommand{\OPfive}{\mathrm{OP5}}
\providecommand{\HC}{}
\renewcommand{\HC}{\mathrm{HC}}
\providecommand{\MM}{}
\renewcommand{\MM}{\mathrm{MM}}
\providecommand{\Dix}{}
\renewcommand{\Dix}{\mathrm{Dix}}
\providecommand{\Ad}{}
\renewcommand{\Ad}{\mathsf{adm}}
\providecommand{\Block}{}
\renewcommand{\Block}{\mathsf{block}}
\providecommand{\Abs}{}
\renewcommand{\Abs}{\mathsf{abs}}
\providecommand{\Route}{}
\renewcommand{\Route}{\mathsf{route}}
\providecommand{\Res}{}
\renewcommand{\Res}{\mathsf{res}}
\providecommand{\Leak}{}
\renewcommand{\Leak}{\mathsf{leak}}
\providecommand{\Debt}{}
\renewcommand{\Debt}{\mathsf{debt}}
\providecommand{\HaloOp}{}
\renewcommand{\HaloOp}{\mathsf{Halo}}
\providecommand{\ExitOp}{}
\renewcommand{\ExitOp}{\mathsf{Exit}}
\providecommand{\NearK}{}
\renewcommand{\NearK}{\mathrm{near}\,K}
% --- begin integrated body for r60 ---



\begin{abstract}
This layer audits the available route-menu surface after the local proof chain. Its
purpose is not to introduce a new algebraic mechanism. It checks whether the
local photon-frame \(\HC\)-halo/Dixon coding chain leaves any declared
admissible near-kernel \(\OPfive\) route outside the finite capture menu
\(A_1,\ldots,A_6\). The result is deliberately scoped: the local proof chain delta
stack is internally exhausted by the local proof chain classification, while the
available monolithic basis source is only the local proof line. Therefore this layer
certifies the declared delta-stack route image and records the exact remaining
condition for a future source-native monolithic freeze: no active OP5 branch in
the full Active Core may remain outside the capture classes or the named
zero-register tests.
\end{abstract}

\section{Audit scope}

The available source surface consists of the local proof line monolithic Active Core
basis and the local proof chain delta layers. Hence performs two separate checks.

\begin{enumerate}
\item \textbf{Delta-stack route audit.} This scans the local proof chain local
photon-frame \(\HC\)-halo, tunnel, finite-atlas, capture, and Dixon-coding
chain.
\item \textbf{Basis-source compatibility scan.} This records the older the local pass
OP5 route language that predates the \(\HaloOp\)/Dixon grammar and must later be
mapped or explicitly declared out of scope during a monolithic integration.
\end{enumerate}

\begin{remark}[No false monolithic claim]
The present file is a source-level audit layer, not a source-native rebuild of a
single full the local proof line Active Core. It therefore proves a local exhaustiveness
claim for the declared the local proof chain grammar and lists the remaining monolithic audit
obligation separately.
\end{remark}

\section{Delta-stack route menu}

The local proof chain reduces the local near-kernel \(\OPfive\) problem as follows:
\[
\begin{gathered}
\text{robust photon chart}\to\text{finite atlas}\to\text{double-}\varepsilon\text{ halo},\\
\text{double-}\varepsilon\text{ halo}\to\text{rank/selector/overlap registers}\to\text{capture classes},\\
\text{capture classes}\to\text{Dixon image audit}.
\end{gathered}
\]

The effective route alphabet after is
\[
A_{\mathrm{loc}}=\{A_1,A_2,A_3,A_4,A_5,A_6\}.
\]
The Dixon audit assigns these classes the following controlled statuses.

\begin{center}
\begin{tabular}{lll}
\toprule
class & route meaning & status\\
\midrule
\(A_1\) & collar interior & \(\Ad\)\\
\(A_2\) & \(S_{12}\) selector switch & \(\Route\)\\
\(A_3\) & \(S_{23}\) selector switch & \(\Route\)\\
\(A_4\) & non-adjacent \(S_{13}\) selector & \(\Block\)\\
\(A_5\) & \(\HC/\MM\) bookkeeping face & \(\Abs\)\\
\(A_6\) & atlas-overlap edge & \(\Route\) or \(\Ad\)\\
\bottomrule
\end{tabular}
\end{center}

\begin{definition}[residual support]
Let \(\Xi_{\Dix}^{\mathrm{loc}}\) be the route-to-Dixon map of and let
\(\pi_{\mathrm{cap}}^{\mathrm{loc}}\) be the capture map. The residual support is
\[
\mathcal I_{\Res}^{\mathrm{loc}}
:=
\left\{
 r\in\mathcal R_{\OPfive}^{\Ad,\NearK}:
 \begin{array}{l}
 \pi_{\mathrm{cap}}^{\mathrm{loc}}(r)\notin \{A_1,\ldots,A_6\},\ \text{or}\\
 \Xi_{\Dix}^{\mathrm{loc}}(r)\in\{\Leak,\Debt,\Res\}
 \end{array}
\right\}.
\]
Thus \(\mathcal I_{\Res}^{\mathrm{loc}}\) records exactly the active near-kernel routes
that are either uncaptured or mapped to a non-controlled Dixon status.
\end{definition}

\begin{proposition}[Declared delta-stack residual vanishing]
Inside the declared the local proof chain photon-frame \(\HC\)-halo/Dixon grammar, the source
scan finds no declared route class beyond \(A_1,\ldots,A_6\), and maps those
classes only to \(\Ad\), \(\Block\), \(\Abs\), or \(\Route\). Hence
\[
\mathcal I_{\Res}^{\mathrm{loc}}=\varnothing
\]
for the declared delta-stack image.
\end{proposition}

\begin{proof}
The tunnel menu is reduced by the local proof chain to radial/HC boundary, selector,
rank, and overlap checks. The normal-collar certificate removes rank-drop
on certified pieces. The selector/overlap table blocks or routes the switch
and overlap classes. The capture table enumerates the six remaining local
near-kernel classes \(A_1,\ldots,A_6\). The Dixon map sends captured active
classes into the 64-channel Dixon carrier or to explicit blocked/leak/debt
statuses. The image table assigns \(A_1\) to admissible, \(A_2,A_3,A_6\)
to routed/admissible, \(A_4\) to blocked, and \(A_5\) to absorbed. No table row
is assigned \(\Leak\), \(\Debt\), or \(\Res\). Therefore the declared local image
has empty residual support.
\end{proof}

\section{Source scan facts}

The source scan of the local proof chain delta files produced the following structural
metadata.

\begin{center}\small
\begin{tabular}{lrrrr}
\toprule
layer & lines & labels & raw lower-case halo & raw lower-case exit\\
\midrule
the local pass & 239 & 0 & 0 & 0\\
the local pass & 245 & 0 & 0 & 0\\
the local pass & 293 & 0 & 0 & 0\\
the local pass & 317 & 0 & 1 & 1\\
the local pass & 276 & 0 & 0 & 0\\
the local pass & 257 & 0 & 0 & 0\\
the local pass & 284 & 0 & 0 & 0\\
the local pass & 296 & 0 & 0 & 0\\
the local pass & 387 & 0 & 0 & 0\\
the local pass & 277 & 0 & 0 & 0\\
\bottomrule
\end{tabular}
\end{center}

The single lower-case occurrences are compatibility aliases in the source;
they are not active rendered notation in the local proof chain line. No duplicate labels
were found across the local proof chain delta files.

The available monolithic basis contains the older selector-gated OP5 route
language and does not yet contain the later \(\HaloOp\), near-kernel, or Dixon
route grammar as an integrated source-native block. The scan counts were:

\begin{center}\small
\begin{tabular}{lr}
\toprule
anchor & count\\
\midrule
\texttt{OP\textasciitilde5} & 567\\
\texttt{OP5} & 133\\
selector-gated reachability & 13\\
TE budgets & 4\\
HC/MM residual & 2\\
post-promotion residual & 7\\
OP\textasciitilde4 & 55\\
Dixon & 8\\
Halo & 0\\
near-kernel & 0\\
\bottomrule
\end{tabular}
\end{center}

This confirms that the local proof chain are a new local closure grammar over the basis,
not yet a source-native monolithic replacement of every older OP5 paragraph.

\section{Compatibility with the OP5 route surface}

The selector-gated OP5(B) statement leaves four route obligations after a
first reachability pass: TE budgets, \(\HC/\MM\) residual compatibility,
post-promotion residual control, and OP4 stability. The local proof chain directly
attacks the near-kernel \(\HC\)-cone route obstruction, not the whole OP5(B)
boundary-promotion package.

\begin{definition}[Monolithic residual register]
For a future full Active-Core integration define
\[
\mathcal R_{\mathrm{mono}}^{\mathrm{loc}}
:=
\mathcal R_{\mathrm{uncap}}
\cup
\mathcal R_{\mathrm{HS4\ leak}}
\cup
\mathcal R_{\mathrm{HS5\ debt}}
\cup
\mathcal R_{L_{\mathrm{unr}}}
\cup
\mathcal R_{B_{\HC/\MM}}
\cup
\mathcal R_{\mathrm{loc\text{-}res}}.
\]
Here \(\mathcal R_{\mathrm{uncap}}\) denotes active near-kernel routes not mapped
by \(\pi_{\mathrm{cap}}^{\mathrm{loc}}\), while \(\mathcal R_{\mathrm{loc\text{-}res}}\)
denotes any older OP5(B) budget, post-promotion, or OP4-stability condition
not identified with the local \(\HC\)-cone grammar.
\end{definition}

\begin{theorem}[route-menu exhaustiveness theorem, scoped form]
Assume the working identification
\[
\OPfive_{\NearK}=\text{photon-frame }\HC\text{-cone source-return reflector}
\]
and assume that the capture map is exhaustive over the declared near-kernel
route grammar. If
\[
\mathcal R_{\mathrm{HS4\ leak}}
=\mathcal R_{\mathrm{HS5\ debt}}
=\mathcal R_{L_{\mathrm{unr}}}
=\mathcal R_{B_{\HC/\MM}}
=\varnothing
\]
and the selector/overlap decisions are preserved under integration, then
there is no uncontrolled near-kernel OP5 exit in the local proof-chain delta stack route
menu.

If, in addition, a future monolithic scan proves
\[
\mathcal R_{\mathrm{uncap}}=\varnothing
\qquad\text{and}\qquad
\mathcal R_{\mathrm{loc\text{-}res}}=\varnothing,
\]
then the same statement may be promoted from a local delta-stack closure theorem
to a monolithic local OP5 closure theorem inside the declared
\(\HC\)-halo/Dixon grammar.
\end{theorem}

\begin{proof}
The first part is exactly Proposition 1 combined with the zero-register
hypotheses: every captured near-kernel route is admissible, blocked, absorbed,
or routed, and no leak/debt/residual status remains. The second part adds the
only missing monolithic conditions. If no older active branch lies outside the
capture map and no residual obligation remains unmapped, the local
exhaustiveness statement covers the integrated route menu for the
declared grammar.
\end{proof}

\begin{corollary}[Current OP5 status after]
The local proof chain local route image is exhausted. The current remaining OP5 burden
is no longer a free route explosion; it is the monolithic integration check
\[
\mathcal R_{\mathrm{mono}}^{\mathrm{loc}}=\varnothing.
\]
Thus OP5 is locally closed inside the declared delta-stack grammar and remains
open only at the full Active-Core integration/exhaustiveness boundary.
\end{corollary}

\section{Next integration rule}

The next source-native step should not add another algebraic layer. It should
merge or append the local proof chain into the active source and run a residual search with
the following pass/fail test:
\[
\operatorname{supp}\bigl(\mathcal R_{\OPfive}^{\Ad,\NearK}\bigr)
\subseteq
\bigcup_{i=1}^{6}(\pi_{\mathrm{cap}}^{\mathrm{loc}})^{-1}(A_i)
\cup
\{\Block,\Abs,\Route,\Ad\}.
\]
Any hit outside this set must be named explicitly as one of: uncaptured route,
HS4 leakage, HS5 cross-side debt, hidden leakage, HC/MM debt, post-promotion
residual, TE-budget residual, or OP4-stability residual.

\section{Conservativity guards}

\begin{itemize}
\item does not claim global OP5 closure outside the declared
photon-frame \(\HC\)-halo/Dixon grammar.
\item does not claim OP1, OP4, QCD/confinement, CKM/PMNS, quark masses, or a
Standard Model spectrum derivation.
\item The Dixon carrier is 64-dimensional; \(2^{64}\) is the Boolean activation
mask count, not the Dixon dimension.
\item \(\HC/\MM\) absorption remains explicit bookkeeping with zero unpaid debt,
not residual deletion.
\item A source-native monolithic rebuild remains required before making a public
release-style OP5 closure statement.
\end{itemize}


% --- end integrated body for r60 ---
\clearpage
\section*{Monolithic OP5 Promotion Audit}
\addcontentsline{toc}{section}{the local proof line Monolithic OP5 Promotion Audit}

\begin{definition}[monolithic OP5 residual register]
Let \(\mathcal R_{\mathrm{mono}}^{\mathrm{loc}}\) denote the set of active OP5-facing near-kernel branches in the integrated Active-Core source which are not assigned to one of the capture classes
\[
A_1,\ldots,A_6,
\]
and are not discharged by the named zero-register tests
\[
L_{\mathrm{unr}}=0,\qquad B_{\mathrm{HC/MM}}=0,
\qquad R_{\mathrm{HS4}}=0,
\qquad D_{\mathrm{HS5}}=0.
\]
A hit is counted as residual only if it is active in the current local proof line OP5-facing grammar, not merely preserved historical witness language.
\end{definition}

\begin{proposition}[integrated delta-stack exhaustion]
In the integrated local proof chain delta stack, the local photon-frame HC-halo/Dixon route image is exhausted by the six classes
\[
A_1,\ldots,A_6.
\]
Their status assignment remains
\[
A_1\mapsto \mathsf{adm},\quad
A_2\mapsto \mathsf{route},\quad
A_3\mapsto \mathsf{route},\quad
A_4\mapsto \mathsf{block},\quad
A_5\mapsto \mathsf{abs},\quad
A_6\mapsto \mathsf{route}\ \text{or}\ \mathsf{adm}.
\]
Consequently the residual Dixon image inside the declared local grammar is empty:
\[
\mathcal I_{\mathrm{res}}^{\mathrm{loc}}=\varnothing.
\]
\end{proposition}

\begin{proof}
This is the source-native integration of the audit conclusion with the local proof chain delta bodies included in the same Active-Core source. The local proof-chain layers provide the robust local chart and finite-atlas candidate, the local proof chain reduce the OP5 route question to halo/tunnel, rank, selector, and overlap registers, supplies the finite capture menu, maps captured routes to the 64-dimensional Dixon carrier with a \(2^{64}\)-state mask space, audits the reachable Dixon image, and verifies that no residual-open class remains in that declared delta image.
\end{proof}

\begin{theorem}[local OP5 closure in the declared HC-halo/Dixon grammar]
Assume the working identification that the OP5-facing near-kernel branch is the photon-frame HC-cone/reflector route, and assume the inherited zero-register conditions
\[
L_{\mathrm{unr}}=0,
\qquad
B_{\mathrm{HC/MM}}=0,
\qquad
R_{\mathrm{HS4}}=0,
\qquad
D_{\mathrm{HS5}}=0.
\]
Then the integrated local proof chain Active-Core delta stack leaves no uncontrolled local OP5 exit inside the declared photon-frame HC-halo/Dixon grammar.
\end{theorem}

\begin{proof}
By the capture layer, every declared admissible near-kernel route in the grammar is assigned to \(A_1,\ldots,A_6\), provided the capture map is exhaustive. By the selector/overlap layer, the active switch and overlap channels are blocked, routed, absorbed, or contraction-preserving. By the Dixon-image layers, every captured route has either a Dixon-admissible 64-channel mask or is classified as blocked, routed, or HC/MM-absorbed. The assumed zero-register conditions remove hidden leakage, unpaid HC/MM absorption, HS(4) meta-septum leakage, and HS(5) cross-side debt. Hence no residual-open OP5 class remains in the declared local grammar.
\end{proof}

\begin{remark}[Scope of the promotion]
This statement promotes the local proof-chain delta stack result into the available Active-Core source as a local OP5 closure theorem inside the declared HC-halo/Dixon grammar. It does not claim global OP5 closure outside that grammar, does not close OP1 or OP4, and does not derive QCD/confinement, CKM/PMNS values, quark masses, or a Standard Model spectrum.
\end{remark}

\begin{remark}[Inherited historical OP5 language]
The older Active-Core basis contains preserved historical OP1--OP5 and OP5(B) statements from earlier release lines. In those blocks remain theorem-memory and genealogy unless explicitly referenced by the current HC-halo/Dixon grammar. The governing current promotion claim is the local theorem above.
\end{remark}


\clearpage
\section*{OP5 Closure Hardening and Old-Branch Mapping Audit}
\addcontentsline{toc}{section}{the OP5 Closure Hardening and Old-Branch Mapping Audit}
\noindent\textbf{Target-version rule.} This audit belongs to the provenance layer. In the reader-facing proof spine it is interpreted semantically, not as a mathematical version marker. It hardens the local OP5 promotion and maps inherited OP5 language into current, historical, guard, or outside-scope status.

\begin{definition}[governing OP5 status]
The governing OP5 status is that OP5 is locally closed inside the declared photon-frame HC-halo/Dixon grammar under \[L_{\mathrm{unr}}=0,\quad B_{\mathrm{HC}/\mathrm{MM}}=0,\quad R_{\mathrm{HS4}}=0,\quad D_{\mathrm{HS5}}=0.\] This is not a global OP5 theorem outside that grammar and does not close OP1 or OP4.
\end{definition}

\begin{proposition}[no-active-old-branch promotion rule]
An inherited OP5-facing paragraph may be read as part of the current OP5 closure claim only if it is explicitly mapped into the HC-halo/Dixon grammar or is one of the named zero-register guards. Otherwise it is historical theorem-memory, a guard, or outside the scope of the local theorem.
\end{proposition}

\begin{theorem}[compact OP5 closure spine]
Inside the declared photon-frame HC-halo/Dixon grammar, OP5 has the finite proof spine
\[
\begin{aligned}
&\text{HC cone}
\Rightarrow
\text{double-}\varepsilon\text{ halo}
\Rightarrow
\text{tunnel reduction}
\Rightarrow
\text{normal collar}\\[0.35em]
&\Rightarrow
\text{selector/overlap routing}
\Rightarrow
\text{Dixon image audit}
\Rightarrow
\text{no residual local OP5 exit}.
\end{aligned}
\]
Under the four zero-register assumptions, the local residual register of this spine is empty.
\end{theorem}

\begin{remark}[No version promotion]
The present document is a v4.3.0 release-final. A later public line should be reserved for broader OP4/OP1 promotion or for a genuinely new architecture beyond the OP5 hardening chain.
\end{remark}


% --- Finite Perturbation Threshold Atlas Insert ---

\section{Finite Perturbation Threshold Atlas for the Triolic Orbit Scan}
\label{sec:finite-perturbation-threshold-atlas}

This section records a compact finite threshold atlas for the alternating
triolic orbit-scan model.  The purpose is not to promote the toy chart to a
physical trajectory simulation, but to make the robustness statement falsifiable
at the finite readout level.

Let the finite scan diagnostics be
\[
    \widetilde E_{N,M},\qquad
    \widetilde R_{N,M}^{\min},\qquad
    C_{95}^{\mathrm{even}},\qquad
    C_{05}^{\mathrm{odd}},
\]
where \(\widetilde E_{N,M}\) is the bifocal defect,
\(\widetilde R_{N,M}^{\min}\) is the smallest guard margin,
\(C_{95}^{\mathrm{even}}\) is the 95-percentile even-sector core-sightline
residual, and \(C_{05}^{\mathrm{odd}}\) is the 5-percentile odd-sector
complementarity residual.

\begin{definition}[Finite threshold predicate]
A perturbed finite scan is threshold-admissible when
\[
    \widetilde E_{N,M}\leq \tau_E,\qquad
    \widetilde R_{N,M}^{\min}>0,\qquad
    C_{95}^{\mathrm{even}}\leq \tau_C,\qquad
    C_{05}^{\mathrm{odd}}\geq \tau_O.
\]
The predicate is a finite readout certificate only. It does not assert a
physical equation of motion.
\end{definition}

\begin{lemma}[Index-level parity is threshold-invariant]
For the shifted Kronecker transport
\[
    \Sigma_L{}^i{}_j=\delta^i_{j+1\ \mathrm{mod}\ L},
\]
the even--odd identities
\[
    \Pi_0\Sigma_L=\Sigma_L\Pi_1,\qquad
    \Pi_1\Sigma_L=\Sigma_L\Pi_0,\qquad
    \Pi_r\Sigma_L^2=\Sigma_L^2\Pi_r
\]
are independent of the geometric perturbation amplitudes.
\end{lemma}

\begin{proposition}[Finite threshold-atlas reading]
The finite threshold atlas separates three layers:
\[
\text{index parity},\qquad
\text{protected-core guard},\qquad
\text{geometric triolic balance}.
\]
The first is exact under the shifted Kronecker transport.  The second is
certified by positive guard margin.  The third is a genuine geometric residual
and can fail under core-offset or antipodal-breaking perturbations.
\end{proposition}

\begin{remark}[Compact numerical atlas instance]
In the compact finite atlas run, the one-parameter scans estimate the following
last-stable / first-unstable transition windows:
\[
\begin{array}{c|c|c}
\text{parameter} & \text{last stable} & \text{first unstable}\\
\hline
\text{core offset} & 0.030 & 0.045\\
\text{antipodal-breaking jitter} & 0.080 & \text{none}\\
\text{lens asymmetry} & 0.315 & 0.360\\
\text{symmetric radial ripple} & 0.030 & 0.045\\
\text{twist amplitude} & 0.080 & \text{none}
\end{array}
\]
The atlas is not a physical proof.  It is a reproducible diagnostic showing
which perturbations preserve the local orbit-scan readout and which ones create
a named obstruction.
\end{remark}

\begin{theorem}[Finite threshold-atlas guard]
A finite orbit-scan atlas may be used as a local readout certificate only when
the threshold predicate holds.  If the guard remains positive but the
even-sector core-sightline residual exceeds threshold, then the failure is not a
through-core event; it is a triolic-balance obstruction.
\end{theorem}


% --- Finite Core-Recentering Calibration Insert ---

\section{Finite Core-Recentering Calibration for the Triolic Orbit Scan}
\label{sec:finite-core-recentering-calibration}

The threshold atlas identifies core offset as the most sensitive finite
obstruction for the triolic orbit scan.  This section records the corresponding
calibration layer.  The purpose is conservative: a displaced declared core is
not interpreted as a permitted through-core motion.  It is treated as a
calibration defect unless the residual remains after recentering.

Let \(\mathcal I_0\) denote the even-sector sample indices and let
\((p_n,q_n)\), \(n\in\mathcal I_0\), be the paired even-sector readout points.
The finite midpoint estimator is
\[
    \widehat K
    :=
    \frac1{|\mathcal I_0|}
    \sum_{n\in\mathcal I_0} \frac{p_n+q_n}{2} .
\]
If \(K_0\) is the declared collapse centre, the recentering displacement is
\[
    \Delta K := \widehat K-K_0 .
\]

\begin{definition}[Recentered triolic residuals]
\label{def:recentered-triolic-residuals}
For each even pair define
\[
    C_n(K)
    :=
    \frac{\|(p_n-K)\times(q_n-K)\|}
         {\|p_n-K\|\,\|q_n-K\|},
    \qquad
    M_n(K):=\left\|K-\frac{p_n+q_n}{2}\right\| .
\]
The pre-calibration residuals are \(C_n(K_0),M_n(K_0)\).  The recentered
residuals are \(C_n(\widehat K),M_n(\widehat K)\).
\end{definition}

\begin{lemma}[Midpoint recovery in the antipodal case]
\label{lem:midpoint-recovery-antipodal}
Assume that all even-sector pairs are antipodal about a common point \(K_*\),
that is
\[
    \frac{p_n+q_n}{2}=K_*
    \qquad(n\in\mathcal I_0).
\]
Then \(\widehat K=K_*\).  In particular, if the declared centre \(K_0\) is
offset from the actual finite balance centre, the estimator recovers the
balance centre from the paired readout data.
\end{lemma}

\begin{proof}
The estimator is the average of the pair midpoints.  If every midpoint equals
\(K_*\), their average equals \(K_*\).
\end{proof}

\begin{proposition}[Finite recentering criterion]
\label{prop:finite-recentering-criterion}
A finite core-offset defect is called calibratable when
\[
    C_{95}(\widehat K)\leq \tau_C,
    \qquad
    R^{\min}(\widehat K)>0,
\]
and the displacement \(\Delta K\) is recorded as a calibration displacement,
not as a through-core route.  If these inequalities fail, the residual is a
persistent triolic-balance obstruction.
\end{proposition}

\begin{remark}[Numerical calibration instance]
\label{rem:numerical-core-recentering-instance}
In the finite calibration probe, the uncalibrated core-offset threshold was
approximately
\[
    \text{last stable}=0.030,
    \qquad
    \text{first unstable}=0.040 .
\]
After midpoint recentering, the scan remained stable through the tested range,
with
\[
    \max C_{95}(\widehat K)
    \approx 6.639e-16,
    \qquad
    \min R^{\min}(\widehat K)
    \approx 1.180>0 .
\]
Thus the tested core-offset failure is a calibratable centre-declaration defect
in this finite readout chart, not a permitted crossing of the protected core.
\end{remark}

\begin{theorem}[Core-recentering guard]
\label{thm:core-recentering-guard}
Finite recentering can repair a triolic scan only as a readout calibration if
three conditions hold: the recentered even-sector residual is below threshold,
the protected-core guard remains positive, and the displacement \(\Delta K\) is
reported as calibration provenance.  Under these conditions the repaired scan
may be used as a local readout certificate, but it does not promote the orbit to
a physical through-core trajectory.
\end{theorem}

\begin{proof}
The residual condition restores the finite even-sector core-sightline test.
The positive guard condition keeps all sampled points outside the protected
core.  Recording \(\Delta K\) prevents the calibration from being mistaken for a
new dynamical claim.  Hence the procedure repairs the declared centre of the
readout chart only; it does not assert a through-core motion.
\end{proof}


% --- Crystal--Fractal Readout Probe Insert ---

\section{Crystal--Fractal Readout Probe from Calibrated Orbit Scans}
\label{sec:crystal-fractal-readout-probe}

The calibrated triolic orbit scan separates a centre-calibration problem from
a residual-geometry problem.  This makes the crystal and fractal readings
operational: the crystal reading concerns coherence across directions, while
the fractal reading concerns recurrence across scales.

\begin{definition}[Ringwise calibrated balance centres]
For each scan ring \(\ell\), define the even-triole balance centre
\[
    \widehat K_\ell
    :=
    \frac{1}{|\mathcal I_0|}
    \sum_{n\in\mathcal I_0}
    \frac{p_{\ell n}+q_{\ell n}}{2}.
\]
The calibrated global centre is
\[
    \widehat K
    :=
    \frac{1}{L}\sum_{\ell=1}^{L}\widehat K_\ell.
\]
The finite crystal-coherence residual is
\[
    \mathcal C_{\mathrm{crys}}
    :=
    \left(
    \frac{1}{L}
    \sum_{\ell=1}^{L}
    \|\widehat K_\ell-\widehat K\|^2
    \right)^{1/2}.
\]
\end{definition}

\begin{definition}[Residual readout set across scales]
Let \(s\in S\) range over a finite set of scan scales.  For each scale, let
\(\Sigma_K^{\Gamma,s}\) be the finite periapsis envelope and let
\[
    \mathcal R_s
    :=
    \{x_{\ell m}^{s}:\mathcal E_F(x_{\ell m}^{s})>\tau_{\mathrm{res}}\}
\]
be the residual readout set above the declared bifocal-defect threshold.  The
multiscale residual readout is
\[
    \mathcal R_{\mathrm{multi}}
    :=
    \bigcup_{s\in S}\mathcal R_s.
\]
\end{definition}

\begin{definition}[Finite residual box-counting diagnostic]
For box length \(\delta>0\), let \(N_{\mathcal R}(\delta)\) be the number of
occupied boxes meeting \(\mathcal R_{\mathrm{multi}}\).  A finite scaling
diagnostic is obtained by fitting
\[
    \log N_{\mathcal R}(\delta)
    \approx
    D_{\mathrm{res}}\log(1/\delta)+b.
\]
The number \(D_{\mathrm{res}}\) is a finite residual-scaling estimate, not a
proof of Hausdorff dimension.
\end{definition}

\begin{proposition}[Crystal--fractal separation]
The calibrated orbit-scan data split into two independent readouts:
\[
    \mathcal C_{\mathrm{crys}}
    \quad\text{and}\quad
    D_{\mathrm{res}}.
\]
Small \(\mathcal C_{\mathrm{crys}}\) certifies coherent directional
centering, i.e. a crystal-like local cell centre.  Nontrivial, scale-persistent
\(\mathcal R_{\mathrm{multi}}\) certifies residual scaling structure, i.e. a
fractal-like defect readout.  Neither statement asserts a physical trajectory
through the protected core.
\end{proposition}

\begin{remark}[Finite diagnostic instance]
In the finite toy diagnostic, the maximum ring-centre RMS spread was
\[
    \mathcal C_{\mathrm{crys}}^{\max}\approx 8.699e-17,
\]
with 30216 residual points above threshold across the multiscale
scan.  The box-counting fit returned
\[
    D_{\mathrm{res}}\approx 2.435.
\]
This is a diagnostic number: it indicates that residual structure can be
measured across scales, not that the protected-core surface has already been
identified with a mathematical fractal.
\end{remark}

\begin{theorem}[Readout guard for crystal and fractal claims]
A crystal claim is admissible only as a statement about coherence of the
calibrated centres \(\widehat K_\ell\).  A fractal claim is admissible only as a
statement about the scaling of residual readout sets across declared scan
scales.  The orbit-scan data do not by themselves promote either claim to a
global physical closure theorem.
\end{theorem}


% --- Residual Classification Atlas Insert ---

\section{Residual Classification Atlas for Crystal--Fractal Readouts}
\label{sec:residual-classification-atlas}

The crystal--fractal readout probe shows that residual sets can be measured
across directions and scales.  The present atlas classifies the residuals into
finite diagnostic classes.  This is not a physical closure theorem; it is a
readout classification layer.

\begin{definition}[Residual classification alphabet]
Let \(\mathcal R_{\mathrm{multi}}\) be the multiscale residual readout set.
The finite classification alphabet is
\[
    \mathfrak A_{\mathcal R}
    =
    \{\mathrm{crys},\mathrm{frac},\mathrm{lens},
      \mathrm{guard},\mathrm{noise}\}.
\]
Here \(\mathrm{crys}\) denotes directional or angular recurrence,
\(\mathrm{frac}\) denotes recurrence across scan scales,
\(\mathrm{lens}\) denotes high bifocal stress,
\(\mathrm{guard}\) denotes proximity to the protected-core guard, and
\(\mathrm{noise}\) denotes the unclassified remainder under the declared
finite rules.
\end{definition}

\begin{definition}[Class tags and primary class]
A residual point \(x\in\mathcal R_{\mathrm{multi}}\) may carry several tags:
\[
    \tau_{\mathrm{crys}}(x),\quad
    \tau_{\mathrm{frac}}(x),\quad
    \tau_{\mathrm{lens}}(x),\quad
    \tau_{\mathrm{guard}}(x).
\]
For disjoint counting, the primary class is assigned by the conservative order
\[
    \mathrm{guard}
    \succ
    \mathrm{lens}
    \succ
    \mathrm{frac}
    \succ
    \mathrm{crys}
    \succ
    \mathrm{noise}.
\]
Thus a point close to the guard is counted as a guard residual even if it also
has a scale-recurrence tag.
\end{definition}

\begin{proposition}[Residual classification separation]
The classification atlas separates four causes of residual structure:
\[
\text{directional recurrence},\quad
\text{scale recurrence},\quad
\text{bifocal stress},\quad
\text{guard proximity}.
\]
Only the first two support crystal--fractal interpretation.  Lens and guard
classes are obstruction classes and must not be promoted to crystal or fractal
evidence without an additional argument.
\end{proposition}

\begin{remark}[Finite classification instance]
In the finite diagnostic instance,
\[
    |\mathcal R_{\mathrm{multi}}|=30216.
\]
The primary class counts are
\[
\begin{array}{c|r}
\text{class} & \text{count}\\
\hline
\mathrm{crys} & 0\\
\mathrm{frac} & 27194\\
\mathrm{lens} & 3022\\
\mathrm{guard} & 0\\
\mathrm{noise} & 0
\end{array}
\]
The finite box-counting estimates by primary class are
\[
\begin{array}{c|c}
\text{class} & D_{\mathrm{box}}\\
\hline
\mathrm{crys} & \text{n/a}\\
\mathrm{frac} & 2.379\\
\mathrm{lens} & 1.874\\
\mathrm{guard} & \text{n/a}\\
\mathrm{noise} & \text{n/a}
\end{array}
\]
These numbers are finite readout diagnostics only.  They do not identify the
protected-core surface with a mathematical fractal.
\end{remark}

\begin{theorem}[Classification guard]
A residual point may be used as crystal evidence only through its directional
recurrence tag and as fractal evidence only through its scale-recurrence tag.
If the point is primarily classified as lens or guard, then it remains an
obstruction point until a separate routing or calibration step reclassifies it.
\end{theorem}

\begin{proof}
The tags are defined by distinct finite predicates: angular recurrence,
scale support, bifocal stress, and guard proximity.  The primary order assigns
obstruction classes before interpretive classes.  Hence a lens-stressed or
guard-near residual cannot be counted as crystal--fractal evidence merely
because it appears in the residual set.
\end{proof}


% --- Finite Hopf-Linking Readout Test Insert ---

\section{Finite Hopf-Linking Readout Test}
\label{sec:finite-hopf-linking-readout-test}

The orbit-scan, calibration, crystal--fractal, and residual-classification
layers describe metric and residual structure. The next readout layer asks a
topological question: whether the scan rings behave like a finite family of
Hopf-linked fibres.

\begin{definition}[Finite Hopf-fibre chart]
Let \(S^3\subset\mathbb C^2\) and let a fibre over a base point
\((\alpha,\beta)\in S^2\) be parametrised by
\[
    z_1(t)=e^{i(t+\beta)}\cos(\alpha/2),
    \qquad
    z_2(t)=e^{it}\sin(\alpha/2),
    \qquad
    t\in[0,2\pi).
\]
Writing \(z_1=X_1+iX_2\) and \(z_2=X_3+iX_4\), the finite chart uses the
stereographic projection
\[
    (X_1,X_2,X_3,X_4)
    \longmapsto
    \frac{(X_1,X_2,X_3)}{1-X_4}
    \in\mathbb R^3.
\]
A finite Hopf readout family is a collection
\[
    \Gamma_H=\{\gamma_1,\ldots,\gamma_N\}
\]
of such projected fibres sampled by a finite polygonal mesh.
\end{definition}

\begin{definition}[Finite Gauss-linking estimator]
For two disjoint closed polygonal curves \(\gamma_i,\gamma_j\), define
\[
    \widehat{\operatorname{Lk}}(\gamma_i,\gamma_j)
    :=
    \frac{1}{4\pi}
    \sum_{a,b}
    \frac{
    (r_a-r_b)\cdot(\Delta r_a\times \Delta r_b)
    }{\|r_a-r_b\|^3},
\]
where \(r_a,r_b\) are segment midpoints and \(\Delta r_a,\Delta r_b\) are the
corresponding segment vectors. This is a finite quadrature readout of the
Gauss linking integral.
\end{definition}

\begin{definition}[Linking stability matrix]
The finite linking matrix is
\[
    L_{ij}:=\widehat{\operatorname{Lk}}(\gamma_i,\gamma_j),
    \qquad
    L_{ii}:=0.
\]
For a declared tolerance \(\tau_L>0\), a pair is called Hopf-stable when
\[
    \bigl||L_{ij}|-1\bigr|\leq \tau_L.
\]
Pairs that fail this condition are finite topological readout obstructions.
\end{definition}

\begin{proposition}[Topological readout separation]
The Hopf-linking test is independent of the residual classification alphabet.
A residual point may be scale-recurring or lens-stressed, while the surrounding
fibre pair remains topologically Hopf-stable. Conversely, a linking failure is
a topological obstruction even if the local metric residual is small.
\end{proposition}

\begin{remark}[Finite diagnostic instance]
In the finite diagnostic instance, \(N=10\) projected Hopf fibres were sampled.
There are 45 unordered fibre pairs. With tolerance
\(\tau_L=0.180\), the run returned
\[
    \#\mathrm{stable\ pairs}=45,
    \qquad
    \#\mathrm{unstable\ pairs}=UN45.
\]
The mean and maximum deviations from unit linking were
\[
    \overline{\Delta}_L\approx 0.0001,
    \qquad
    \Delta_L^{\max}\approx 0.0001.
\]
These numbers are finite readout diagnostics. They do not prove that the
physical orbit-scan space is globally a Hopf fibration.
\end{remark}

\begin{theorem}[Hopf-linking guard]
A finite Hopf-linking readout may support a local topological-organisation
claim only when the sampled fibre pairs are disjoint and the linking stability
matrix is unit-stable within the declared tolerance. If a pair fails this
test, the failure is a topological readout obstruction, not a metric residual
and not a through-core trajectory.
\end{theorem}

\begin{proof}
The estimator is computed from closed curve geometry and compares pairwise
linking. Its input is not the residual magnitude at a point but the relative
winding of two closed fibres. Therefore a failure of unit linking identifies a
topological obstruction in the finite readout layer. Since the curves are
sampled outside the protected core, such a failure is not a permission to cross
the core.
\end{proof}


% --- Hopf-Stable Residual Carrier Routing Insert ---

\section{Hopf-Stable Residual Carrier Routing}
\label{sec:hopf-stable-residual-carrier-routing}

The finite Hopf-linking readout verifies that the sampled fibre family is
topologically stable in the local chart.  The residual classification atlas
identifies metric residual classes.  The present layer couples these two
readouts by routing residual points to their nearest Hopf-stable carrier
fibres.

\begin{definition}[Normalised residual-carrier assignment]
Let \(\mathcal R_{\mathrm{multi}}\) be the multiscale residual set and let
\[
    \Gamma_H=\{\gamma_1,\ldots,\gamma_N\}
\]
be the finite Hopf-fibre family.  After applying a common centred
normalisation to the residual cloud and the Hopf-fibre cloud, define
\[
    c(x)
    :=
    \operatorname*{arg\,min}_{1\leq j\leq N}
    \operatorname{dist}(x,\gamma_j).
\]
The carrier distance is
\[
    d_H(x):=\operatorname{dist}(x,\gamma_{c(x)}).
\]
For a declared threshold \(\tau_H\), a residual is called carrier-resolved when
\[
    d_H(x)\leq \tau_H.
\]
\end{definition}

\begin{definition}[Hopf-stable residual carrier]
A carrier-resolved residual \(x\) is called Hopf-stably carried when its
assigned fibre belongs to a unit-stable linking matrix:
\[
    \bigl||L_{ij}|-1\bigr|\leq\tau_L
    \quad
    \text{for all tested }i\neq j.
\]
Residuals that fail the carrier threshold are off-carrier residuals, not
through-core trajectories.
\end{definition}

\begin{proposition}[Metric residual versus topological carrier]
Residual class and carrier stability are logically distinct.  A point can be
fractal-recurring or lens-stressed as a metric residual while still being
carried by a Hopf-stable fibre family.  Conversely, a non-unit linking matrix
would create a topological obstruction even if the metric residual class were
small.
\end{proposition}

\begin{remark}[Finite carrier-routing instance]
In the finite diagnostic instance,
\[
    |\mathcal R_{\mathrm{multi}}|=30216,
    \qquad
    N=10.
\]
The Hopf-linking matrix had
\[
    45
\]
stable pairs out of
\[
    45
\]
tested pairs.  With carrier threshold
\[
    \tau_H=0.6835,
\]
the carrier routing returned
\[
    \#\mathrm{carrier\ resolved}=28705,
    \qquad
    \#\mathrm{off\ carrier}=1511.
\]
The carrier-resolved fraction was
\[
    \rho_H\approx 0.950.
\]
The normalised carrier-load entropy was
\[
    H_H\approx 0.986.
\]
\end{remark}

\begin{theorem}[Hopf-stable residual carrier guard]
If the Hopf-linking matrix is unit-stable and a residual is carrier-resolved,
then the residual may be read as metric or scale structure carried by stable
local Hopf topology.  It may not be promoted to global physical closure.  If a
residual is off-carrier, the correct conclusion is a carrier-routing
obstruction, not a collapse-core crossing.
\end{theorem}

\begin{proof}
The Hopf-linking matrix tests pairwise topological organisation of the carrier
fibres.  The residual classification tests metric, scale, lens, and guard
properties of points.  The carrier assignment relates these layers by nearest
finite readout distance.  Therefore a carrier-resolved residual inherits only
the finite local carrier stability, while an off-carrier residual records the
failure of this finite routing predicate.
\end{proof}


% --- Off-Carrier Residual Obstruction Split Insert ---

\section{Off-Carrier Residual Obstruction Split}
\label{sec:off-carrier-residual-obstruction-split}

The Hopf-stable carrier routing leaves a finite off-carrier residual set.  This
section classifies that remainder.  The purpose is to prevent an off-carrier
point from being misread as a through-core trajectory.  It is a routing
obstruction inside the finite readout chart.

\begin{definition}[Off-carrier residual set]
Let
\[
    \mathcal R_{\mathrm{off}}
    :=
    \{x\in\mathcal R_{\mathrm{multi}}: d_H(x)>\tau_H\}
\]
be the residuals that fail the declared Hopf-carrier distance threshold.  These
points are not discarded.  They are routed into an obstruction alphabet.
\end{definition}

\begin{definition}[Off-carrier obstruction alphabet]
The finite obstruction alphabet is
\[
    \mathfrak A_{\mathrm{off}}
    =
    \{\mathrm{lens\text{-}extreme},
      \mathrm{scale\text{-}boundary},
      \mathrm{projection\text{-}chart},
      \mathrm{missing\text{-}fiber},
      \mathrm{persistent}\}.
\]
The meanings are:
\[
\begin{array}{c|l}
\text{class} & \text{readout meaning}\\
\hline
\mathrm{lens\text{-}extreme} & \text{high bifocal defect or lens stress}\\
\mathrm{scale\text{-}boundary} & \text{finite scale-atlas boundary effect}\\
\mathrm{projection\text{-}chart} & \text{large chart distance / possible projection artifact}\\
\mathrm{missing\text{-}fiber} & \text{near-threshold point requiring denser carrier family}\\
\mathrm{persistent} & \text{unexplained carrier obstruction}
\end{array}
\]
\end{definition}

\begin{proposition}[Off-carrier split guard]
Every off-carrier residual must be classified before it can be used as
evidence for a topological defect.  Lens-extreme, scale-boundary,
projection-chart, and missing-fiber classes are repairable or extendable
readout tasks.  Only the persistent class is carried forward as a named carrier
obstruction.
\end{proposition}

\begin{remark}[Finite off-carrier diagnostic instance]
In the finite carrier-routing instance,
\[
    |\mathcal R_{\mathrm{off}}|=1511.
\]
The primary obstruction counts are
\[
\begin{array}{c|r}
\text{class} & \text{count}\\
\hline
\mathrm{lens\text{-}extreme} & 315\\
\mathrm{scale\text{-}boundary} & 633\\
\mathrm{projection\text{-}chart} & 0\\
\mathrm{missing\text{-}fiber} & 563\\
\mathrm{persistent} & 0
\end{array}
\]
The box-counting estimates for the nonempty classes are
\[
\begin{array}{c|c}
\text{class} & D_{\mathrm{box}}\\
\hline
\mathrm{lens\text{-}extreme} & 1.330\\
\mathrm{scale\text{-}boundary} & 1.508\\
\mathrm{projection\text{-}chart} & \text{n/a}\\
\mathrm{missing\text{-}fiber} & 1.588\\
\mathrm{persistent} & \text{n/a}
\end{array}
\]
These values are finite diagnostics only.
\end{remark}

\begin{theorem}[Off-carrier obstruction contract]
The off-carrier set admits the following conservative contract.  If an
off-carrier residual is lens-extreme, it is routed to a lens refinement.  If it
is scale-boundary, the scale atlas must be extended.  If it is projection-chart,
a second chart must be tested.  If it is missing-fiber, the Hopf carrier family
must be refined.  Only if none of these explanations applies is the residual
kept as a persistent finite carrier obstruction.
\end{theorem}

\begin{proof}
The classes are defined by finite predicates that test bifocal stress, scale
support, carrier distance, and near-threshold carrier failure.  Each predicate
corresponds to a specific repair or refinement task.  Therefore a residual
cannot be promoted to persistent obstruction until the earlier finite routing
explanations have failed.
\end{proof}


% --- Closure Ledger and Release-Final Audit Insert ---

\section{Closure Ledger and Release-Final Audit for the Orbit--Hopf Chain}
\label{sec:closure-ledger-release-final-audit}

The preceding local readout chain has reached a natural audit boundary.  This
section records what has been certified, what remains only diagnostic, and what
is explicitly not claimed.  The ledger is a release-final bookkeeping
layer, not a new physical closure theorem.

\begin{definition}[Finite orbit--Hopf readout chain]
\label{def:finite-orbit-hopf-readout-chain}
The finite chain audited here is
\[
\begin{aligned}
&\text{alternating triolic orbit scan}
\longrightarrow
\text{core re-centring calibration}
\longrightarrow
\text{crystal--fractal readout}\\
&\longrightarrow
\text{residual classification}
\longrightarrow
\text{Hopf-linking readout}
\longrightarrow
\text{Hopf-stable residual carrier routing}\\
&\longrightarrow
\text{off-carrier obstruction split}.
\end{aligned}
\]
Each arrow is a finite readout transition under declared chart, sampling, and
guard assumptions.
\end{definition}

\begin{definition}[Closure ledger states]
\label{def:closure-ledger-states}
A stage may receive one of the following ledger states:
\[
\mathrm{defined},\quad
\mathrm{diagnostic},\quad
\mathrm{calibration\text{-}resolved},\quad
\mathrm{classified},\quad
\mathrm{carrier\text{-}resolved},\quad
\mathrm{persistent}.
\]
A state of \(\mathrm{persistent}\) is reserved for a named finite obstruction
that is not explained by calibration, scale extension, chart correction, lens
routing, or carrier refinement.
\end{definition}

\begin{proposition}[Finite local ledger result]
\label{prop:finite-local-ledger-result}
In the present finite audit instance, the chain records no persistent
finite-carrier obstruction:
\[
    \#\mathcal R_{\mathrm{persistent}} = 0.
\]
The carrier routing resolves
\[
    28705
\]
residuals and leaves
\[
    1511
\]
off-carrier residuals, all of which are classified as lens-extreme,
scale-boundary, or missing-fibre tasks.
\end{proposition}

\begin{proof}
The residual classification layer splits the multiscale residual set into
finite primary classes.  The Hopf-linking layer verifies unit-stable linking for
\[
    45/45
\]
tested fibre pairs.  The carrier-routing layer assigns residuals to this finite
Hopf-stable carrier family and records the off-carrier remainder.  The
off-carrier split then classifies that remainder as
\[
\begin{array}{c|r}
\text{class} & \text{count}\\
\hline
\mathrm{lens\text{-}extreme} & 315\\
\mathrm{scale\text{-}boundary} & 633\\
\mathrm{projection\text{-}chart} & 0\\
\mathrm{missing\text{-}fiber} & 563\\
\mathrm{persistent} & 0
\end{array}
\]
so no member of the off-carrier set is carried forward as a persistent finite
carrier obstruction.
\end{proof}

\begin{remark}[Numerical readout summary]
\label{rem:numerical-readout-summary}
The finite crystal--fractal and Hopf readouts record
\[
    \mathcal C_{\mathrm{crys}}^{\max}\approx 8.699e-17,
    \qquad
    D_{\mathrm{res}}\approx 2.435117,
\]
with
\[
    |\mathcal R_{\mathrm{multi}}|=30216.
\]
The residual classification counts are
\[
\begin{array}{c|r}
\text{primary class} & \text{count}\\
\hline
\mathrm{frac} & 27194\\
\mathrm{lens} & 3022\\
\mathrm{guard} & 0\\
\mathrm{noise} & 0
\end{array}
\]
The Hopf readout records mean and maximum unit-linking deviations
\[
    \overline{\Delta}_L\approx 1.1382e-04,
    \qquad
    \Delta_L^{\max}\approx 1.4636e-04.
\]
The carrier-resolved fraction is
\[
    \rho_H\approx 0.950,
\]
with normalised carrier-load entropy
\[
    H_H\approx 0.986.
\]
\end{remark}

\begin{theorem}[Release-final guard]
\label{thm:release-final-guard}
The finite orbit--Hopf chain is locally reviewable as a release-final
readout block with no persistent finite-carrier obstruction.  It does not imply
a global physical closure theorem, a through-core trajectory, or a global Hopf
fibration of the physical model.
\end{theorem}

\begin{proof}
Every stage in the chain is finite and chart-dependent.  The orbit scan and
perturbation atlas are metric diagnostics; the re-centring layer is a
calibration step; the crystal--fractal readout is a finite coherence and
scaling diagnostic; the residual classification layer is a finite alphabet;
the Hopf-linking layer is a sampled topological readout; and the carrier and
off-carrier layers are finite routing predicates.  Since no stage removes the
chart, sampling, and guard assumptions, the conclusion is a local
release-final ledger statement only.
\end{proof}

\begin{corollary}[Next release tasks]
\label{cor:next-release-tasks}
The remaining non-closure tasks are:
\[
\begin{array}{ll}
(1) & \text{extend the scale atlas for scale-boundary residuals},\\
(2) & \text{increase Hopf-carrier density for missing-fibre residuals},\\
(3) & \text{route lens-extreme residuals through bifocal/lens-stress refinement},\\
(4) & \text{insert a final version-specific DOI only at actual release},\\
(5) & \text{optionally rebuild the Historical Witness in the same release packet}.
\end{array}
\]
\end{corollary}


% --- v4.4 Roadmap Source Packet Insert ---

\section{v4.4 Roadmap}
\label{sec:v44-roadmap-refinement-op-audit}

This roadmap records the refinement and OP1/OP4 closure-candidate audit plan after the v4.3.0 local
orbit--Hopf readout release.  It is a roadmap and audit-contract layer, not a
new closure claim.  The v4.4 target is
\[
    \mathrm{v4.4}
    =
    \text{Refinement}
    +
    \text{OP1/OP4 closure-candidate audit}.
\]
The corresponding v4.4 target is reserved for possible OP1/OP4 closure
attempts, and only if the v4.4 audit contracts are satisfied.

\subsection{Starting point from the v4.3.0 local readout chain}
\label{subsec:v44-starting-point-v43-chain}

The v4.3.0 local chain consists of
\[
\begin{aligned}
&\text{Orbit Scan}
\to \text{Core Recentering}
\to \text{Crystal/Fractal Readout}\\
{}&\to \text{Residual Classification}
\to \text{Hopf Linking}\\
{}&\to \text{Carrier Routing}
\to \text{Off-Carrier Split}.
\end{aligned}
\]
Its finite ledger reports no persistent finite carrier obstruction in the
declared local readout chart:
\[
    \#\mathcal R_{\mathrm{persistent}}=0.
\]
The remaining explanatory classes are not ignored.  They become the v4.4
refinement tasks:
\[
    \#\mathcal R_{\mathrm{missing\ fiber}}=563,
    \qquad
    \#\mathcal R_{\mathrm{scale\ boundary}}=633,
    \qquad
    \#\mathcal R_{\mathrm{lens\ extreme}}=315.
\]
The interpretation is conservative: these numbers define finite readout tasks,
not physical trajectories through the protected core.

\subsection{Strand A: local refinement of the orbit--Hopf chain}
\label{subsec:v44-strand-a-refinement}

The first v4.4 strand refines the explanatory classes left by v4.3.0.

\paragraph{Dense Hopf carrier refinement.}
The missing-fibre class is tested by increasing the carrier family,
for example
\[
    N=10\longrightarrow 20\longrightarrow 40,
\]
while rechecking the finite Hopf-linking matrix.  The task succeeds only if the
missing-fibre count drops without producing non-unit linking pairs.

\paragraph{Scale atlas extension.}
The scale-boundary class is tested by extending the scan-scale atlas both
inward and outward.  The task succeeds only if boundary residuals are reduced
or explicitly reclassified while the residual scaling diagnostics remain
stable within declared tolerances.

\paragraph{Lens-extreme local refitting.}
The lens-extreme class is routed through local bifocal refitting.  A local
focus-pair correction
\[
    (F_-,F_+)
    \longmapsto
    (\widehat F_-,\widehat F_+)
\]
may reduce bifocal stress.  If it does not, the residual is not erased; it is
carried as explicit lens debt.

\paragraph{Dual projection chart audit.}
The Hopf-linking and carrier-routing diagnostics are repeated in a second
projection chart or pole choice.  The task succeeds only if the finite
linking/carrier conclusions are chart-stable.

\paragraph{v4.4 refinement ledger.}
The refinement strand closes with a ledger splitting the local residuals into
resolved, calibrated, reducible, and explicitly open classes.  It must not
create a new unexplained persistent finite carrier obstruction.

\subsection{Strand B: OP1/OP4 ingress and closure-candidate audit}
\label{subsec:v44-strand-b-op-audit}

The second v4.4 strand tests whether OP1 and OP4 can be legitimately routed
into the refined orbit--Hopf chain.  It does not close OP1 or OP4.

\begin{definition}[OP ingress contract]
\label{def:v44-op-ingress-contract}
For \(j\in\{1,4\}\), an OP\(j\) ingress contract is an explicitly declared
map
\[
    \mathcal D_{\mathrm{OP}j}
    \longrightarrow
    \mathcal R_{\mathrm{orbit/Hopf}}^{\mathrm{OP}j}
\]
from the OP\(j\) obstruction grammar into the refined finite orbit--Hopf
readout chain.  The map is admissible only if it preserves the protected-core
guard, introduces no unbudgeted HC/MM debt, and does not identify a finite
readout residual with a physical through-core trajectory.
\end{definition}

\begin{definition}[OP closure-candidate status]
\label{def:v44-op-closure-candidate-status}
After ingress and routing, OP\(j\) receives exactly one status:
\[
    \mathrm{closed\text{-}candidate},
    \quad
    \mathrm{reduced},
    \quad
    \mathrm{blocked},
    \quad
    \mathrm{not\ applicable}.
\]
The status \(\mathrm{closed\text{-}candidate}\) is permitted only when
\[
    \mathcal R_{\mathrm{persistent}}^{\mathrm{OP}j}=0
\]
under the refined carrier, scale, lens, and dual-chart audits, and when no new
unbudgeted debt is introduced.
\end{definition}

\begin{proposition}[v4.4 does not close OP1 or OP4]
\label{prop:v44-no-op-closure-claim}
The v4.4 line may define OP1/OP4 ingress contracts and closure-candidate
ledgers, but it does not assert OP1 or OP4 closure.
\end{proposition}

\begin{proof}
The v4.4 tasks are refinement, routing, and audit tasks.  Their output is a
status ledger.  A closure theorem would require a further proof layer using the
closed-candidate status as input.  That proof layer is reserved for v4.4.
\end{proof}

\subsection{v4.4 handoff criteria}
\label{subsec:v45-handoff-criteria}

A v4.4 closure attempt for OP1 or OP4 is allowed only when the v4.4 status is
\(\mathrm{closed\text{-}candidate}\).  The minimal handoff criteria are:
\[
\begin{array}{ll}
\text{(H1)} & \text{protected-core guard preserved;}\\
\text{(H2)} & \text{no new unbudgeted HC/MM, lens, carrier, scale, or guard debt;}\\
\text{(H3)} & \mathcal R_{\mathrm{persistent}}^{\mathrm{OP}j}=0
              \text{ after refined routing;}\\
\text{(H4)} & \text{dual-chart / carrier stability holds;}\\
\text{(H5)} & \text{all remaining classes are resolved, calibrated, or explicitly open.}
\end{array}
\]
If any criterion fails, v4.4 may still report reduction or blockage, but not
closure.

\begin{theorem}[v4.4 roadmap guard]
\label{thm:v44-roadmap-guard}
The v4.4 roadmap is a conservative preparation layer.  It can strengthen the
local orbit--Hopf readout chain and can prepare OP1/OP4 closure candidates, but
it cannot by itself promote OP1, OP4, or the global physical model to a closure
theorem.
\end{theorem}

\begin{proof}
Each v4.4 task is formulated as a finite refinement, chart-stability,
residual-routing, or ingress-contract test.  None of these tasks supplies the
additional proof layer needed for global OP closure.  Therefore the only
admissible v4.4 conclusion is a ledger status and a v4.4 handoff decision.
\end{proof}


% --- Dense Hopf Carrier Refinement Insert ---

\section{Dense Hopf Carrier Refinement}
\label{sec:dense-hopf-carrier-refinement}

This section begins the v4.4 refinement branch.  The v4.3 carrier routing left a
finite missing-fibre class.  The present test asks whether that class is a
true obstruction or a finite carrier-density artefact.

\begin{definition}[Dense carrier family]
Let
\[
    \Gamma_H^{(N)}=\{\gamma_1,\ldots,\gamma_N\}
\]
be a finite Hopf-carrier family sampled in the same stereographic chart as the
v4.3 finite Hopf-linking readout.  The dense refinement tests
\[
    N\in\{10,20,40\}
\]
while keeping the v4.3 carrier threshold \(\tau_H\) fixed.  A residual is
resolved at density \(N\) when
\[
    d_H^{(N)}(x)\leq \tau_H.
\]
\end{definition}

\begin{definition}[Missing-fibre refinement predicate]
Let \(\mathcal R_{\mathrm{miss}}^{(10)}\) denote the v4.3 missing-fibre
residual class.  The dense refinement residue is
\[
    \mathcal R_{\mathrm{miss}}^{(N)}
    :=
    \{x\in\mathcal R_{\mathrm{miss}}^{(10)}:d_H^{(N)}(x)>\tau_H\}.
\]
The missing-fibre class is called density-resolved at level \(N\) if
\[
    |\mathcal R_{\mathrm{miss}}^{(N)}|=0.
\]
\end{definition}

\begin{proposition}[Carrier-density refinement guard]
Dense carrier refinement may reduce a missing-fibre residual class only if the
Hopf-linking stability matrix remains unit-stable at the refined carrier
density.  Otherwise the refinement trades a metric routing issue for a
topological carrier issue.
\end{proposition}

\begin{remark}[Finite dense-carrier diagnostic]
Using the fixed v4.3 threshold
\[
    \tau_H = 0.6835,
\]
the finite dense-carrier run produced the following ledger:
\[
\begin{array}{c|r|r|r|r}
N & \text{stable pairs} & \text{all off-carrier} &
\text{old missing remaining} & \max\Delta_L\\
\hline
10 & 45/45 & 1512 & 563 & 0.0016\\
20 & 190/190 & 1283 & 444 & 0.0016\\
40 & 780/780 & 1159 & 391 & 0.0020
\end{array}
\]
The test is finite and local.  It does not assert OP1 or OP4 closure.
\end{remark}

\begin{theorem}[Dense carrier handoff]
If the missing-fibre class decreases under increasing \(N\) while the Hopf
linking matrix remains unit-stable, then the missing-fibre class is a carrier
resolution task rather than a persistent topology obstruction.  Any remaining
points are passed to the v4.4 refinement ledger, not to OP closure.
\end{theorem}

\begin{proof}
The threshold \(\tau_H\) is held fixed from the v4.3 carrier readout.  Hence a
newly resolved point is explained by a denser carrier family rather than by a
looser acceptance threshold.  Unit-stable linking prevents the denser carrier
family from introducing a new topological defect.  Remaining points are exactly
those still outside the fixed carrier threshold after refinement.
\end{proof}


% --- Scale Atlas Extension Insert ---

\section{Scale Atlas Extension for Boundary Residuals}
\label{sec:scale-atlas-extension-boundary-residuals}

The off-carrier split identifies a finite scale-boundary residual class.  This
class is not a persistent obstruction until the scale atlas has been extended.
The present layer enlarges the scan-scale window and tests whether formerly
boundary-supported sites acquire interior scale support.

\begin{definition}[Extended scale atlas]
Let the original scale window be
\[
    S_0=\{1.00,0.82,0.67,0.55,0.45,0.37\}.
\]
The v4.4 extension uses
\[
    S_1=\{1.30,1.15,1.00,0.82,0.67,0.55,0.45,0.37,0.30,0.25,0.20\}.
\]
For a residual site \((\ell,m)\), define its extended support
\[
    \sigma_{S_1}(\ell,m)
    :=
    \#\{s\in S_1: x_{\ell m}^{s}\in\mathcal R_s\}.
\]
A former scale-boundary site is called extension-resolved when
\[
    \sigma_{S_1}(\ell,m)>4.
\]
\end{definition}

\begin{definition}[Scale-window dimension stability]
Let \(D_{S_0}\) and \(D_{S_1}\) be finite box-counting estimates for the
residual readout set in the original and extended scale windows.  The scale
stability span is
\[
    \Delta D_{\mathrm{scale}}
    :=
    \max(D_{S_0},D_{S_1},D_{\mathrm{outer}},D_{\mathrm{inner}})
    -
    \min(D_{S_0},D_{S_1},D_{\mathrm{outer}},D_{\mathrm{inner}}),
\]
where \(D_{\mathrm{outer}}\) and \(D_{\mathrm{inner}}\) are the corresponding
estimates for the newly added outer and inner scale windows.
\end{definition}

\begin{proposition}[Scale-boundary guard]
A residual classified as scale-boundary cannot be promoted to persistent
obstruction until it remains boundary-supported after an explicit scale-atlas
extension.  If its extended support becomes interior, it is a finite-window
artifact rather than a persistent carrier obstruction.
\end{proposition}

\begin{remark}[Finite scale-extension instance]
In the finite v4.4 diagnostic, the old off-carrier scale-boundary class had
\[
    N_{\mathrm{scale}}^{0}=633
\]
rows, supported on
\[
    N_{\mathrm{sites}}^{0}=633
\]
distinct residual sites.  After extension to \(S_1\),
\[
    N_{\mathrm{rows,resolved}}=633,
    \qquad
    N_{\mathrm{rows,unresolved}}=0.
\]
At the site level,
\[
    N_{\mathrm{sites,resolved}}=633,
    \qquad
    N_{\mathrm{sites,unresolved}}=0.
\]
The box-counting estimates were
\[
\begin{array}{c|c}
\text{window} & D_{\mathrm{box}}\\
\hline
S_0 & 2.435\\
S_1 & 2.520\\
\text{outer added} & 2.171\\
\text{inner added} & 2.291
\end{array}
\]
with scale-stability span
\[
    \Delta D_{\mathrm{scale}}\approx 0.348.
\]
\end{remark}

\begin{theorem}[Scale-extension contract]
The scale atlas extension may reduce a scale-boundary class to a finite-window
artifact or leave it as an explicitly unresolved scale task.  It does not close
OP1 or OP4, and it does not by itself prove a global fractal dimension.  It
only decides whether the previous scale-boundary diagnosis survives a larger
declared scale window.
\end{theorem}


% --- Lens-Extreme Local Refitting Insert ---

\section{Lens-Extreme Local Refitting}
\label{sec:lens-extreme-local-refitting}

The off-carrier split contains a finite lens-extreme class.  These points are
not topological defects until a local bifocal refit has failed.  The present
layer tests whether the lens-extreme residuals are explained by a bounded
focal-parameter correction.

\begin{definition}[Bounded symmetric bifocal refit]
For a scan point \(x\) on ring \(\ell\) and scale \(s\), let
\[
    F_-=-c\,e_\ell,\qquad F_+=c\,e_\ell
\]
be the symmetric foci along the ring axis.  A bounded refit chooses
\[
    c=\kappa a_s,\qquad \kappa\in[\kappa_{\min},\kappa_{\max}],
\]
and an effective semi-major value \(\widehat a_s\) so as to minimise the finite
bifocal residual
\[
    \Delta_{\mathrm{lens}}(x;\widehat a_s,\widehat c_s)
    :=
    \bigl|d(x,\widehat F_-)+d(x,\widehat F_+)-2\widehat a_s\bigr|.
\]
The refit is scale-local when \((\widehat a_s,\widehat c_s)\) is fitted per
scale, and ring-local when it is fitted per \((s,\ell)\) group.
\end{definition}

\begin{definition}[Lens-refit statuses]
A lens-extreme point is called:
\[
\begin{array}{c|l}
\text{status} & \text{finite condition}\\
\hline
\mathrm{calibrated} &
\Delta_{\mathrm{lens}}\leq \tau_{\mathrm{lens}}\\
\mathrm{reduced} &
\Delta_{\mathrm{lens}}\text{ is reduced by at least }50\%\text{ but remains above threshold}\\
\mathrm{persistent} &
\text{neither calibrated nor reduced}
\end{array}
\]
The guard margin must remain positive in every case.
\end{definition}

\begin{proposition}[Lens-refit guard]
A lens-extreme residual is a bifocal-chart task before it is a carrier
obstruction.  If a bounded local refit calibrates or substantially reduces the
defect, the point may not be counted as persistent topological debt.
\end{proposition}

\begin{remark}[Finite lens-refit instance]
In the finite v4.4 diagnostic, the lens-extreme class had
\[
    N_{\mathrm{lens}}=315
\]
points.  The scale-local refit used
\[
    N_{\mathrm{scale\ groups}}=1
\]
scale groups and the ring-local fallback used
\[
    N_{\mathrm{ring\ groups}}=0
\]
groups.  The final status counts were
\[
\begin{array}{c|r}
\text{status} & \text{count}\\
\hline
\mathrm{calibrated} & 315\\
\mathrm{reduced} & 0\\
\mathrm{persistent} & 0
\end{array}
\]
The mean original bifocal defect was
\[
    \overline{\Delta}_{\mathrm{orig}}\approx 0.0764,
\]
and the mean scale-refit defect was
\[
    \overline{\Delta}_{\mathrm{scale}}\approx 0.0033.
\]
The median scale-refit improvement was
\[
    I_{\mathrm{scale}}\approx 0.968.
\]
\end{remark}

\begin{theorem}[Lens-extreme refinement contract]
The lens-extreme class is resolved for the finite v4.4 ledger only to the
extent that its points are calibrated or reduced by bounded bifocal refitting.
Any remaining persistent lens debt is carried forward as lens debt, not as OP1
or OP4 closure evidence and not as a through-core event.
\end{theorem}


% --- Dual Projection Chart Audit Insert ---

\section{Dual Projection Chart Audit}
\label{sec:dual-projection-chart-audit}

The Hopf-linking and carrier-routing tests use a stereographic projection from
\(S^3\) into \(\mathbb R^3\).  A single projection chart is not enough to
exclude chart artefacts.  The v4.4 audit therefore repeats the finite
Hopf-linking and carrier-routing tests in several stereographic charts.

\begin{definition}[Projection-chart family]
Let \(X=(X_1,X_2,X_3,X_4)\in S^3\).  For a coordinate \(q\in\{1,2,3,4\}\)
and sign \(\sigma\in\{+1,-1\}\), define the stereographic chart
\[
    \pi_{q,\sigma}(X)
    :=
    \frac{(X_i)_{i\neq q}}{1-\sigma X_q}
    \in\mathbb R^3.
\]
A chart is admissible for the finite fibre family when
\[
    \min(1-\sigma X_q)>0
\]
on all sampled fibre points.
\end{definition}

\begin{definition}[Chart-stability predicate]
For every tested chart \(c\), let
\[
    L_{ij}^{(c)}
\]
be the finite linking matrix and let
\[
    \rho_H^{(c)}
\]
be the carrier-resolved residual fraction.  The chart audit is passed when:
\[
    \bigl||L_{ij}^{(c)}|-1\bigr|\leq\tau_L
    \quad\text{for all tested }i\neq j,c,
\]
the projection denominators stay positive, and
\[
    \max_c\rho_H^{(c)}-\min_c\rho_H^{(c)}
\]
remains below the declared chart-variation tolerance.
\end{definition}

\begin{proposition}[Projection-chart guard]
If the Hopf-linking matrix remains unit-stable and the carrier-resolved
fraction remains stable across admissible charts, then the finite carrier
routing is not explained solely by the initial stereographic projection.  This
is a chart audit only.  It is not a global Hopf theorem and not OP1 or OP4
closure.
\end{proposition}

\begin{remark}[Finite chart-audit instance]
The v4.4 finite audit tested the charts
\[
    X4^+,\quad X4^-,\quad X3^+,\quad X3^-.
\]
Across all charts, the unit-linking condition returned
\[
    \mathrm{unstable\ pairs}=0.
\]
The carrier-resolved fraction ranged from
\[
    \rho_H^{\min}\approx 0.950
\]
to
\[
    \rho_H^{\max}\approx 0.953,
\]
with span
\[
    \Delta\rho_H\approx 0.003.
\]
The maximum linking deviation over all charts was
\[
    \Delta_L^{\max}\approx 0.0001.
\]
The minimum projection denominator was
\[
    d_{\pi}^{\min}\approx 0.3596.
\]
The declared chart-stability predicate returned
\[
    \mathrm{chart\ stable}=true.
\]
\end{remark}

\begin{theorem}[Dual-chart audit contract]
A carrier-routing result may be used in the v4.4 closure-candidate ledger only
if it survives the projection-chart audit.  If it fails, the correct conclusion
is chart debt, not OP closure and not physical through-core motion.
\end{theorem}


% --- v4.4 Refinement Ledger Insert ---

\section{v4.4 Refinement Ledger for the Orbit--Hopf Chain}
\label{sec:v44-refinement-ledger-orbit-hopf-chain}

The v4.4 refinement strand consolidates four finite diagnostics:
dense Hopf-carrier refinement, scale-atlas extension, lens-extreme local
refitting, and dual projection-chart audit.  This ledger is a bookkeeping
layer.  It does not close OP1 or OP4.

\begin{definition}[Refinement-ledger statuses]
For each finite residual class, the ledger records:
\[
    N_{\mathrm{in}},\qquad
    N_{\mathrm{resolved}},\qquad
    N_{\mathrm{remaining}}.
\]
A class is called \emph{resolved} when its residuals are explained by the
declared finite refinement.  It is called \emph{reduced-open} when the
refinement reduces but does not eliminate it.  It is called \emph{passed} when
a guard or audit predicate is satisfied.
\end{definition}

\begin{proposition}[v4.4 Strang-A refinement ledger]
The v4.4 refinement diagnostics give the following finite ledger:
\[
\begin{array}{c|r|r|r|c}
\text{class} & N_{\mathrm{in}} & N_{\mathrm{resolved}} & N_{\mathrm{remaining}} & \text{status}\\
\hline
\mathrm{missing\text{-}fiber} & 563 & 172 & 391 & \mathrm{reduced\text{-}open}\\
\mathrm{scale\text{-}boundary} & 633 & 633 & 0 & \mathrm{resolved}\\
\mathrm{lens\text{-}extreme} & 315 & 315 & 0 & \mathrm{resolved}\\
\mathrm{projection\text{-}chart} & 4 & 4 & 0 & \mathrm{passed}
\end{array}
\]
Thus the scale-boundary and lens-extreme classes are resolved by the v4.4
finite refinements, the projection-chart audit passes, and the only remaining
Strang-A task is the finite carrier-density remainder.
\end{proposition}

\begin{remark}[Numerical diagnostics]
For dense Hopf-carrier refinement, the \(N=40\) test retained
\[
    780
\]
stable Hopf pairs out of
\[
    780
\]
tested pairs, with maximum linking deviation
\[
    \Delta_L^{\max}\approx 0.0020.
\]
For the scale-atlas extension,
\[
    D_{S_0}\approx 2.435,\qquad
    D_{S_1}\approx 2.520,\qquad
    \Delta D_{\mathrm{scale}}\approx 0.348.
\]
For lens-extreme refitting, the mean bifocal defect moved from
\[
    \overline{\Delta}_{\mathrm{orig}}\approx 0.0764
\]
to
\[
    \overline{\Delta}_{\mathrm{refit}}\approx 0.0033,
\]
with median improvement
\[
    I_{\mathrm{lens}}\approx 0.968.
\]
The dual projection-chart audit returned
\[
    \mathrm{chart\ stable}=true,
\]
with carrier-resolved fraction span
\[
    \Delta\rho_H\approx 0.0029.
\]
\end{remark}

\begin{theorem}[Refinement-ledger guard]
The v4.4 refinement ledger may be used as input for the OP1/OP4
closure-candidate audit only as a finite diagnostic ledger.  It does not
itself close OP1 or OP4.  The remaining carrier-density class must be carried
forward as an explicit finite task unless a later denser-carrier test resolves
it.
\end{theorem}

\begin{proof}
Each entry in the ledger is produced by a finite diagnostic predicate:
carrier-density refinement, scale-window extension, bounded lens refit, and
projection-chart stability.  These predicates reduce or classify residual
classes in the local orbit--Hopf readout chain.  None of them is an OP1 or OP4
closure theorem.  Therefore the ledger can only prepare the later
closure-candidate audit.
\end{proof}


% --- OP1/OP4 Ingress Contract Insert ---

\section{OP1/OP4 Ingress Contract}
\label{sec:v44-op1-op4-ingress-contract}

The v4.4 refinement ledger prepares, but does not close, the OP1/OP4 line.
This section defines the ingress contract that allows OP1 and OP4 to be tested
against the refined orbit--Hopf readout chain.  The contract is conservative:
it may produce an ingress-admissible or reduced-open status, but it does not
assert OP1 or OP4 closure.

\begin{definition}[Refined orbit--Hopf ledger available to OP1/OP4]
Let
\[
    \mathcal L_A^{4.4}
\]
be the v4.4 Strand-A refinement ledger.  It records the finite residual
reductions
\[
    \mathcal R_{\mathrm{scale}}: 633\to0,\qquad
    \mathcal R_{\mathrm{lens}}:315\to0,
\]
the chart-stability audit
\[
    \mathcal R_{\mathrm{chart}}:4\to0,
\]
and the remaining carrier-density task
\[
    \mathcal R_{\mathrm{carrier}}^{\mathrm{rem}}=391.
\]
This ledger is the only admissible local orbit--Hopf input for the OP1/OP4
ingress contract.
\end{definition}

\begin{definition}[OP1 ingress side]
The OP1 ingress side is the selector-domain or pre-route side of the coupled
OP1/OP4 selector problem.  An OP1 ingress datum is a guarded map
\[
    \iota_1:\mathcal D_{\mathrm{OP1}}
    \longrightarrow
    P(\rho_{1/4}^{\mathrm{pre}})
\]
together with a selector-domain component
\[
    C_{1/4}^{\mathrm{sel,OP1}}.
\]
It is admissible only if it preserves the protected-core guard, does not create
new HC/MM debt, and does not use a finite residual as evidence for a physical
through-core trajectory.
\end{definition}

\begin{definition}[OP4 ingress side]
The OP4 ingress side is the backtranslation and post-residual stability side.
An OP4 ingress datum is a guarded map
\[
    \iota_4:\mathcal D_{\mathrm{OP4}}
    \longrightarrow
    P(\rho_{1/4}^{\mathrm{shift}})
\]
together with a \(B_4\)-compatible post-selector residual
\[
    R_{1/4}^{\mathrm{post,OP4}}.
\]
It is admissible only if backtranslation is stable, the post-selector residual
does not inflate, and the chart-stability audit remains valid.
\end{definition}

\begin{definition}[Coupled OP1/OP4 selector interface]
The coupled selector interface is the two-sided comparison object
\[
    C_{1/4}^{\mathrm{sel}}:
    P(\rho_{1/4}^{\mathrm{pre}})
    \times
    P(\rho_{1/4}^{\mathrm{shift}})
    \longrightarrow
    E_{1/4}^{\mathrm{sel}}.
\]
It is a comparison tensor, not a raw same-variance contraction.  Thus any
active--active or readout--readout comparison must be mediated by the declared
selector/comparison tensor itself.
\end{definition}

\begin{definition}[OP1/OP4 ingress statuses]
After applying the ingress contract to \(j\in\{1,4\}\), the status must be one
of
\[
    \mathrm{ingress\text{-}admissible},
    \quad
    \mathrm{reduced\text{-}open},
    \quad
    \mathrm{blocked},
    \quad
    \mathrm{not\ applicable}.
\]
The stronger status \(\mathrm{closed\text{-}candidate}\) is not licensed unless
\[
    \mathcal R_{\mathrm{carrier}}^{\mathrm{rem}}=0
\]
or unless the nonzero carrier-density remainder is separately certified
harmless for the specific OP ingress being tested.
\end{definition}

\begin{proposition}[v4.4 OP1/OP4 ingress ledger]
With the current Strand-A refinement ledger, OP1 and OP4 receive the following
ingress statuses:
\[
\begin{array}{c|c|c}
\text{OP} & \text{ingress side} & \text{status}\\
\hline
\mathrm{OP1} & \text{selector-domain / pre-route} &
\mathrm{ingress\text{-}admissible/reduced\text{-}open}\\
\mathrm{OP4} & \text{backtranslation / post-residual} &
\mathrm{ingress\text{-}admissible/reduced\text{-}open}\\
\mathrm{OP1/OP4} & \text{coupled selector interface} &
\mathrm{closed\text{-}candidate\ not\ licensed}
\end{array}
\]
The obstruction is not scale debt, lens debt, or chart debt.  The only named
blocker is the remaining finite carrier-density task.
\end{proposition}

\begin{proof}
By the v4.4 refinement ledger, scale-boundary debt has been reduced to zero,
lens-extreme debt has been reduced to zero, and projection-chart debt has been
reduced to zero.  The remaining task is
\[
    \mathcal R_{\mathrm{carrier}}^{\mathrm{rem}}=391.
\]
Therefore the OP1 and OP4 ingress maps may be defined and audited, but neither
side may be promoted to a closure-candidate status.  The nonzero carrier-density
remainder must be carried explicitly.
\end{proof}

\begin{theorem}[OP1/OP4 ingress guard]
The OP1/OP4 ingress contract licenses an audit, not a closure.  A v4.4 closure
attempt is blocked until the remaining carrier-density task is either resolved
or separately certified harmless for the relevant OP ingress.  No v4.4
statement in this section closes OP1 or OP4.
\end{theorem}


% --- OP1/OP4 Guard Preservation Audit Insert ---

\section{OP1/OP4 Guard Preservation Audit}
\label{sec:v44-op1-op4-guard-preservation-audit}

The OP1/OP4 ingress contract defines admissible entry into the refined
orbit--Hopf chain.  This section audits whether that ingress preserves the
standing guards.  It does not close OP1 or OP4.

\begin{definition}[Guard-preservation register]
For \(j\in\{1,4\}\), an OP\(j\) ingress is guard-preserving when all of the
following finite predicates hold:
\[
\begin{array}{ll}
(G_1) & \text{protected-core guard is preserved;}\\
(G_2) & \text{no scale-boundary debt is imported;}\\
(G_3) & \text{no persistent lens debt is imported;}\\
(G_4) & \text{projection-chart stability is preserved;}\\
(G_5) & \text{no new HC/MM debt is created.}
\end{array}
\]
OP1 additionally requires selector-domain admissibility.  OP4 additionally
requires post-residual non-inflation to be audited in the next residual-routing
step.
\end{definition}

\begin{definition}[Carrier-density visibility condition]
The remaining carrier-density task
\[
    \mathcal R_{\mathrm{carrier}}^{\mathrm{rem}}=391
\]
is not a guard failure, provided it is kept explicit and is not used as
closure evidence.  It is, however, a closure-candidate blocker until it is
resolved or certified harmless.
\end{definition}

\begin{proposition}[Guard-preservation outcome]
In the finite v4.4 audit, OP1 and OP4 preserve the inherited guards:
\[
\begin{array}{c|c|c|c}
\text{OP} & \text{guards tested} & \text{failed guards} & \text{status}\\
\hline
\mathrm{OP1} & 6 & 0 & \mathrm{guard\text{-}preserved/closure\text{-}blocked}\\
\mathrm{OP4} & 6 & 0 & \mathrm{guard\text{-}preserved/closure\text{-}blocked}
\end{array}
\]
The blocker is not scale debt, lens debt, chart debt, or new HC/MM debt.  The
blocker is the visible carrier-density remainder.
\end{proposition}

\begin{proof}
The v4.4 refinement ledger gives
\[
    \mathcal R_{\mathrm{scale}}=0,\qquad
    \mathcal R_{\mathrm{lens}}=0,\qquad
    \mathcal R_{\mathrm{chart}}=0.
\]
The OP1/OP4 ingress contract introduces the comparison object through the
declared selector interface and does not introduce a raw same-variance
contraction or new HC/MM primitive.  Hence no new debt is created by ingress.
The remaining task is
\[
    \mathcal R_{\mathrm{carrier}}^{\mathrm{rem}}=391,
\]
which is explicit.  Therefore the guards are preserved, but closure is blocked.
\end{proof}

\begin{theorem}[Guard-preservation handoff]
The v4.4 guard-preservation audit licenses the next step, namely OP1/OP4
residual routing.  It does not license OP1 or OP4 closure.  A v4.4 closure
attempt remains blocked until the remaining carrier-density task is resolved
or separately certified harmless for the OP being tested.
\end{theorem}


% --- OP1/OP4 Residual Routing Audit Insert ---

\section{OP1/OP4 Residual Routing Audit}
\label{sec:v44-op1-op4-residual-routing-audit}

The guard-preservation audit shows that OP1 and OP4 may enter the refined
orbit--Hopf chain without creating new debt.  The remaining task is to route
the explicit carrier-density remainder.  This section classifies that
remainder against OP1 and OP4.  It does not close either OP.

\begin{definition}[Carrier-density routing remainder]
Let
\[
    \mathcal R_{\mathrm{carrier}}^{\mathrm{rem}}
\]
be the remaining finite carrier-density class after the v4.4 refinement
ledger.  In the present audit,
\[
    |\mathcal R_{\mathrm{carrier}}^{\mathrm{rem}}|=391.
\]
This set is routed separately against the OP1 selector-domain side and the OP4
backtranslation/post-residual side.
\end{definition}

\begin{definition}[OP1 residual-routing statuses]
For OP1, a carrier-density residual is assigned one of two statuses:
\[
\mathrm{OP1\text{-}bounded\text{-}open},
\qquad
\mathrm{OP1\text{-}blocking}.
\]
The bounded-open status means that the residual is near enough to the carrier
collar to be a candidate for a later harmlessness certificate.  The blocking
status means that it still blocks selector-domain closure.  Neither status is
closure evidence.
\end{definition}

\begin{definition}[OP4 residual-routing statuses]
For OP4, a carrier-density residual is assigned one of two statuses:
\[
\mathrm{OP4\text{-}harmless\text{-}candidate},
\qquad
\mathrm{OP4\text{-}open}.
\]
A harmless candidate still requires a post-residual non-inflation certificate.
An open residual remains a carrier-density task.
\end{definition}

\begin{proposition}[Finite residual-routing ledger]
The finite OP1/OP4 residual-routing audit returns
\[
\begin{array}{c|r|r|r|c}
\text{side} & \text{input} & \text{bounded/harmless candidate} &
\text{blocking/open} & \text{closure candidate}\\
\hline
\mathrm{OP1} & 391 & 196 & 195 & \mathrm{no}\\
\mathrm{OP4} & 391 & 313 & 78 & \mathrm{no}\\
\mathrm{coupled} & 391 & 196 & 195 & \mathrm{no}
\end{array}
\]
Thus residual routing refines the remaining carrier-density class but does not
license OP1 or OP4 closure.
\end{proposition}

\begin{proof}
The routing audit does not remove residuals.  It assigns finite statuses to the
remaining carrier-density points.  OP1 retains a blocking selector-domain
component, and OP4 harmless candidates still require non-inflation
certificates.  The coupled selector cannot ignore the OP1 side.  Therefore no
closure-candidate status is licensed.
\end{proof}

\begin{theorem}[Residual-routing handoff]
The residual-routing audit creates the next v4.4 task ledger.  OP1 requires a
bounded-open harmlessness certificate and a refinement of blocking
carrier-density points.  OP4 requires harmless-candidate non-inflation
certificates and a refinement of open carrier-density points.  v4.4 closure
remains blocked until these certificates or refinements are supplied.
\end{theorem}


% --- OP4 Harmless-Candidate Certificate Test Insert ---

\section{OP4 Harmless-Candidate Certificate Test}
\label{sec:v44-op4-harmless-candidate-certificate-test}

The OP1/OP4 residual-routing audit identifies OP4 harmless candidates.  These
candidates are not automatically harmless.  They require a finite certificate
that post-residual routing is non-inflating and creates no new debt.

\begin{definition}[OP4 harmless-candidate certificate predicates]
Let \(x\) be an OP4 harmless candidate from the residual-routing ledger.  The
finite certificate predicates are:
\[
\begin{array}{ll}
(C_1) & \text{no scale, lens, or chart debt is imported;}\\
(C_2) & \text{the protected-core guard remains positive;}\\
(C_3) & \text{the post-residual routing proxy is non-inflating;}\\
(C_4) & \text{the point remains inside the route-stability collar.}
\end{array}
\]
If \(C_1\)-\(C_4\) all hold, the point is OP4-harmless-certified in the local
finite ledger.
\end{definition}

\begin{definition}[Coupled harmlessness condition]
An OP4-harmless-certified point is not automatically coupled-certified.  For
the coupled OP1/OP4 selector, it must also fail to be OP1-blocking.  Thus the
coupled condition adds
\[
    (C_5):\quad x\notin \mathcal R_{\mathrm{OP1,blocking}}.
\]
\end{definition}

\begin{proposition}[Finite OP4 certificate outcome]
The finite v4.4 certificate test returned
\[
\begin{array}{c|r}
\text{class} & \text{count}\\
\hline
\mathrm{OP4\ harmless\ candidates} & 313\\
\mathrm{OP4\ harmless\ certified} & 251\\
\mathrm{OP4\ uncertified\ candidates} & 62\\
\mathrm{OP4\ open\ carrier\ density} & 78\\
\mathrm{coupled\ harmless\ certified} & 196
\end{array}
\]
The certified OP4 points are local post-residual harmlessness certificates.
They are not OP4 closure and not coupled OP1/OP4 closure.
\end{proposition}

\begin{proof}
The certificate requires \(C_1\)-\(C_4\).  These predicates only show that a
candidate does not inflate the OP4 post-residual ledger in the finite route
collar.  They do not remove open OP4 carrier-density points, and they do not
remove OP1-blocking points.  Hence no closure-candidate status follows.
\end{proof}

\begin{theorem}[OP4 harmless-certificate handoff]
The OP4 harmless-candidate certificate test reduces the OP4 side by certifying
local harmlessness for the points satisfying \(C_1\)-\(C_4\).  A v4.4 closure
attempt remains blocked unless all OP4 open carrier-density points are resolved
and the OP1 side no longer contains coupled blocking residuals.
\end{theorem}


% --- OP1 Blocking Carrier-Density Refinement Insert ---

\section{OP1 Blocking Carrier-Density Refinement}
\label{sec:v44-op1-blocking-carrier-density-refinement}

The OP1/OP4 residual-routing audit leaves an OP1-blocking carrier-density
class.  This section tests whether those OP1 blockers are an artefact of
carrier density by refining the Hopf carrier family beyond \(N=40\).  This is
a reduction test, not an OP1 closure theorem.

\begin{definition}[OP1 dense-carrier refinement]
Let
\[
    \mathcal R_{\mathrm{OP1,block}}^{40}
\]
be the OP1-blocking carrier-density class after the \(N=40\) refinement.  A
dense-carrier refinement tests carrier families
\[
    N\in\{40,60,80,120\}
\]
and reroutes the same finite residual points against the denser carrier
families.  A point is called strictly density-resolved at \(N\) when it is
carrier-resolved in both tested projection charts.
\end{definition}

\begin{definition}[OP1 refinement statuses]
At the final tested carrier density, a point receives one of three statuses:
\[
\begin{array}{c|l}
\text{status} & \text{meaning}\\
\hline
\mathrm{density\text{-}resolved} & \text{resolved in both tested charts}\\
\mathrm{chart\text{-}dependent} & \text{resolved in one chart only}\\
\mathrm{persistent\text{-}blocking} & \text{unresolved in both charts}
\end{array}
\]
Only the first status is locally resolved.  The other two remain open tasks.
\end{definition}

\begin{proposition}[Finite OP1 density-refinement outcome]
The finite v4.4 OP1 refinement started from
\[
    |\mathcal R_{\mathrm{OP1,block}}^{40}|=195.
\]
At \(N=120\), the result was
\[
\begin{array}{c|r}
\text{status} & \text{count}\\
\hline
\mathrm{density\text{-}resolved} & 195\\
\mathrm{chart\text{-}dependent} & 0\\
\mathrm{persistent\text{-}blocking} & 0
\end{array}
\]
The sampled Hopf-linking guard at \(N=120\) returned
\[
    320
\]
stable pairs out of
\[
    320
\]
tested pairs.
\end{proposition}

\begin{proof}
The test changes the carrier density while keeping the residual input fixed.
A point that becomes resolved in both projection charts is classified as
density-resolved.  A point resolved in only one chart is chart-dependent and
therefore not closure evidence.  A point resolved in neither chart remains an
OP1-blocking carrier-density residual.  Since nonzero open tasks remain, OP1
closure is not licensed.
\end{proof}

\begin{theorem}[OP1 density-refinement handoff]
The OP1 blocking carrier-density refinement reduces or classifies the OP1
blocking class.  It does not close OP1.  v4.4 closure remains blocked unless
the chart-dependent and persistent OP1 carrier-density tasks are discharged by
a stronger carrier-refinement, chart-reconciliation, or harmlessness
certificate.
\end{theorem}


% --- OP4 Open and Uncertified Residual Refinement Insert ---

\section{OP4 Open and Uncertified Residual Refinement}
\label{sec:v44-op4-open-uncertified-residual-refinement}

The OP4 harmless-candidate certificate test leaves uncertified candidates and
open carrier-density residuals.  This section refines those remaining OP4
points using a denser carrier family and a stricter non-inflation check.  The
test reduces or certifies finite local OP4 residuals.  It does not close OP4.

\begin{definition}[OP4 unresolved refinement input]
Let
\[
    \mathcal R_{\mathrm{OP4,unres}}
\]
be the union of uncertified OP4 harmless candidates and OP4 open
carrier-density residuals.  In the finite v4.4 audit,
\[
    |\mathcal R_{\mathrm{OP4,unres}}|=140,
\]
with
\[
    62
\]
uncertified candidates and
\[
    78
\]
open carrier-density residuals.
\end{definition}

\begin{definition}[Refined OP4 certificate]
A point of \(\mathcal R_{\mathrm{OP4,unres}}\) is refined
OP4-harmless-certified at carrier density \(N\) when it is carrier-resolved in
both tested charts and its maximum new routing distance is non-inflating
relative to the previous route proxy.  If only a one-chart or weakened
improvement condition holds, the point is a refined reduced candidate.  If no
such condition holds, it remains persistent open.
\end{definition}

\begin{proposition}[Finite OP4 unresolved-refinement outcome]
At \(N=160\), the OP4 unresolved-refinement ledger returned
\[
\begin{array}{c|r}
\text{status} & \text{count}\\
\hline
\mathrm{refined\ harmless\ certified} & 140\\
\mathrm{refined\ reduced\ candidate} & 0\\
\mathrm{persistent\ open} & 0
\end{array}
\]
The sampled Hopf-linking guard at \(N=160\) returned
\[
    320
\]
stable pairs out of
\[
    320
\]
tested pairs.
\end{proposition}

\begin{proof}
The refinement keeps the OP4 unresolved input fixed and changes only the
carrier density and chart test.  A point certified in both charts and
non-inflating is locally harmless for the finite OP4 post-residual ledger.
One-chart or weakly improved points remain candidates.  Points failing both
conditions remain open.  Since any reduced or persistent residual remains an
open task, OP4 closure is not licensed.
\end{proof}

\begin{theorem}[OP4 unresolved-refinement handoff]
The OP4 unresolved-refinement test may add local harmlessness certificates to
the OP4 ledger.  It does not close OP4.  Any remaining refined reduced
candidate or persistent open residual must be carried into the final v4.4
closure-candidate ledger as an explicit blocker or open task.
\end{theorem}


% --- v4.4 OP1/OP4 Closure-Candidate Ledger Insert ---

\section{v4.4 OP1/OP4 Closure-Candidate Ledger}
\label{sec:v44-op1-op4-closure-candidate-ledger}

This ledger consolidates the OP1 and OP4 audits in v4.4.  It is a
closure-candidate ledger only.  It does not close OP1 or OP4.  Its purpose is
to decide whether a v4.4 closure attempt is allowed.

\begin{definition}[Closure-attempt gate]
The v4.4 closure-attempt gate is open only if the v4.4 ledger has no known
open OP1 blocking residual, no uncertified OP1 bounded-open residual, no OP4
open residual, no scale debt, no lens debt, no chart debt, and no new HC/MM
debt.  Opening this gate does not prove closure.  It only licenses the next
version to attempt closure.
\end{definition}

\begin{proposition}[OP1/OP4 closure-candidate ledger]
The finite v4.4 OP1/OP4 ledger is
\[
\begin{array}{c|r|r|r|c}
\text{ledger item} & N_{\mathrm{in}} & N_{\mathrm{cert/res}} &
N_{\mathrm{open}} & \text{status}\\
\hline
\mathrm{OP1\ bounded\text{-}open} & 196 & 196 & 0 & \mathrm{certified\text{-}local}\\
\mathrm{OP1\ blocking} & 195 & 195 & 0 & \mathrm{density\text{-}resolved}\\
\mathrm{OP4\ first\ candidates} & 313 & 251 & 62 & \mathrm{forwarded}\\
\mathrm{OP4\ open/uncertified} & 140 & 140 & 0 & \mathrm{refined\text{-}certified}\\
\mathrm{scale/lens/chart/new\ debt} & 0 & 0 & 0 & \mathrm{zero\text{-}debt}
\end{array}
\]
The coupled known-open remainder is
\[
    R_{\mathrm{known\ open}}^{1/4}=0.
\]
\end{proposition}

\begin{proof}
The OP1 blocking side is resolved by dense carrier refinement.  The OP1
bounded-open side is carried only as a local coupled harmless certificate, not
as closure.  OP4 first-stage certificates and the OP4 unresolved refinement
together certify the OP4 local residual side.  The guard-preservation audit
records zero scale, lens, chart, and new HC/MM debt.  Therefore the known-open
ledger is \(0\).  If this number is zero, v4.4 may attempt closure.
No closure is asserted here.
\end{proof}

\begin{theorem}[v4.4 handoff condition]
In the present v4.4 ledger, the v4.4 closure-attempt handoff status is
\[
    \mathrm{V45\_HANDOFF}.
\]
This means only that the known v4.4 finite blockers have been discharged or
locally certified.  It does not mean that OP1 or OP4 are closed.  The v4.4
document must still prove the actual closure theorem, if possible.
\end{theorem}

\begin{remark}[Numerical linking guards]
The OP1 \(N=120\) dense refinement sampled
\[
    320/320
\]
stable Hopf-linking pairs, with maximum deviation
\[
    \Delta_{L,\mathrm{OP1}}^{\max}\approx 0.0044.
\]
The OP4 \(N=160\) refinement sampled
\[
    320/320
\]
stable Hopf-linking pairs, with maximum deviation
\[
    \Delta_{L,\mathrm{OP4}}^{\max}\approx 0.0144.
\]
These are finite guard checks, not global Hopf theorems.
\end{remark}


% --- v4.4 Final Audit / Source-Freeze Candidate Insert ---

\section{v4.4 Final Audit and Source-Freeze Candidate}
\label{sec:v44-final-audit-source-freeze-candidate}

This section records the v4.4 final audit for the working source line.  It is a
source-freeze candidate, not a release freeze.  No DOI is assigned here, and no
OP1 or OP4 closure is claimed.

\begin{definition}[v4.4 source-freeze candidate]
A v4.4 source-freeze candidate is a working source state in which the v4.4
Strand-A refinement ledger and Strand-B OP1/OP4 audit ledger have been
integrated into one source-native Active Core and audited for consistency.  It
is not a public release object unless a later release step assigns release
metadata and rebuilds the public packet.
\end{definition}

\begin{proposition}[v4.4 final finite-audit ledger]
The v4.4 final finite-audit ledger records
\[
    R_{\mathrm{known\ open}}^{1/4}=0,
\]
with
\[
    R_{\mathrm{scale}}=0,\qquad
    R_{\mathrm{lens}}=0,\qquad
    R_{\mathrm{chart}}=0,\qquad
    B_{\mathrm{HC/MM,new}}=0.
\]
The v4.4 handoff status is
\[
    \mathrm{v4.4\ closure\text{-}attempt\ handoff}=\mathrm{allowed}.
\]
This means that the known v4.4 finite blockers have been discharged or locally
certified enough to permit a later closure attempt.  It does not mean closure.
\end{proposition}

\begin{theorem}[No-closure source-freeze guard]
The v4.4 source-freeze candidate may say that v4.4 is allowed to attempt OP1
and OP4 closure if the handoff gate is open.  It may not say that v4.4 closes
OP1 or OP4.  Any theorem in a later v4.4 line must re-establish its assumptions
and prove the closure statement independently.
\end{theorem}

\begin{proof}
The v4.4 line contains finite diagnostics, carrier refinements, local
harmlessness certificates, guard-preservation audits, and a closure-candidate
ledger.  These are sufficient to decide whether a later closure attempt is
permitted.  They are not identical to an OP1 or OP4 closure theorem.  Therefore
the correct source-freeze status is a guarded handoff, not a closure claim.
\end{proof}

\begin{remark}[Release guard]
This source-freeze candidate has no new DOI and no release-final claim.  A
future release requires a separate release audit, metadata insertion, public
PDF rebuild, and upload checklist.
\end{remark}


% --- v4.4 OP1/OP4 Closure-Attempt Intake Insert ---

\section{v4.4 OP1/OP4 Closure-Attempt Intake}
\label{sec:v45-op1-op4-closure-attempt-intake}

The v4.4 source-freeze candidate opens the gate for a v4.4 OP1/OP4 closure
attempt.  This section starts that attempt.  It does not close OP1 or OP4.
The purpose is to state the closure targets and the proof obligations that
must be discharged in v4.4.

\begin{definition}[v4.4 closure-attempt basis]
The v4.4 closure attempt is based on the v4.4 handoff ledger
\[
    R_{\mathrm{known\ open}}^{1/4}=0,
\]
with
\[
    R_{\mathrm{scale}}=0,\qquad
    R_{\mathrm{lens}}=0,\qquad
    R_{\mathrm{chart}}=0,\qquad
    B_{\mathrm{HC/MM,new}}=0.
\]
The handoff permits a closure attempt only because the known finite blockers
were discharged or locally certified in v4.4.  The handoff is not itself a
closure theorem.
\end{definition}

\begin{definition}[OP1 closure target]
The OP1 closure target is the selector-domain statement that the OP1 ingress
side carries no unaccounted selector-domain obstruction after applying the
declared orbit--Hopf refinement and coupled selector interface:
\[
    \mathcal D_{\mathrm{OP1}}
    \xrightarrow{\ \iota_1\ }
    P(\rho_{1/4}^{\mathrm{pre}})
    \xrightarrow{\ C_{1/4}^{\mathrm{sel}}\ }
    E_{1/4}^{\mathrm{sel}}
\]
with zero unaccounted OP1 residual.
\end{definition}

\begin{definition}[OP4 closure target]
The OP4 closure target is the post-residual/backtranslation statement that the
OP4 ingress side carries no unaccounted post-selector residual after applying
the declared orbit--Hopf refinement and coupled selector interface:
\[
    \mathcal D_{\mathrm{OP4}}
    \xrightarrow{\ \iota_4\ }
    P(\rho_{1/4}^{\mathrm{shift}})
    \xrightarrow{\ C_{1/4}^{\mathrm{sel}}\ }
    E_{1/4}^{\mathrm{sel}}
\]
with zero unaccounted OP4 residual.
\end{definition}

\begin{definition}[v4.4 proof-obligation ledger]
The v4.4 closure attempt must discharge the proof-obligation ledger
\[
\begin{array}{c|l|c}
\text{id} & \text{obligation} & \text{status}\\
\hline
P1 & \text{OP1 selector-domain promotion} & \mathrm{open}\\
P2 & \text{OP4 post-residual promotion} & \mathrm{open}\\
P3 & \text{coupled selector compatibility} & \mathrm{open}\\
P4 & \text{no-new-debt promotion} & \mathrm{open}\\
P5 & \text{finite-to-declared-grammar closure guard} & \mathrm{open}
\end{array}
\]
A closure theorem is not licensed until all five obligations are discharged.
\end{definition}

\begin{proposition}[Closure-attempt licensed, closure not proven]
The v4.4 handoff licenses the v4.4 closure attempt:
\[
    \mathrm{closure\text{-}attempt\ licensed}=true.
\]
However, at intake time the number of open v4.4 proof obligations is
\[
    N_{\mathrm{open}}^{P}=5.
\]
Therefore
\[
    \mathrm{OP1/OP4\ closure\ proven}=\mathrm{false}.
\]
\end{proposition}

\begin{proof}
The v4.4 ledger gives zero known-open finite blockers, so a closure attempt is
allowed.  But the v4.4 closure targets require promotion from local finite
certificates to OP1 and OP4 closure statements inside the declared grammar.
Those promotions are precisely the open proof obligations \(P1\)-\(P5\).
Hence no closure is proved at intake.
\end{proof}

\begin{theorem}[v4.4 no-overclaim guard]
The v4.4 line may attempt OP1 and OP4 closure only by discharging \(P1\)-\(P5\).
Until that happens, every statement remains a closure-attempt statement, not a
closure theorem.
\end{theorem}


% --- v4.4 P1 OP1 Selector-Domain Promotion Insert ---

\section{P1: OP1 Selector-Domain Promotion}
\label{sec:v45-p1-op1-selector-domain-promotion}

This section addresses proof obligation \(P1\) from the v4.4 closure-attempt
ledger.  It promotes the v4.4 OP1 local certificates to the OP1
selector-domain side.  It discharges \(P1\) only.  It does not close OP1 or
OP4, because \(P2\)-\(P5\) remain open.

\begin{definition}[OP1 selector-domain residual]
Let the OP1 selector-domain residual be
\[
    R_{\mathrm{OP1}}^{\mathrm{sel}}
    :=
    R_{\mathrm{OP1,bounded}}^{\mathrm{uncert}}
    \cup
    R_{\mathrm{OP1,blocking}}^{\mathrm{rem}}
    \cup
    R_{\mathrm{scale}}
    \cup
    R_{\mathrm{lens}}
    \cup
    R_{\mathrm{chart}}
    \cup
    B_{\mathrm{HC/MM,new}}.
\]
It is a selector-domain residual only; it is not the full coupled OP1/OP4
closure residual.
\end{definition}

\begin{definition}[P1 discharge predicate]
The proof obligation \(P1\) is discharged when
\[
    R_{\mathrm{OP1,bounded}}^{\mathrm{uncert}}=0,\qquad
    R_{\mathrm{OP1,blocking}}^{\mathrm{rem}}=0,
\]
and no scale, lens, chart, or new HC/MM debt is imported into the OP1
selector-domain side.
\end{definition}

\begin{proposition}[OP1 selector-domain evidence ledger]
The v4.4-to-v4.4 OP1 evidence ledger gives
\[
\begin{array}{c|r|r|r}
\text{OP1 component} & N_{\mathrm{in}} & N_{\mathrm{cert/res}} &
N_{\mathrm{rem}}\\
\hline
\mathrm{bounded\text{-}open} & 196 & 196 & 0\\
\mathrm{blocking} & 195 & 195 & 0
\end{array}
\]
and the imported debts are
\[
    R_{\mathrm{scale}}=0,\quad
    R_{\mathrm{lens}}=0,\quad
    R_{\mathrm{chart}}=0,\quad
    B_{\mathrm{HC/MM,new}}=0.
\]
Consequently
\[
    |R_{\mathrm{OP1}}^{\mathrm{sel}}|=0.
\]
\end{proposition}

\begin{proof}
The bounded-open OP1 points are locally certified by the coupled harmlessness
certificate, leaving \(0\) uncertified bounded-open points.  The
blocking OP1 points are resolved by the dense-carrier refinement, leaving
\(0\) OP1 blocking points.  The guard ledger imports no scale,
lens, chart, or new HC/MM debt.  Hence the OP1 selector-domain residual has
the stated size.
\end{proof}

\begin{theorem}[P1 discharge]
In the present v4.4 ledger,
\[
    R_{\mathrm{OP1}}^{\mathrm{sel}}=\varnothing.
\]
Therefore the OP1 selector-domain promotion obligation \(P1\) is
\[
    \mathrm{P1}=\mathrm{discharged}.
\]
This is not an OP1 closure theorem and not an OP1/OP4 closure theorem.  The
remaining proof obligations are
\[
    P2,\quad P3,\quad P4,\quad P5.
\]
\end{theorem}

\begin{remark}[Hopf-linking guard inherited by P1]
The dense OP1 carrier-refinement step sampled
\[
    320/320
\]
stable Hopf-linking pairs, with maximum linking deviation
\[
    \Delta_{L,\mathrm{OP1}}^{\max}\approx 0.0044.
\]
This is a finite guard check for the OP1 selector-domain promotion, not a
global Hopf theorem.
\end{remark}


% --- v4.4 P2 OP4 Post-Residual Promotion Insert ---

\section{P2: OP4 Post-Residual Promotion}
\label{sec:v45-p2-op4-post-residual-promotion}

This section addresses proof obligation \(P2\) from the v4.4 closure-attempt
ledger.  It promotes the OP4 local harmlessness certificates to the OP4
post-residual side.  It discharges \(P2\) only.  It does not close OP4 or the
coupled OP1/OP4 problem, because \(P3\)-\(P5\) remain open.

\begin{definition}[OP4 post-residual ledger]
Let the OP4 post-residual ledger be
\[
    R_{\mathrm{OP4}}^{\mathrm{post}}
    :=
    R_{\mathrm{OP4,uncert}}^{\mathrm{first}}
    \cup
    R_{\mathrm{OP4,open}}^{\mathrm{first}}
    \cup
    R_{\mathrm{scale}}
    \cup
    R_{\mathrm{lens}}
    \cup
    R_{\mathrm{chart}}
    \cup
    B_{\mathrm{HC/MM,new}},
\]
after applying the declared OP4 refinement certificates.  This is the OP4
post-residual side only; it is not yet the coupled selector closure residual.
\end{definition}

\begin{definition}[P2 discharge predicate]
The proof obligation \(P2\) is discharged when all first-stage uncertified
OP4 candidates and all first-stage OP4 open carrier-density residuals are
absorbed by refined harmlessness certificates, and when no scale, lens, chart,
or new HC/MM debt is imported into the OP4 post-residual side.
\end{definition}

\begin{proposition}[OP4 post-residual evidence ledger]
The OP4 evidence ledger gives
\[
\begin{array}{c|r|r|r}
\text{OP4 component} & N_{\mathrm{in}} & N_{\mathrm{cert/res}} &
N_{\mathrm{rem}}\\
\hline
\mathrm{first\ harmless\ candidates} & 313 & 251 & 62\\
\mathrm{open/uncertified\ refinement} & 140 & 140 & 0
\end{array}
\]
and the imported debts are
\[
    R_{\mathrm{scale}}=0,\quad
    R_{\mathrm{lens}}=0,\quad
    R_{\mathrm{chart}}=0,\quad
    B_{\mathrm{HC/MM,new}}=0.
\]
Consequently
\[
    |R_{\mathrm{OP4}}^{\mathrm{post}}|=0.
\]
\end{proposition}

\begin{proof}
The first OP4 certificate test locally certifies \(251\) of the
first-stage candidates and forwards \(62\) uncertified
candidates together with \(78\) open carrier-density residuals to
the refinement test.  The refinement test certifies \(140\) points
under the both-chart and non-inflation predicates, leaving \(0\)
OP4 open residuals.  The guard ledger imports no scale, lens, chart, or new
HC/MM debt.  Hence the OP4 post-residual ledger has the stated size.
\end{proof}

\begin{theorem}[P2 discharge]
In the present v4.4 ledger,
\[
    R_{\mathrm{OP4}}^{\mathrm{post}}=\varnothing.
\]
Therefore the OP4 post-residual promotion obligation \(P2\) is
\[
    \mathrm{P2}=\mathrm{discharged}.
\]
This is not an OP4 closure theorem and not an OP1/OP4 closure theorem.  The
remaining proof obligations are
\[
    P3,\quad P4,\quad P5.
\]
\end{theorem}

\begin{remark}[Hopf-linking guard inherited by P2]
The OP4 unresolved-refinement step sampled
\[
    320/320
\]
stable Hopf-linking pairs, with maximum linking deviation
\[
    \Delta_{L,\mathrm{OP4}}^{\max}\approx 0.0144.
\]
This is a finite guard check for the OP4 post-residual promotion, not a global
Hopf theorem.
\end{remark}


% --- v4.4 P3 Coupled Selector Compatibility Insert ---

\section{P3: Coupled Selector Compatibility}
\label{sec:v45-p3-coupled-selector-compatibility}

This section addresses proof obligation \(P3\).  Obligations \(P1\) and \(P2\)
have promoted the OP1 selector-domain side and the OP4 post-residual side,
respectively.  \(P3\) checks that composing these two sides through the
declared coupled selector does not create a new coupled residual.  This
discharges \(P3\) only.  It does not close OP1/OP4, because \(P4\) and \(P5\)
remain open.

\begin{definition}[Coupled selector residual]
Let
\[
    C_{1/4}^{\mathrm{sel}}:
    P(\rho_{1/4}^{\mathrm{pre}})
    \times
    P(\rho_{1/4}^{\mathrm{shift}})
    \longrightarrow
    E_{1/4}^{\mathrm{sel}}
\]
be the declared coupled selector interface.  Define the coupled selector
residual by
\[
    R_{1/4}^{\mathrm{sel}}
    :=
    R_{\mathrm{OP1}}^{\mathrm{sel}}
    \cup
    R_{\mathrm{OP4}}^{\mathrm{post}}
    \cup
    R_{\mathrm{cross}}^{\mathrm{sel}},
\]
where \(R_{\mathrm{cross}}^{\mathrm{sel}}\) is the residual created, if any, by
composing the two promoted sides through \(C_{1/4}^{\mathrm{sel}}\).
\end{definition}

\begin{definition}[P3 discharge predicate]
The proof obligation \(P3\) is discharged when
\[
    R_{\mathrm{OP1}}^{\mathrm{sel}}=\varnothing,\qquad
    R_{\mathrm{OP4}}^{\mathrm{post}}=\varnothing,\qquad
    R_{\mathrm{cross}}^{\mathrm{sel}}=\varnothing.
\]
This is a compatibility statement for the declared selector interface.  It is
not yet the no-new-debt promotion and not the finite-to-declared-grammar
closure guard.
\end{definition}

\begin{proposition}[Coupled compatibility evidence ledger]
The coupled-selector evidence ledger gives
\[
\begin{array}{c|r|c}
\text{component} & \text{residual size} & \text{status}\\
\hline
R_{\mathrm{OP1}}^{\mathrm{sel}} & 0 & \mathrm{P1\ discharged}\\
R_{\mathrm{OP4}}^{\mathrm{post}} & 0 & \mathrm{P2\ discharged}\\
R_{\mathrm{known\ open}}^{1/4} & 0 & \mathrm{v4.4\ handoff}\\
R_{\mathrm{cross}}^{\mathrm{sel}} & 0 & \mathrm{composition\ audit}
\end{array}
\]
Consequently
\[
    |R_{1/4}^{\mathrm{sel}}|=0.
\]
\end{proposition}

\begin{proof}
By \(P1\), the OP1 selector-domain residual is empty.  By \(P2\), the OP4
post-residual ledger is empty.  The v4.4 coupled handoff has zero known-open
remainder.  Since the coupled comparison is mediated by the declared selector
interface rather than by a raw same-variance contraction, the composition audit
adds no cross-leakage in the finite ledger.  Therefore the coupled selector
residual is empty.
\end{proof}

\begin{theorem}[P3 discharge]
In the present v4.4 ledger,
\[
    R_{1/4}^{\mathrm{sel}}=\varnothing.
\]
Therefore the coupled selector compatibility obligation \(P3\) is
\[
    \mathrm{P3}=\mathrm{discharged}.
\]
This is not an OP1/OP4 closure theorem.  The remaining proof obligations are
\[
    P4,\quad P5.
\]
\end{theorem}


% --- v4.4 P4 No-New-Debt Promotion Insert ---

\section{P4: No-New-Debt Promotion}
\label{sec:v45-p4-no-new-debt-promotion}

This section addresses proof obligation \(P4\).  Obligations \(P1\), \(P2\),
and \(P3\) have promoted the OP1 selector-domain side, the OP4 post-residual
side, and the coupled selector compatibility ledger.  \(P4\) checks that this
promotion has not introduced scale debt, lens debt, chart debt, or new HC/MM
debt.  It discharges \(P4\) only.  It does not close OP1/OP4, because \(P5\)
remains open.

\begin{definition}[Promoted debt ledger]
The promoted no-new-debt ledger is
\[
    B_{\mathrm{new}}^{1/4}
    :=
    R_{\mathrm{scale}}
    \cup
    R_{\mathrm{lens}}
    \cup
    R_{\mathrm{chart}}
    \cup
    B_{\mathrm{HC/MM,new}}
    \cup
    B_{\mathrm{comp}}^{P1-P3},
\]
where \(B_{\mathrm{comp}}^{P1-P3}\) is any new debt introduced by composing the
\(P1\), \(P2\), and \(P3\) promotions.
\end{definition}

\begin{definition}[P4 discharge predicate]
The proof obligation \(P4\) is discharged when
\[
    R_{\mathrm{scale}}=0,\qquad
    R_{\mathrm{lens}}=0,\qquad
    R_{\mathrm{chart}}=0,\qquad
    B_{\mathrm{HC/MM,new}}=0,\qquad
    B_{\mathrm{comp}}^{P1-P3}=0.
\]
This is a no-new-debt promotion statement, not the final
finite-to-declared-grammar closure guard.
\end{definition}

\begin{proposition}[No-new-debt evidence ledger]
The promoted debt ledger gives
\[
\begin{array}{c|r}
\text{debt channel} & \text{promoted debt}\\
\hline
R_{\mathrm{scale}} & 0\\
R_{\mathrm{lens}} & 0\\
R_{\mathrm{chart}} & 0\\
B_{\mathrm{HC/MM,new}} & 0\\
B_{\mathrm{comp}}^{P1-P3} & 0
\end{array}
\]
and therefore
\[
    |B_{\mathrm{new}}^{1/4}|=0.
\]
\end{proposition}

\begin{proof}
The v4.4 guard ledger records zero scale, lens, chart, and new HC/MM debt.  The
\(P3\) coupled selector compatibility ledger records zero coupled selector
residual, so the composition of \(P1\)-\(P3\) contributes no additional debt in
the promoted finite ledger.  Hence the promoted debt total is zero.
\end{proof}

\begin{theorem}[P4 discharge]
In the present v4.4 ledger,
\[
    B_{\mathrm{new}}^{1/4}=\varnothing.
\]
Therefore the no-new-debt promotion obligation \(P4\) is
\[
    \mathrm{P4}=\mathrm{discharged}.
\]
This is not an OP1/OP4 closure theorem.  The remaining proof obligation is
\[
    P5.
\]
\end{theorem}


% --- v4.4 P5 Finite-to-Declared-Grammar Closure Guard Insert ---

\section{P5: Finite-to-Declared-Grammar Closure Guard}
\label{sec:v45-p5-finite-to-declared-grammar-closure-guard}

This section addresses proof obligation \(P5\), the final guard in the v4.4
OP1/OP4 closure attempt.  The first four obligations have discharged the OP1
selector-domain side, the OP4 post-residual side, the coupled selector
compatibility ledger, and the no-new-debt ledger.  \(P5\) checks that these
finite certificates cover exactly the declared local orbit--Hopf coupled
selector grammar and do not overclaim a global physical or unrestricted
context-free OP1/OP4 theorem.

\begin{definition}[Declared local OP1/OP4 grammar]
The declared local OP1/OP4 grammar is the finite guarded grammar generated by
\[
    \mathcal G_{1/4}^{\mathrm{loc}}
    :=
    \bigl(
    \mathcal D_{\mathrm{OP1}},
    \mathcal D_{\mathrm{OP4}},
    P(\rho_{1/4}^{\mathrm{pre}}),
    P(\rho_{1/4}^{\mathrm{shift}}),
    C_{1/4}^{\mathrm{sel}},
    E_{1/4}^{\mathrm{sel}}
    \bigr),
\]
together with the v4.4/v4.4 carrier, lens, scale, chart, and debt ledgers.
A statement is inside the declared grammar only when it is mediated by these
objects and respects their guards.
\end{definition}

\begin{definition}[Finite-to-declared-grammar coverage residual]
The coverage residual is
\[
    R_{\mathrm{cov}}^{1/4}
    :=
    R_{\mathrm{known\ open}}^{1/4}
    \cup
    R_{1/4}^{\mathrm{sel}}
    \cup
    B_{\mathrm{new}}^{1/4}
    \cup
    B_{\mathrm{comp}}^{P1-P3}.
\]
The finite-to-declared-grammar guard is passed when
\[
    R_{\mathrm{cov}}^{1/4}=\varnothing
\]
and the claim is explicitly restricted to \(\mathcal G_{1/4}^{\mathrm{loc}}\).
\end{definition}

\begin{definition}[No-overclaim boundary]
The following are not claimed by this theorem:
\[
\begin{array}{ll}
(1) & \text{a global physical closure theorem for QFT, GR, or all readouts;}\\
(2) & \text{an unrestricted context-free OP1/OP4 theorem outside }
      \mathcal G_{1/4}^{\mathrm{loc}};\\
(3) & \text{a global Hopf theorem beyond the finite audited carrier families;}\\
(4) & \text{a release-final claim or DOI-bearing public release.}
\end{array}
\]
\end{definition}

\begin{proposition}[P5 guard ledger]
The finite-to-declared-grammar guard ledger gives
\[
\begin{array}{c|r}
\text{component} & \text{residual/debt}\\
\hline
R_{\mathrm{known\ open}}^{1/4} & 0\\
R_{1/4}^{\mathrm{sel}} & 0\\
B_{\mathrm{new}}^{1/4} & 0\\
B_{\mathrm{comp}}^{P1-P3} & 0\\
R_{\mathrm{cov}}^{1/4} & 0
\end{array}
\]
and the scope guard returns
\[
    \mathrm{declared\text{-}grammar\ only}=\mathrm{true}.
\]
\end{proposition}

\begin{proof}
By the v4.4/v4.4 ledger, the known-open coupled remainder is zero.  By \(P3\),
the coupled selector residual is zero.  By \(P4\), the promoted no-new-debt
ledger is zero.  The composition debt introduced by \(P1\)-\(P3\) is also zero.
Therefore the coverage residual is empty.  The no-overclaim boundary restricts
the resulting theorem to the declared local grammar, so the finite
certificates are promoted only within \(\mathcal G_{1/4}^{\mathrm{loc}}\).
\end{proof}

\begin{theorem}[Declared-grammar OP1/OP4 closure theorem]
Inside the declared local orbit--Hopf coupled selector grammar
\(\mathcal G_{1/4}^{\mathrm{loc}}\),
\[
    R_{\mathrm{OP1}}^{\mathrm{sel}}=\varnothing,\qquad
    R_{\mathrm{OP4}}^{\mathrm{post}}=\varnothing,\qquad
    R_{1/4}^{\mathrm{sel}}=\varnothing,\qquad
    B_{\mathrm{new}}^{1/4}=\varnothing,
\]
and
\[
    R_{\mathrm{cov}}^{1/4}=\varnothing.
\]
Thus all v4.4 proof obligations \(P1\)-\(P5\) are discharged and the OP1/OP4
closure attempt succeeds inside the declared local grammar:
\[
    \mathrm{Closure}_{\mathcal G_{1/4}^{\mathrm{loc}}}(\mathrm{OP1},\mathrm{OP4}).
\]
This is a declared-grammar closure theorem only.  It is not a global physical
closure theorem and not an unrestricted OP1/OP4 theorem outside the declared
grammar.
\end{theorem}


% --- v4.4 Closure Summary / Release-Candidate Audit Insert ---


\section{v4.4 Closure Summary and Release-Candidate Audit}
\label{sec:v45-closure-summary-release-candidate-audit}

This section consolidates the v4.4 OP1/OP4 closure attempt after \(P1\)-\(P5\).
It is a release-candidate audit section, not a release-final section.  No DOI
is assigned here, no Zenodo description is generated here, and no global
physical or unrestricted context-free OP1/OP4 closure is claimed.

\begin{definition}[v4.4 declared-grammar closure scope]
The v4.4 closure theorem is scoped to the declared local orbit--Hopf
coupled-selector grammar
\[
    \mathcal G_{1/4}^{\mathrm{loc}}.
\]
It uses only the declared OP1/OP4 ingress objects, the coupled selector
\(C_{1/4}^{\mathrm{sel}}\), and the finite carrier, lens, scale, chart, and
debt ledgers audited in v4.4/v4.4.
\end{definition}

\begin{proposition}[Proof-obligation summary]
The v4.4 proof-obligation ledger after \(P5\) is
\[
\begin{array}{c|c}
\text{obligation} & \text{status}\\
\hline
P1\ \text{OP1 selector-domain promotion} & \mathrm{discharged}\\
P2\ \text{OP4 post-residual promotion} & \mathrm{discharged}\\
P3\ \text{coupled selector compatibility} & \mathrm{discharged}\\
P4\ \text{no-new-debt promotion} & \mathrm{discharged}\\
P5\ \text{finite-to-declared-grammar closure guard} & \mathrm{discharged}
\end{array}
\]
Thus
\[
    N_{\mathrm{discharged}}=5,\qquad
    N_{\mathrm{open}}=0,\qquad
    N_{\mathrm{blocked}}=0.
\]
This release-candidate audit is not itself a release.
\end{proposition}

\begin{proposition}[Residual and debt summary]
The final v4.4 local closure ledger records
\[
\begin{array}{c|r}
\text{ledger component} & \text{residual/debt}\\
\hline
R_{\mathrm{known\ open}}^{1/4} & 0\\
R_{1/4}^{\mathrm{sel}} & 0\\
B_{\mathrm{new}}^{1/4} & 0\\
B_{\mathrm{comp}}^{P1-P3} & 0\\
R_{\mathrm{cov}}^{1/4} & 0
\end{array}
\]
All entries vanish in the declared local grammar.
\end{proposition}

\begin{theorem}[v4.4 declared-grammar OP1/OP4 closure summary]
Inside \(\mathcal G_{1/4}^{\mathrm{loc}}\), the v4.4 closure attempt has
discharged \(P1\)-\(P5\), and the final coverage residual is empty:
\[
    R_{\mathrm{cov}}^{1/4}=\varnothing.
\]
Therefore the declared-grammar OP1/OP4 closure theorem is established:
\[
    \mathrm{Closure}_{\mathcal G_{1/4}^{\mathrm{loc}}}(\mathrm{OP1},\mathrm{OP4}).
\]
This statement is restricted to the declared local grammar.
\end{theorem}

\begin{remark}[No-overclaim boundary]
The v4.4 declared-grammar closure summary does not assert any of the following:
\[
\begin{array}{ll}
(1) & \text{global physical closure for QFT, GR, all readouts, or all models;}\\
(2) & \text{unrestricted context-free OP1/OP4 closure outside }
      \mathcal G_{1/4}^{\mathrm{loc}};\\
(3) & \text{a global Hopf theorem beyond the audited finite carrier families;}\\
(4) & \text{release-final status, a DOI assignment, or Zenodo publication.}
\end{array}
\]
\end{remark}

\begin{proposition}[Release-candidate audit status]
The release-candidate audit status is
\[
    \mathrm{RC\ audit\ candidate}=true.
\]
This means the source is ready for a later release decision.  It is not itself
a release.  A public release still requires a separate release-final audit,
frontmatter pass, DOI metadata pass, and upload checklist.
\end{proposition}


% --- v4.4 Release-Final Preflight / Metadata Intake Insert ---

\section{v4.4 Release-Final Preflight and Metadata Intake Gate}
\label{sec:v45-release-final-preflight-metadata-intake}

This section records the release-final preflight state after the v4.4
release-candidate audit.  It is not a release-final section and assigns no
DOI.  It exists to make the remaining release blockers explicit.

\begin{definition}[Release-final preflight gate]
The release-final preflight gate separates mathematical release-candidate
readiness from public release readiness.  The v4.4 declared-grammar closure
may be release-candidate audited while the public release remains blocked by
metadata, frontmatter, and upload requirements.
\end{definition}

\begin{proposition}[Mathematical RC status]
The mathematical release-candidate state is
\[
\begin{array}{c|c}
\text{gate} & \text{status}\\
\hline
\mathrm{RC\ audit\ pass} & true\\
\mathrm{declared\ grammar\ closure} & true\\
\mathrm{global/unrestricted\ closure} & false\\
R_{\mathrm{cov}}^{1/4} & 0\\
\mathrm{open\ proof\ obligations} & 0
\end{array}
\]
Thus the mathematical RC side is ready for release metadata intake.
\end{proposition}

\begin{proposition}[Release blockers]
The release-final gate remains blocked by the following public-release tasks:
\[
\begin{array}{c|c}
\text{release task} & \text{status}\\
\hline
\mathrm{version\ DOI} & \mathrm{missing}\\
\mathrm{release\text{-}final\ frontmatter\ pass} & \mathrm{not\ run}\\
\mathrm{Zenodo\ description} & \mathrm{not\ generated}\\
\mathrm{public\ PDF\ rebuild} & \mathrm{not\ run}\\
\mathrm{upload\ checklist} & \mathrm{not\ run}
\end{array}
\]
Consequently
\[
    \mathrm{release\text{-}final\ allowed}=false.
\]
\end{proposition}

\begin{theorem}[Preflight handoff]
The v4.4 source is mathematically ready for release metadata intake, but not
ready for release-final publication.  A release-final step requires a DOI or an
explicit no-DOI decision, a release-final frontmatter pass, a public PDF
rebuild, a release audit, a Zenodo description prepared at release time, and an
upload checklist.
\end{theorem}

\begin{remark}[No release claim]
This preflight section does not assign a DOI, does not generate a Zenodo
description, and does not mark the source as release-final.
\end{remark}


% --- v4.4 Clean RC Consolidation / Release-Readiness Audit Insert ---

\section{v4.4 Clean RC Consolidation and Release-Readiness Audit}
\label{sec:v45-clean-rc-consolidation-release-readiness-audit}

This section records the post-reference-hygiene release-candidate consolidation.
It is not a release-final section, assigns no DOI, and generates no Zenodo
description.  It only records that the current v4.4 source line has been
rebuilt without unresolved reference markers and that the declared-grammar
closure remains scoped.

\begin{definition}[Clean RC consolidation state]
A clean RC consolidation state is a release-candidate working state in which
the declared-grammar closure summary, the no-overclaim guards, and the
reference-hygiene rebuild are all simultaneously visible in the source-native
Active Core.
\end{definition}

\begin{proposition}[Reference-hygiene status]
The clean v4.4 reference-hygiene rebuild records
\[
\begin{array}{c|r}
\text{audit item} & \text{count}\\
\hline
\text{visible double-question reference markers} & 0\\
\text{visible citation question markers} & 0\\
\text{undefined reference warnings} & 0\\
\text{undefined citation warnings} & 0\\
\text{fatal/LaTeX errors} & 0
\end{array}
\]
Thus the reference-hygiene side is clean for the current working RC line.
\end{proposition}

\begin{proposition}[Closure-scope status]
The v4.4 closure summary remains scoped as follows:
\[
\begin{array}{c|c}
\text{claim} & \text{status}\\
\hline
\text{declared local grammar OP1/OP4 closure} & true\\
\text{global/unrestricted OP1/OP4 closure} & false\\
R_{\mathrm{cov}}^{1/4} & 0
\end{array}
\]
The closure claim remains restricted to the declared local orbit--Hopf
coupled-selector grammar.
\end{proposition}

\begin{proposition}[Release-readiness gates]
The clean RC line currently has
\[
    N_{\mathrm{passed}}=3,
    \qquad
    N_{\mathrm{blocked}}=5
\]
release-readiness gates.  Since at least the version DOI, release-final
frontmatter pass, Zenodo description, public release rebuild, and upload
checklist are not yet complete,
\[
    \mathrm{release\text{-}final\ ready}=false.
\]
\end{proposition}

\begin{theorem}[Clean RC handoff]
The current v4.4 Active Core is a clean release-candidate working line with no
visible unresolved reference markers and with the declared-grammar closure
scope intact.  It is ready for release metadata intake, but it is not
release-final.  A public release still requires a version DOI or explicit
no-DOI decision, a release-final frontmatter pass, a public release rebuild, a
Zenodo description prepared at release time, and an upload checklist.
\end{theorem}

\begin{remark}[Build discipline]
Future builds of this large source must use repeated LaTeX passes, preferably
three passes, before a PDF is audited for references.  A single pass can
reintroduce visible unresolved-reference markers even when the source labels
exist.
\end{remark}


\clearpage
\section{v4.4 Release-Final DOI and Public Upload Audit}
\label{sec:v44-release-final-doi-public-upload-audit}

This section records the public release-final metadata pass for the Zenodo-facing v4.4 release.  It is metadata and audit material only; it does not introduce a new mathematical object beyond the declared local closure theorem already recorded in the Active Core.

\begin{definition}[v4.4 public release-final package]
The v4.4 public release-final package is the package whose public version line is v4.4 and whose version-specific Zenodo DOI is
\[
    \href{https://doi.org/10.5281/zenodo.20068747}{10.5281/zenodo.20068747}.
\]
The Companion concept DOI remains
\[
    \href{https://doi.org/10.5281/zenodo.19210728}{10.5281/zenodo.19210728}.
\]
All DOI records are release metadata only.  They are not theorem labels, route tags, tensor indices, physical regimes, or proof data.
\end{definition}

\begin{proposition}[Public version alignment]
The public release version, title-page version, public filenames, manifest, checksums, and Zenodo version field are aligned to v4.4.  Internal v4.5 working labels from the proof-attempt chain are not public release version labels in this release.
\end{proposition}

\begin{theorem}[Release-final conservativity]
The insertion of the v4.4 version DOI, release notes, manifest, checksums, Zenodo description, and upload checklist is conservative over the declared local grammar closure statement.  It changes release metadata only and does not promote the scoped theorem to a global physical theorem or to unrestricted OP1/OP4 closure outside \(\mathcal G_{1/4}^{\mathrm{loc}}\).
\end{theorem}

\begin{remark}[No-overclaim boundary]
This release asserts the declared local orbit--Hopf coupled-selector grammar closure only.  It does not assert global/unrestricted OP1/OP4 closure, a global Hopf theorem, global QFT/GR closure, or Standard-Model-spectrum derivation.
\end{remark}


% --- v4.5.0-r01 OP5 Local-to-Global Promotion Intake ---

\section{OP5 Local-to-Global Promotion Intake Layer}
\label{sec:v450-r01-op5-local-to-global-promotion-intake}

This section starts the v4.5.0 working line.  The public release basis remains
v4.4.  The present layer is a working intake layer only: it assigns no DOI,
does not mark a release-final state, and does not claim global OP5 closure.

\begin{definition}[OP5 local carrier-closure import]
The v4.5.0 line imports OP5 only in its local carrier/invariant form.  We write
\[
    \mathrm{Closure}_{\mathrm{loc}}(\mathrm{OP5})
\]
for the already available carrier-side statement that OP5 is closed inside its
admissible local return grammar.  This import is not an unrestricted global
OP5 theorem.
\end{definition}

\begin{definition}[OP5 promotion target]
The OP5 local-to-global promotion target is the problem of deciding whether the
local carrier/invariant OP5 closure can be promoted through the HC/MM-chart
route-to-global interface without introducing new coupled debt.  The promotion
target is represented schematically by
\[
    \mathrm{Closure}_{\mathrm{loc}}(\mathrm{OP5})
    \xrightarrow{\ \Pi_{\mathrm{HC/MM\leftrightarrow chart}}\ }
    \mathrm{Closure}_{\mathrm{prom}}(\mathrm{OP5}).
\]
The arrow is not assumed to exist.  It is the object of the v4.5.0 witness
search.
\end{definition}

\begin{definition}[OP5 promotion residual]
Let
\[
    R_{\mathrm{OP5}}^{\mathrm{prom}}
\]
be the residual measuring the gap between local OP5 carrier-closure and an
admissible promoted OP5 closure in the larger HC/MM-chart grammar.  At intake,
the residual is initialized by the named persistent route-to-global obstruction
\[
    \left|\mathfrak D_{\mathrm{HC/MM}\leftrightarrow\mathrm{chart}}^{\mathrm{str}}\right|=1.
\]
Thus
\[
    |R_{\mathrm{OP5}}^{\mathrm{prom}}|=1
\]
at the start of v4.5.0-r01.
\end{definition}

\begin{definition}[HC/MM witness menu for OP5 promotion]
The first witness family for the promotion attempt is
\[
    W_{\mathrm{HC/MM}}^{\mathrm{OP5}}
    =
    \left(
    W_{\mathrm{route}}^{\mathrm{HC/MM}},
    W_{\mathrm{chart}}^{\mathrm{bridge}},
    W_{\mathrm{carrier}}^{\mathrm{OP5}}
    \right).
\]
A successful witness must produce strict coupled-debt reduction
\[
    \Delta_{\mathrm{cpl}}<0
\]
without hidden leakage, residual reopening, or new chart debt.  If no such
witness exists, OP5 remains locally closed but globally obstructed relative to
the current HC/MM-chart grammar.
\end{definition}

\begin{proposition}[v4.5.0-r01 intake ledger]
The OP5 promotion-intake ledger is
\[
\begin{array}{c|c}
\text{item} & \text{status}\\
\hline
\text{public release basis} & \mathrm{v4.4}\\
\text{working line} & \mathrm{v4.5.0\mbox{-}r01}\\
\text{local OP5 carrier-closure} & \mathrm{imported}\\
\text{global OP5 closure} & \mathrm{not\ claimed}\\
\text{OP5 promotion residual} & |R_{\mathrm{OP5}}^{\mathrm{prom}}|=1\\
\text{witness search} & \mathrm{declared,\ not\ executed}
\end{array}
\]
\end{proposition}

\begin{theorem}[No-overclaim guard for OP5 promotion intake]
The v4.5.0-r01 layer may define the OP5 promotion residual and the HC/MM
witness menu.  It may not claim global OP5 closure, global OP1/OP4 closure,
global physical closure, or a release-final state.  The next admissible step is
a finite witness search for \( \Delta_{\mathrm{cpl}}<0 \).
\end{theorem}

\begin{remark}[Relation to v4.4]
The v4.4 release remains the public release basis.  The v4.5.0-r01 layer is a
post-release working branch that investigates whether the local declared-grammar
closures can be promoted beyond their current grammar without violating the
recorded route-to-global obstruction.
\end{remark}


% --- v4.5.0-r02 HC/MM Coupled-Debt Witness Search ---

\section{HC/MM Coupled-Debt Witness Search Layer}
\label{sec:v450-r02-hcmm-coupled-debt-witness-search}

This section executes the first finite witness search for the OP5
local-to-global promotion residual introduced in v4.5.0-r01. The search is
conservative: it only tests the declared first witness family and does not
claim global OP5 closure.

\begin{definition}[Coupled-debt delta]
For a witness \(W\), define
\[
    \Delta_{\mathrm{cpl}}(W)
    :=
    d_{\mathrm{after}}(W)
    + \ell_{\mathrm{hidden}}(W)
    + d_{\mathrm{chart,new}}(W)
    + r_{\mathrm{reopen}}(W)
    - d_{\mathrm{before}}(W).
\]
A strict reduction witness must satisfy
\[
    \Delta_{\mathrm{cpl}}(W)<0.
\]
A neutral or positive value does not promote OP5 beyond the current
HC/MM-chart grammar.
\end{definition}

\begin{definition}[First finite OP5 witness family]
The first finite witness family is
\[
    W_{\mathrm{OP5}}^{(1)}
    =
    \{W_0,W_1,W_2,W_3,W_4,W_5\}.
\]
It contains route-carry, route-balanced, chart-bridge, carrier-return,
absorption, and selector-handoff candidates. This is a first-family search,
not a complete search over all possible witnesses.
\end{definition}

\begin{proposition}[r02 witness register]
The first witness register gives
\[
\begin{array}{c|r|r|r|r|r}
\text{witness} & d_{\mathrm{before}} & d_{\mathrm{after}} &
\ell_{\mathrm{hidden}} & d_{\mathrm{chart,new}} &
\Delta_{\mathrm{cpl}}\\
\hline
W_0 & 1 & 1 & 0 & 0 & 0\\
W_1 & 1 & 1 & 0 & 0 & 0\\
W_2 & 1 & 1 & 0 & 1 & 1\\
W_3 & 1 & 1 & 0 & 0 & 1\\
W_4 & 1 & 0 & 1 & 0 & 0\\
W_5 & 1 & 1 & 0 & 0 & 0
\end{array}
\]
with the carrier-return reopening contribution included in \(W_3\).
Consequently,
\[
    \min_{W\in W_{\mathrm{OP5}}^{(1)}}
    \Delta_{\mathrm{cpl}}(W)
    =
    0.
\]
\end{proposition}

\begin{proof}
The register is evaluated by the coupled-debt delta formula. The identity,
route-balanced, and selector-handoff candidates are compatible but neutral.
The chart-bridge candidate introduces new chart debt. The carrier-return
candidate reopens a return-side residual. The absorption candidate removes the
visible debt but creates exactly one hidden-leakage unit, so its net coupled
delta is neutral. No tested witness has negative coupled delta.
\end{proof}

\begin{theorem}[First-family OP5 promotion obstruction]
No witness in \(W_{\mathrm{OP5}}^{(1)}\) satisfies
\[
    \Delta_{\mathrm{cpl}}<0.
\]
Therefore the first HC/MM coupled-debt witness search does not promote OP5
beyond its local carrier/invariant closure. The OP5 promotion residual remains
\[
    |R_{\mathrm{OP5}}^{\mathrm{prom}}|=1.
\]
This is a first-family obstruction result only. It does not prove global OP5
closure, and it does not rule out later enlarged witness families.
\end{theorem}


% --- v4.5.0-r03 OP5 Promotion Residual Classifier ---

\section{OP5 Promotion Residual Classifier}
\label{sec:v450-r03-op5-promotion-residual-classifier}

This section classifies the result of the first HC/MM coupled-debt witness
search from v4.5.0-r02.  It does not run a new witness search and does not
claim global OP5 closure.  Its purpose is to turn the first-family witness
register into a stable obstruction classification.

\begin{definition}[OP5 promotion classifier alphabet]
The r03 classifier uses the alphabet
\[
\begin{aligned}
    \mathcal A_{\mathrm{OP5}}^{\mathrm{prom}}
    =\{&\mathrm{red},\mathrm{route\_neutral},\mathrm{block\_chart},\\
       &\mathrm{block\_reopen},\mathrm{leak\_neutral},\\
       &\mathrm{persistent\_first\_family\_obstruction}\}.
\end{aligned}
\]
Here \(\mathrm{red}\) means strict reduction, \(\mathrm{route\_neutral}\) means
compatible but non-reducing carry, \(\mathrm{block\_chart}\) means new chart
debt, \(\mathrm{block\_reopen}\) means return-side reopening,
\(\mathrm{leak\_neutral}\) means apparent debt removal exactly offset by
hidden leakage, and \(\mathrm{persistent\_first\_family\_obstruction}\) means
that the first finite family contains no strict reducer while the promotion
residual remains.
\end{definition}

\begin{proposition}[r03 classifier register]
The r02 witness family is classified as
\[
\begin{array}{c|r}
\text{class} & \text{count}\\
\hline
\mathrm{red} & 0\\
\mathrm{route\_neutral} & 3\\
\mathrm{block\_chart} & 1\\
\mathrm{block\_reopen} & 1\\
\mathrm{leak\_neutral} & 1
\end{array}
\]
and the minimum coupled-debt change remains
\[
    \min \Delta_{\mathrm{cpl}} = 0.
\]
\end{proposition}

\begin{proof}
The classifier is obtained by applying the r03 alphabet to the r02 witness
register.  No row has \(\Delta_{\mathrm{cpl}}<0\), so no row is a strict
reduction witness.  Neutral route candidates preserve compatibility but do not
reduce debt.  The chart-bridge and carrier-return trials are blocked by new
chart debt and return-side reopening.  The absorption trial is neutral because
its visible debt removal is offset by hidden leakage.
\end{proof}

\begin{theorem}[First-family persistent OP5 promotion obstruction]
For the first finite witness family,
\[
    |R_{\mathrm{OP5}}^{\mathrm{prom}}| = 1
\]
after r02 and r03.  Therefore the r03 classifier assigns
\[
    R_{\mathrm{OP5}}^{\mathrm{prom}}
    \in
    \mathrm{persistent\_first\_family\_obstruction}.
\]
Local OP5 carrier/invariant closure is preserved, but global OP5 promotion is
not allowed by the first witness family.  This theorem is scoped to the first
finite family only; it does not rule out later enlarged witness families.
\end{theorem}

\begin{remark}[No-overclaim guard]
The r03 result should be read as: OP5 remains locally closed and first-family
globally obstructed.  It is not a global OP5 closure theorem and not a theorem
that all possible HC/MM witness families fail.
\end{remark}


% --- v4.5.0-r04 Enlarged Witness-Family Decision ---

\section{Enlarged Witness-Family Decision}
\label{sec:v450-r04-enlarged-witness-family-decision}

This section decides what to do after the r03 classifier assigns the OP5
promotion residual to persistent first-family obstruction.  It does not execute
a new search and does not claim global OP5 closure.

\begin{definition}[Sharpness versus enlargement fork]
After a first-family obstruction result, there are two admissible routes:
\[
\begin{array}{ll}
\mathrm{freeze} & \text{if first-family sharpness is proved;}\\
\mathrm{extend} & \text{if first-family sharpness is not proved.}
\end{array}
\]
The r03 classifier proved only first-family obstruction.  It did not prove
that all possible enlarged witness families fail.
\end{definition}

\begin{proposition}[r04 decision ledger]
The r04 decision ledger gives
\[
\begin{array}{c|c}
\text{gate} & \text{value}\\
\hline
\text{persistent first-family obstruction} & true\\
\text{strict first-family reducer exists} & false\\
\text{first-family sharpness proved} & false\\
\text{local OP5 closure preserved} & true\\
\text{global OP5 closure claimed} & false
\end{array}
\]
Therefore the admissible r04 decision is
\[
    \mathrm{decision}_{r04}=\mathrm{extend\_witness\_family}.
\]
\end{proposition}

\begin{proof}
The r02 and r03 layers show that the first finite witness family contains no
strict coupled-debt reducer and leaves
\[
    |R_{\mathrm{OP5}}^{\mathrm{prom}}|=1.
\]
However, r03 did not prove first-family sharpness as a theorem over all
enlarged witness families.  Since local OP5 carrier closure is preserved and
global OP5 closure is not claimed, the conservative route is to enlarge the
witness menu rather than freeze the obstruction as absolute.
\end{proof}

\begin{definition}[Second OP5 witness family]
The enlarged second witness family is
\[
    W_{\mathrm{OP5}}^{(2)}
    =
    \{U_0,U_1,U_2,U_3,U_4,U_5\}.
\]
Its channels are:
\[
\begin{array}{c|l}
U_0 & \text{two-sided tensor bridge}\\
U_1 & \text{chart-tilt reroute}\\
U_2 & \text{carrier-return damping}\\
U_3 & \text{explicit leakage account}\\
U_4 & \text{Hopf return-phase probe}\\
U_5 & \text{composite route-chart-carrier witness}
\end{array}
\]
The family is declared here but not executed here.
\end{definition}

\begin{theorem}[r04 enlarged-family handoff]
The v4.5.0-r04 layer selects the extension route:
\[
    W_{\mathrm{OP5}}^{(1)}
    \longrightarrow
    W_{\mathrm{OP5}}^{(2)}.
\]
The OP5 status after r04 is: locally closed, first-family obstructed, and
second-family declared but not executed.  No global OP5 closure is claimed.
\end{theorem}


% --- v4.5.0-r05 Execute W_OP5^(2) Enlarged Witness Search ---

\section{Enlarged OP5 Witness-Family Execution}
\label{sec:v450-r05-execute-enlarged-op5-witness-family}

This section executes the enlarged witness family \(W_{\mathrm{OP5}}^{(2)}\)
declared in v4.5.0-r04.  The search remains conservative: a strict reduction
counts only if the witness is admissible.  In particular, a negative formal
\(\Delta_{\mathrm{cpl}}\) is not enough if the comparison uses an unlicensed
same-variance contraction.

\begin{definition}[Admissible strict OP5 promotion witness]
A witness \(U\in W_{\mathrm{OP5}}^{(2)}\) is an admissible strict OP5
promotion witness when
\[
    \Delta_{\mathrm{cpl}}(U)<0,
\]
and simultaneously
\[
    \ell_{\mathrm{hidden}}(U)=0,\qquad
    d_{\mathrm{chart,new}}(U)=0,\qquad
    r_{\mathrm{reopen}}(U)=0,\qquad
    s_{\mathrm{samevar}}^{\mathrm{unlic}}(U)=0.
\]
Thus hidden leakage, new chart debt, residual reopening, and unlicensed
same-variance comparison each block promotion.
\end{definition}

\begin{proposition}[r05 enlarged witness register]
The \(W_{\mathrm{OP5}}^{(2)}\) search gives
\[
\begin{array}{c|r|r|r|r|r|r}
\text{witness} & d_{\mathrm{after}} & \ell_{\mathrm{hidden}} &
d_{\mathrm{chart,new}} & r_{\mathrm{reopen}} &
s_{\mathrm{samevar}}^{\mathrm{unlic}} & \Delta_{\mathrm{cpl}}\\
\hline
U_0 & 1 & 0 & 0 & 0 & 0 & 0\\
U_1 & 1 & 0 & 1 & 0 & 0 & 1\\
U_2 & 1 & 0 & 0 & 0 & 0 & 0\\
U_3 & 0 & 1 & 0 & 0 & 0 & 0\\
U_4 & 1 & 0 & 0 & 0 & 0 & 0\\
U_5 & 0 & 0 & 0 & 0 & 1 & -1
\end{array}
\]
where every row has \(d_{\mathrm{before}}=1\).
The raw minimum is
\[
    \min \Delta_{\mathrm{cpl}} = -1,
\]
but the minimum over admissible witnesses is
\[
    \min_{\mathrm{adm}}\Delta_{\mathrm{cpl}} = 0.
\]
\end{proposition}

\begin{proof}
The first, second, and fourth neutral channels preserve compatibility but do
not reduce the residual.  The chart-tilt reroute creates new chart debt.  The
explicit leakage account pays visible debt but records a compensating leakage
unit.  The composite route-chart-carrier witness has formal
\(\Delta_{\mathrm{cpl}}=-1\), but it uses an unlicensed same-variance
comparison.  Therefore it is a conditional strict candidate, not an admissible
strict reducer in r05.
\end{proof}

\begin{theorem}[r05 conditional candidate without admissible promotion]
The enlarged witness search \(W_{\mathrm{OP5}}^{(2)}\) contains no admissible
strict reducer:
\[
    N_{\mathrm{strict,adm}}=0.
\]
It does contain one conditional strict candidate whose only blocking marker is
the unlicensed same-variance comparison:
\[
    U_5 \in \mathrm{conditional\_strict\_candidate\_unlicensed}.
\]
Consequently the OP5 promotion residual remains
\[
    |R_{\mathrm{OP5}}^{\mathrm{prom}}|=1.
\]
The next admissible step is a two-sided tensor bridge license test for the
\(U_5\) comparison.  No global OP5 closure is claimed in r05.
\end{theorem}


% --- v4.5.0-r06 Two-Sided Tensor Bridge License Test ---

\section{Two-Sided Tensor Bridge License Test}
\label{sec:v450-r06-two-sided-tensor-bridge-license-test}

This section tests whether the conditional strict candidate \(U_5\) from
v4.5.0-r05 can be licensed by a two-sided tensor bridge.  The test is not a
global OP5 closure claim.  It is a grammar-license test for one previously
blocked comparison.

\begin{definition}[Two-sided HC/MM-chart comparison bridge]
A two-sided comparison bridge is a typed comparison object
\[
    \mathfrak B_{\mathrm{HC/MM}\leftrightarrow\mathrm{chart}}^{2s}
    :
    \mathcal R_{\mathrm{readout}}
    \times
    \mathcal S_{\mathrm{structure}}
    \longrightarrow
    \mathcal C_{\mathrm{cmp}},
\]
which separates readout-facing indices from structure-facing indices before
comparison.  A same-variance comparison is licensed only if it factors through
\(\mathfrak B_{\mathrm{HC/MM}\leftrightarrow\mathrm{chart}}^{2s}\).  Raw
same-variance contraction remains forbidden.
\end{definition}

\begin{definition}[Bridge license residual]
For the \(U_5\) comparison, define the bridge license residual
\[
    L_{2s}(U_5)
\]
as the finite ledger of missing license conditions: untyped readout side,
untyped structure side, missing comparison tensor, unmatched variance,
remaining raw same-variance contraction, hidden leakage, new chart debt, or
return-side reopening.
\end{definition}

\begin{proposition}[r06 bridge license ledger]
The \(U_5\) bridge-license ledger gives
\[
\begin{array}{c|r}
\text{license component} & \text{residual}\\
\hline
\text{untyped readout side} & 0\\
\text{untyped structure side} & 0\\
\text{missing comparison tensor} & 0\\
\text{unmatched variance after bridge} & 0\\
\text{raw same-variance contraction} & 0\\
\text{hidden leakage} & 0\\
\text{new chart debt} & 0\\
\text{return-side reopening} & 0
\end{array}
\]
and hence
\[
    |L_{2s}(U_5)|=0.
\]
\end{proposition}

\begin{proof}
The \(U_5\) comparison is rewritten as a two-sided comparison: the readout
side and the structure side are separately typed, the comparison tensor is
declared, the variance matching is mediated by the bridge, and the raw
same-variance contraction is removed.  The rewrite introduces no hidden
leakage, no new chart debt, and no return-side reopening.  Therefore the
license residual vanishes.
\end{proof}

\begin{theorem}[U5 becomes an admissible strict OP5 promotion witness]
In r05, \(U_5\) had raw
\[
    \Delta_{\mathrm{cpl}}(U_5)=-1
\]
but was blocked by an unlicensed same-variance comparison.  In r06,
\[
    |L_{2s}(U_5)|=0,
\]
so \(U_5\) becomes an admissible strict reducer inside the declared two-sided
HC/MM-chart comparison grammar.  Consequently
\[
    N_{\mathrm{strict,adm}}^{r06}=1,
    \qquad
    |R_{\mathrm{OP5}}^{\mathrm{prom}}|=0
\]
inside this declared two-sided bridge grammar.
\end{theorem}

\begin{remark}[Scope guard]
The r06 result is a promoted OP5 closure result inside the declared two-sided
HC/MM-chart comparison grammar.  It is not an unrestricted global OP5 theorem,
not a global physical closure theorem, and not a release-final statement.
\end{remark}


% --- v4.5.0-r07 OP5 Two-Sided Promoted-Closure Theorem ---

\section{OP5 Two-Sided Promoted-Closure Theorem}
\label{sec:v450-r07-op5-two-sided-promoted-closure-theorem}

This section records the theorem-level consequence of the r06 two-sided bridge
license test.  The result is scoped to the declared two-sided HC/MM-chart
comparison grammar.  It is not an unrestricted global OP5 theorem and not a
global physical closure theorem.

\begin{definition}[Declared two-sided OP5 promotion grammar]
Let
\[
    \mathcal G_{\mathrm{OP5}}^{2s}
\]
be the declared two-sided OP5 promotion grammar consisting of: the local OP5
carrier/invariant closure import, the \(W_{\mathrm{OP5}}^{(1)}\) and
\(W_{\mathrm{OP5}}^{(2)}\) witness registers, the two-sided
HC/MM-chart comparison bridge
\[
    \mathfrak B_{\mathrm{HC/MM}\leftrightarrow\mathrm{chart}}^{2s},
\]
and the no-leakage, no-new-chart-debt, no-reopening, and no-raw-same-variance
guards.
\end{definition}

\begin{definition}[Two-sided OP5 promoted closure]
Inside \(\mathcal G_{\mathrm{OP5}}^{2s}\), OP5 is two-sided promoted-closed
when the OP5 local carrier closure is preserved and the promotion residual
vanishes:
\[
    \mathrm{Closure}_{\mathrm{loc}}(\mathrm{OP5})
    \quad\text{and}\quad
    R_{\mathrm{OP5}}^{\mathrm{prom}}=\varnothing.
\]
This definition is explicitly grammar-scoped.
\end{definition}

\begin{proposition}[r01-r06 dependency ledger]
The r07 theorem depends on the following finite ledger:
\[
\begin{array}{c|c}
\text{stage} & \text{status}\\
\hline
r01 & \text{local OP5 carrier/invariant closure imported}\\
r02 & \text{first-family HC/MM witness search executed}\\
r03 & \text{first-family obstruction classified}\\
r04 & W_{\mathrm{OP5}}^{(2)}\text{ declared}\\
r05 & U_5\text{ found as conditional strict candidate}\\
r06 & U_5\text{ licensed by the two-sided bridge}
\end{array}
\]
The r06 result gives
\[
    \Delta_{\mathrm{cpl}}(U_5)=-1,
    \qquad
    |L_{2s}(U_5)|=0,
    \qquad
    N_{\mathrm{strict,adm}}=1.
\]
\end{proposition}

\begin{theorem}[OP5 promoted closure inside the two-sided grammar]
Inside the declared two-sided HC/MM-chart comparison grammar
\(\mathcal G_{\mathrm{OP5}}^{2s}\), the licensed witness \(U_5\) is an
admissible strict reducer.  Therefore
\[
    |R_{\mathrm{OP5}}^{\mathrm{prom}}|=0.
\]
Consequently OP5 is promoted-closed inside
\(\mathcal G_{\mathrm{OP5}}^{2s}\):
\[
    \mathrm{Closure}_{\mathcal G_{\mathrm{OP5}}^{2s}}(\mathrm{OP5}).
\]
\end{theorem}

\begin{proof}
The r01 intake imported local OP5 carrier/invariant closure without claiming
global OP5 closure.  The r02 and r03 layers showed that the first witness
family was obstructed.  The r04 layer therefore enlarged the witness family.
The r05 layer found \(U_5\) with raw \(\Delta_{\mathrm{cpl}}=-1\), blocked
only by an unlicensed same-variance comparison.  The r06 layer introduced the
two-sided HC/MM-chart comparison bridge and proved
\[
    |L_{2s}(U_5)|=0.
\]
Thus the only r05 blocker for \(U_5\) is removed without hidden leakage, new
chart debt, or return-side reopening.  Hence \(U_5\) is an admissible strict
reducer in \(\mathcal G_{\mathrm{OP5}}^{2s}\), and the OP5 promotion residual
vanishes in that grammar.
\end{proof}

\begin{remark}[No-overclaim boundary]
The theorem proves OP5 promoted closure only inside the declared two-sided
HC/MM-chart comparison grammar.  It does not prove unrestricted global OP5
closure, does not prove global physical closure, does not close QFT/GR, and
does not mark a release-final state.
\end{remark}


% --- v4.5.0-r08 No-Overclaim / Route-to-Global Boundary Audit ---

\section{No-Overclaim and Route-to-Global Boundary Audit}
\label{sec:v450-r08-no-overclaim-route-to-global-boundary-audit}

This section audits the boundary of the r07 OP5 two-sided promoted-closure
theorem.  The audit preserves the r07 result inside the declared two-sided
HC/MM-chart comparison grammar, but blocks unrestricted global OP5, all-readout
OP5, global physical, and release-final claims.

\begin{definition}[Route-to-global boundary residual]
Let
\[
    B_{\mathrm{global}}^{\mathrm{OP5}}
\]
denote the boundary residual between the declared two-sided OP5 promotion
grammar \(\mathcal G_{\mathrm{OP5}}^{2s}\) and an unrestricted global OP5
claim.  It is not the same object as \(R_{\mathrm{OP5}}^{\mathrm{prom}}\).
The r07 theorem proves
\[
    R_{\mathrm{OP5}}^{\mathrm{prom}}=0
    \quad\text{inside}\quad
    \mathcal G_{\mathrm{OP5}}^{2s},
\]
but it does not prove
\[
    B_{\mathrm{global}}^{\mathrm{OP5}}=0.
\]
\end{definition}

\begin{proposition}[Allowed and blocked claims]
The r08 claim ledger is
\[
\begin{array}{c|c|c}
\text{claim} & \text{allowed} & \text{scope}\\
\hline
\mathrm{Closure}_{\mathcal G_{\mathrm{OP5}}^{2s}}(\mathrm{OP5})
& \mathrm{yes} & \text{declared two-sided grammar}\\
\text{unrestricted global OP5 closure}
& \mathrm{no} & \text{not established}\\
\text{all-readout OP5 closure}
& \mathrm{no} & \text{not established}\\
\text{global physical closure}
& \mathrm{no} & \text{not established}\\
\text{release-final status}
& \mathrm{no} & \text{working line only}
\end{array}
\]
\end{proposition}

\begin{proposition}[Route-to-global boundary ledger]
The boundary ledger gives
\[
\begin{array}{c|r}
\text{gate} & \text{value}\\
\hline
\text{two-sided OP5 gates passed} & 1\\
\text{unrestricted global gates passed} & 0\\
B_{\mathrm{global}}^{\mathrm{OP5}} & 1
\end{array}
\]
where \(B_{\mathrm{global}}^{\mathrm{OP5}}=1\) is a scope-boundary marker, not
a reopening of \(R_{\mathrm{OP5}}^{\mathrm{prom}}\) inside
\(\mathcal G_{\mathrm{OP5}}^{2s}\).
\end{proposition}

\begin{theorem}[r08 no-overclaim boundary theorem]
The r07 theorem may be stated as
\[
    \mathrm{Closure}_{\mathcal G_{\mathrm{OP5}}^{2s}}(\mathrm{OP5}).
\]
It may not be strengthened in r08 to unrestricted global OP5 closure, all-readout
OP5 closure, global physical closure, or release-final status.  The route from
\(\mathcal G_{\mathrm{OP5}}^{2s}\) to a larger global claim remains guarded by
\[
    B_{\mathrm{global}}^{\mathrm{OP5}}=1.
\]
\end{theorem}

\begin{remark}[Interpretation]
The boundary residual in r08 does not undo the two-sided OP5 promoted closure.
It records exactly where the next promotion problem begins.  In plain terms:
OP5 is promoted closed in the declared two-sided grammar, but the unrestricted
global route remains outside the proved scope.
\end{remark}


% --- v4.5.0-r09 OP5 / OP1-OP4 Compatibility Audit ---

\section{OP5 / OP1--OP4 Compatibility Audit}
\label{sec:v450-r09-op5-op14-compatibility-audit}

This section audits compatibility between the public v4.4 OP1/OP4
declared-grammar closure and the v4.5.0 OP5 two-sided promoted-closure result.
It asks whether the OP5 promotion reopens any OP1/OP4 ledger channel.  It does
not claim unified global closure.

\begin{definition}[Compatibility product grammar]
Let
\[
    \mathcal G_{\mathrm{cmp}}^{1/4,5}
    :=
    \mathcal G_{1/4}^{\mathrm{loc}}
    \boxtimes
    \mathcal G_{\mathrm{OP5}}^{2s}
\]
be the compatibility product of the local OP1/OP4 declared grammar and the
two-sided OP5 promotion grammar.  The product symbol means coexistence with
explicit cross-reopen guards; it does not mean unrestricted global collapse of
the two grammars.
\end{definition}

\begin{definition}[OP5-to-OP1/OP4 reopen residual]
Define
\[
    R_{\mathrm{reopen}}^{1/4\leftarrow 5}
\]
as the finite residual measuring whether the OP5 two-sided promoted closure
reopens any OP1/OP4 local closure channel: scale, lens, chart, debt, selector,
or coverage.
\end{definition}

\begin{proposition}[r09 compatibility ledger]
The compatibility ledger gives
\[
\begin{array}{c|r}
\text{OP1/OP4 reopen channel} & \text{residual}\\
\hline
\text{scale} & 0\\
\text{lens} & 0\\
\text{chart} & 0\\
\text{debt} & 0\\
\text{selector} & 0\\
\text{coverage} & 0\\
\hline
R_{\mathrm{reopen}}^{1/4\leftarrow 5} & 0
\end{array}
\]
so the OP5 two-sided promotion does not reopen the v4.4 OP1/OP4 local closure
ledger.
\end{proposition}

\begin{proof}
The OP5 theorem from r07 is scoped to
\(\mathcal G_{\mathrm{OP5}}^{2s}\), while r08 explicitly keeps the
route-to-global boundary outside the theorem.  The boundary marker
\(B_{\mathrm{global}}^{\mathrm{OP5}}=1\) is not imported as scale, lens, chart,
debt, selector, or coverage debt inside \(\mathcal G_{1/4}^{\mathrm{loc}}\).
Therefore none of the v4.4 OP1/OP4 closure channels reopens.
\end{proof}

\begin{theorem}[Combined scoped compatibility surface]
The following two scoped closure statements are compatible:
\[
    \mathrm{Closure}_{\mathcal G_{1/4}^{\mathrm{loc}}}(\mathrm{OP1},\mathrm{OP4})
\]
and
\[
    \mathrm{Closure}_{\mathcal G_{\mathrm{OP5}}^{2s}}(\mathrm{OP5}).
\]
Equivalently, inside the compatibility product grammar,
\[
    \mathrm{Closure}_{\mathcal G_{\mathrm{cmp}}^{1/4,5}}
    (\mathrm{OP1},\mathrm{OP4};\mathrm{OP5})
\]
is allowed with
\[
    R_{\mathrm{reopen}}^{1/4\leftarrow 5}=0.
\]
This is a combined scoped surface only, not unified global closure.
\end{theorem}

\begin{remark}[Global boundary]
The route-to-global boundary remains
\[
    B_{\mathrm{global}}^{\mathrm{OP5}}=1.
\]
This boundary marker prevents the combined scoped surface from being rewritten
as unrestricted global OP1/OP4/OP5 closure.
\end{remark}


% --- v4.5.0-r10 Two-Sided Tensor Grammar Consolidation ---

\section{Two-Sided Tensor Grammar Consolidation}
\label{sec:v450-r10-two-sided-tensor-grammar-consolidation}

This section consolidates the notation used by the v4.4 OP1/OP4 local closure
and the v4.5.0 OP5 two-sided promoted closure.  It is a grammar and notation
layer only.  It preserves the scoped compatibility surface from r09 and does
not claim unified global closure.

\begin{definition}[r10 grammar normal form]
Set
\[
    \mathcal G_L := \mathcal G_{1/4}^{\mathrm{loc}},
    \qquad
    \mathcal G_T := \mathcal G_{\mathrm{OP5}}^{2s}.
\]
Here \(\mathcal G_L\) is the local declared OP1/OP4 grammar and
\(\mathcal G_T\) is the two-sided OP5 HC/MM-chart comparison grammar.  Their
compatibility product is written
\[
    \mathcal G_C
    :=
    \mathcal G_L \boxtimes_{\mathfrak C} \mathcal G_T,
\]
where \(\mathfrak C\) is an explicit comparison tensor.  The symbol
\(\boxtimes_{\mathfrak C}\) denotes guarded compatibility, not unrestricted
global tensor contraction.
\end{definition}

\begin{definition}[Two-sided comparison tensor]
The comparison tensor
\[
    \mathfrak C:
    \mathcal R_{\mathrm{readout}}
    \times
    \mathcal S_{\mathrm{structure}}
    \longrightarrow
    \mathcal C_{\mathrm{cmp}}
\]
separates readout-facing indices from structure-facing indices before any
comparison is made.  A same-variance comparison is admissible only when it
factors through \(\mathfrak C\).  Raw same-variance contraction remains
forbidden.
\end{definition}

\begin{proposition}[r10 grammar rules]
The consolidated grammar obeys the following rules:
\[
\begin{array}{c|l}
\text{rule} & \text{content}\\
\hline
\mathrm{TSG1} & \text{readout and structure indices remain typed}\\
\mathrm{TSG2} & \text{same-variance comparison requires }\mathfrak C\\
\mathrm{TSG3} & \boxtimes_{\mathfrak C}\text{ is not global grammar collapse}\\
\mathrm{TSG4} & B_G\text{ remains outside the scoped closure}\\
\mathrm{TSG5} & R_{\mathrm{reopen}}^{1/4\leftarrow 5}=0
\end{array}
\]
where \(B_G\) is the route-to-global boundary marker inherited from r08.
\end{proposition}

\begin{theorem}[Consolidated scoped grammar surface]
In the r10 normal form, the two scoped closure statements are
\[
    \mathrm{Closure}_{\mathcal G_L}(\mathrm{OP1},\mathrm{OP4})
\]
and
\[
    \mathrm{Closure}_{\mathcal G_T}(\mathrm{OP5}).
\]
The compatibility audit from r09 gives
\[
    R_{\mathrm{reopen}}^{1/4\leftarrow 5}=0,
\]
so they define the combined scoped surface
\[
    \mathrm{Closure}_{\mathcal G_C}(\mathrm{OP1},\mathrm{OP4};\mathrm{OP5}).
\]
This theorem is a consolidated notation theorem for the scoped product
\(\mathcal G_C\).  It is not a unified global closure theorem.
\end{theorem}

\begin{remark}[Boundary preserved]
The global boundary marker remains
\[
    B_G = 1.
\]
It does not reopen the scoped product closure.  It marks exactly the route
from the combined scoped surface to any future larger global claim.
\end{remark}


% --- v4.5.0-r11 Compatibility Theorem / No-Reopen Summary ---

\section{Compatibility Theorem and No-Reopen Summary}
\label{sec:v450-r11-compatibility-theorem-no-reopen-summary}

This section summarizes the compatibility result after the r10 two-sided tensor
grammar consolidation.  It also records a short notation guide for the symbol
\(\boxtimes_{\mathfrak C}\), since this symbol is not meant as an ordinary
multiplication sign.

\begin{definition}[Guarded compatibility product]
The symbol
\[
    \boxtimes_{\mathfrak C}
\]
is read as a guarded compatibility product over the comparison tensor
\(\mathfrak C\).  In the present setting
\[
    \mathcal G_C
    =
    \mathcal G_L \boxtimes_{\mathfrak C} \mathcal G_T
\]
means: combine the local OP1/OP4 grammar \(\mathcal G_L\) and the two-sided
OP5 grammar \(\mathcal G_T\) only through the explicit comparison tensor
\(\mathfrak C\).  It does not mean ordinary multiplication, raw tensor
contraction, or unrestricted global grammar collapse.
\end{definition}

\begin{proposition}[r11 no-reopen ledger]
The r11 no-reopen summary is
\[
\begin{array}{c|r}
\text{component} & \text{residual}\\
\hline
\mathrm{Closure}_{\mathcal G_L}(\mathrm{OP1},\mathrm{OP4}) & 0\\
\mathrm{Closure}_{\mathcal G_T}(\mathrm{OP5}) & 0\\
R_{\mathrm{reopen}}^{1/4\leftarrow 5} & 0\\
\mathrm{Closure}_{\mathcal G_C}(\mathrm{OP1},\mathrm{OP4};\mathrm{OP5}) & 0\\
B_G & 1
\end{array}
\]
where \(B_G=1\) is the retained route-to-global boundary marker.
\end{proposition}

\begin{theorem}[Combined scoped no-reopen theorem]
In the r10 normal form, the local OP1/OP4 closure and the OP5 two-sided
promoted closure coexist in the guarded compatibility grammar
\[
    \mathcal G_C
    =
    \mathcal G_L \boxtimes_{\mathfrak C} \mathcal G_T.
\]
Moreover,
\[
    R_{\mathrm{reopen}}^{1/4\leftarrow 5}=0.
\]
Therefore
\[
    \mathrm{Closure}_{\mathcal G_C}(\mathrm{OP1},\mathrm{OP4};\mathrm{OP5})
\]
is a valid combined scoped closure surface.
\end{theorem}

\begin{proof}
The v4.4 line provides the local declared OP1/OP4 closure in
\(\mathcal G_L\).  The r07 line provides OP5 promoted closure in
\(\mathcal G_T\).  The r09 compatibility audit shows that the OP5 closure
does not reopen any OP1/OP4 scale, lens, chart, debt, selector, or coverage
channel.  The r10 grammar consolidation then rewrites the combined statement
through the explicit comparison tensor \(\mathfrak C\).  Hence the combined
scoped surface is valid with zero cross-reopen residual.
\end{proof}

\begin{remark}[Boundary retained]
The theorem is intentionally scoped.  The boundary
\[
    B_G=1
\]
is retained.  Thus the theorem is not
\[
    \mathrm{Closure}_{\mathrm{global}}(\mathrm{OP1},\mathrm{OP4},\mathrm{OP5}).
\]
It is only the combined scoped closure theorem inside \(\mathcal G_C\).
\end{remark}


% --- v4.5.0-r13 Route-to-Global Boundary Problem Statement ---

\section{Route-to-Global Boundary Problem Statement}
\label{sec:v450-r12-route-to-global-boundary-problem-statement}

This section states the next problem after the r11 combined scoped no-reopen
theorem.  The scoped theorem remains valid:
\[
    \mathrm{Closure}_{\mathcal G_C}(\mathrm{OP1},\mathrm{OP4};\mathrm{OP5}),
    \qquad
    R_{\mathrm{reopen}}^{1/4\leftarrow 5}=0.
\]
However, the route-to-global boundary marker remains
\[
    B_G=1.
\]
The purpose of r13 is to name the exact boundary problem, not to solve it.

\begin{definition}[Route-to-global promotion problem]
The route-to-global promotion problem is the problem of deciding whether the
scoped compatibility grammar
\[
    \mathcal G_C
    =
    \mathcal G_L \boxtimes_{\mathfrak C} \mathcal G_T
\]
can be embedded into a larger grammar \(\mathcal G_{\mathrm{global}}\) without
creating new debt, losing comparison typing, or reopening any OP1/OP4 or OP5
closure channel.  The problem is denoted
\[
    \mathrm{Promote}(\mathcal G_C\to \mathcal G_{\mathrm{global}}).
\]
No such promotion is assumed in r12.
\end{definition}

\begin{definition}[Global-boundary gate ledger]
The boundary marker \(B_G\) is discharged only if the following gates are
closed:
\[
\begin{array}{c|l}
\text{gate} & \text{meaning}\\
\hline
\mathrm{GLOB1} & \text{grammar exhaustion}\\
\mathrm{GLOB2} & \text{all-readout invariance}\\
\mathrm{GLOB3} & \text{boundary transport without new debt}\\
\mathrm{GLOB4} & \text{no-new-debt under grammar extension}\\
\mathrm{GLOB5} & \text{functorial comparison stability}
\end{array}
\]
The physical-interpretation guard is separate: even a mathematical grammar
promotion would not by itself be a global physical closure theorem.
\end{definition}

\begin{proposition}[r13 boundary ledger]
At r12, the global-boundary ledger is
\[
\begin{array}{c|r}
\text{quantity} & \text{value}\\
\hline
\text{required global gates} & 5\\
\text{passed global gates} & 0\\
\text{open global gates} & 5\\
B_G & 1
\end{array}
\]
Therefore the route-to-global boundary remains open.
\end{proposition}

\begin{theorem}[Scoped closure with open global boundary]
The r11 scoped no-reopen theorem remains valid:
\[
    \mathrm{Closure}_{\mathcal G_C}(\mathrm{OP1},\mathrm{OP4};\mathrm{OP5})
\]
and
\[
    R_{\mathrm{reopen}}^{1/4\leftarrow 5}=0.
\]
But r13 does not prove
\[
    \mathrm{Closure}_{\mathrm{global}}(\mathrm{OP1},\mathrm{OP4},\mathrm{OP5}).
\]
Instead, it records the precise open boundary:
\[
    B_G=1.
\]
\end{theorem}

\begin{remark}[Interpretation]
In plain language, v4.5.0 has reached a combined scoped closure surface.  The
next problem is no longer OP5 local promotion; it is the route from the scoped
surface to any larger/global grammar.  That route requires new gates, not just
renaming the scoped theorem as global.
\end{remark}


% --- v4.5.0-r14 Reader-Facing Notation Cleanup ---

\section{Reader-Facing Notation Cleanup}
\label{sec:v450-r14-reader-facing-notation-cleanup}

This section is an expository notation cleanup.  It introduces no new closure
claim.  Its purpose is to make the symbols from the v4.5.0 working line readable
without weakening the no-overclaim guards.

\begin{definition}[Boxed-times compatibility product]
The symbol
\[
    \boxtimes_{\mathfrak C}
\]
is read as a \emph{guarded compatibility product over the comparison tensor}
\(\mathfrak C\).  Thus
\[
    \mathcal G_C
    =
    \mathcal G_L \boxtimes_{\mathfrak C} \mathcal G_T
\]
means that the local OP1/OP4 grammar \(\mathcal G_L\) and the two-sided OP5
grammar \(\mathcal G_T\) are compared only through the explicit bridge
\(\mathfrak C\).  It does not mean ordinary multiplication, raw tensor
contraction, or unrestricted global grammar collapse.
\end{definition}

\begin{definition}[Reader glossary]
The working symbols should be read as follows:
\begin{center}
\begin{tabular}{>{\raggedright\arraybackslash}p{0.22\linewidth}|>{\raggedright\arraybackslash}p{0.68\linewidth}}
\textbf{symbol} & \textbf{reader meaning}\\
\hline
\(\mathcal G_L\) & the local declared OP1/OP4 grammar from the v4.4 line\\
\(\mathcal G_T\) & the two-sided OP5 HC/MM-chart comparison grammar\\
\(\mathfrak C\) & the explicit comparison tensor separating readout and structure sides\\
\(\mathcal G_C\) & the combined scoped grammar \(\mathcal G_L\boxtimes_{\mathfrak C}\mathcal G_T\)\\
\(B_G\) & the retained route-to-global boundary marker\\
\(R_{\mathrm{reopen}}^{1/4\leftarrow 5}\) & the cross-reopen residual from OP5 back to OP1/OP4\\
\(\mathrm{GLOB1}\text{--}\mathrm{GLOB5}\) & the five open gates for any future route-to-global promotion
\end{tabular}
\end{center}
\end{definition}

\begin{proposition}[Reader-safe separation of residuals]
The two residual markers
\[
    R_{\mathrm{reopen}}^{1/4\leftarrow 5}
    \qquad\text{and}\qquad
    B_G
\]
have different jobs.  The first asks whether OP5 reopens the OP1/OP4 local
closure ledger.  In the r11 summary this residual is zero.  The second records
that the route from the scoped grammar \(\mathcal G_C\) to a larger/global
grammar is still open.  In r13 this marker remains
\[
    B_G=1.
\]
Thus the global boundary marker is not a reopened scoped residual.
\end{proposition}

\begin{remark}[Plain-language reading]
The expression
\[
    \mathrm{Closure}_{\mathcal G_C}(\mathrm{OP1},\mathrm{OP4};\mathrm{OP5})
\]
should be read as: OP1/OP4 and OP5 are closed together on the declared guarded
compatibility surface.  It should not be read as: all global readouts and all
physical interpretations are closed.  The latter would require the open gates
\(\mathrm{GLOB1}\) through \(\mathrm{GLOB5}\), which remain open at this stage.
\end{remark}

\begin{theorem}[Reader-clean notation handoff]
The r14 cleanup preserves the r11/r13 state:
\[
    \mathrm{Closure}_{\mathcal G_C}(\mathrm{OP1},\mathrm{OP4};\mathrm{OP5}),
    \qquad
    R_{\mathrm{reopen}}^{1/4\leftarrow 5}=0,
    \qquad
    B_G=1.
\]
It introduces no global closure theorem, no global physical closure theorem,
and no release-final state.
\end{theorem}


% --- v4.5.0-r15 Global-Boundary Witness-Menu Design ---

\section{Global-Boundary Witness-Menu Design}
\label{sec:v450-r15-global-boundary-witness-menu-design}

This section designs finite witness menus for the five open route-to-global
gates from r12.  It does not execute the menus.  It does not discharge
\(B_G\), and it does not claim global closure.

\begin{definition}[Global-boundary witness families]
For the five open gates
\[
    \mathrm{GLOB1},\ldots,\mathrm{GLOB5},
\]
define the witness families
\[
\begin{array}{c|c|l}
\text{gate} & \text{family} & \text{role}\\
\hline
\mathrm{GLOB1} & W_{\mathrm{exh}} & \text{grammar exhaustion}\\
\mathrm{GLOB2} & W_{\mathrm{inv}} & \text{all-readout invariance}\\
\mathrm{GLOB3} & W_{\mathrm{tr}} & \text{boundary transport}\\
\mathrm{GLOB4} & W_{\mathrm{debt}} & \text{no-new-debt extension}\\
\mathrm{GLOB5} & W_{\mathrm{fun}} & \text{functorial comparison stability}
\end{array}
\]
These families are finite menus for future execution, not proof objects in r15.
\end{definition}

\begin{proposition}[r15 witness-menu ledger]
The r15 design ledger is
\[
\begin{array}{c|r}
\text{quantity} & \text{value}\\
\hline
\text{witness menus designed} & 5\\
\text{witness menus executed} & 0\\
\text{new global gates passed} & 0\\
\text{open global gates after r15} & 5\\
B_G & 1
\end{array}
\]
Therefore r15 is a design step only.
\end{proposition}

\begin{definition}[Execution order]
The intended execution order is
\[
    W_{\mathrm{exh}}
    \longrightarrow
    W_{\mathrm{inv}}
    \longrightarrow
    W_{\mathrm{tr}}
    \longrightarrow
    W_{\mathrm{debt}}
    \longrightarrow
    W_{\mathrm{fun}}.
\]
Grammar exhaustion comes first because all-readout invariance and boundary
transport are underspecified if admissible branches have not first been
classified.
\end{definition}

\begin{theorem}[No-discharge theorem for r15]
The r15 layer defines the finite witness menus needed for the route-to-global
program, but it does not execute any witness family.  Hence
\[
    B_G=1
\]
remains unchanged, and no global OP1/OP4/OP5 closure is claimed.
\end{theorem}

\begin{remark}[Plain reading]
The scoped closure surface from r11 is still intact.  r15 only turns the open
global-boundary problem into a concrete work queue.
\end{remark}


% --- v4.5.0-r16 Execute W_exh Grammar Exhaustion ---

\section{Grammar Exhaustion Witness Execution}
\label{sec:v450-r16-execute-wexh-grammar-exhaustion}

This section executes the first global-boundary witness family from r15:
\[
    W_{\mathrm{exh}}=\{E_0,E_1,E_2\}.
\]
It tests only the grammar-exhaustion gate \(\mathrm{GLOB1}\).  It does not
test all-readout invariance, boundary transport, no-new-debt under extension,
or functorial comparison stability.

\begin{definition}[Grammar-exhaustion witness family]
The family \(W_{\mathrm{exh}}\) consists of:
\[
\begin{array}{c|l}
E_0 & \text{atlas-coverage witness}\\
E_1 & \text{readout-factorization witness}\\
E_2 & \text{no-orphan-branch witness}
\end{array}
\]
The family passes if every admissible branch of the current scoped surface is
represented in \(\mathcal G_C\) or routed to a declared open boundary, and no
admissible branch remains orphan.
\end{definition}

\begin{proposition}[r16 branch classification register]
The branch register is
\[
\begin{array}{c|c}
\text{branch} & \text{classification}\\
\hline
C_0\ \text{OP1/OP4 local} & \text{represented}\\
C_1\ \text{OP5 two-sided} & \text{represented}\\
C_2\ \text{compatibility product} & \text{represented}\\
C_3\ \text{route-to-global boundary} & \text{routed to }B_G\\
C_4\ \text{raw same-variance comparison} & \text{blocked, not admissible}\\
C_5\ \text{global physical interpretation} & \text{blocked, not admissible}
\end{array}
\]
Thus
\[
    N_{\mathrm{orphan}}=0,
    \qquad
    N_{\mathrm{new\,debt}}=0.
\]
\end{proposition}

\begin{theorem}[GLOB1 grammar-exhaustion gate passed]
The witness family \(W_{\mathrm{exh}}\) discharges the grammar-exhaustion gate:
\[
    \mathrm{GLOB1}=1.
\]
After r16 the global-boundary ledger is
\[
\begin{array}{c|r}
\text{quantity} & \text{value}\\
\hline
\text{passed global gates} & 1\\
\text{open global gates} & 4\\
B_G & 1
\end{array}
\]
The boundary marker remains \(B_G=1\) because \(\mathrm{GLOB2}\) through
\(\mathrm{GLOB5}\) remain open.
\end{theorem}

\begin{remark}[No global closure claim]
The r16 result is a coverage/exhaustion result only.  It does not prove global
promotion and does not claim global OP1/OP4/OP5 closure.
\end{remark}


% --- v4.5.0-r17 Execute W_inv Declared Readout Invariance ---

\section{Declared Readout Invariance Witness Execution}
\label{sec:v450-r17-execute-winv-readout-invariance}

This section executes the second global-boundary witness family:
\[
    W_{\mathrm{inv}}=\{I_0,I_1,I_2\}.
\]
The result is scoped to the declared admissible readout transformations in the
current grammar.  It is not an unrestricted all-possible-readouts theorem.

\begin{definition}[Declared readout-invariance witness family]
The family \(W_{\mathrm{inv}}\) consists of:
\[
\begin{array}{c|l}
I_0 & \text{declared chart-change witness}\\
I_1 & \text{declared readout-change witness}\\
I_2 & \text{declared return-phase-change witness}
\end{array}
\]
It passes if the scoped closure surface is preserved, no OP1/OP4 or OP5
channel reopens, no new debt is created, and the boundary marker \(B_G\) is not
rewritten as a solved global claim.
\end{definition}

\begin{proposition}[r17 readout-invariance register]
The declared readout-invariance register is
\[
\begin{array}{c|c|r|r|r}
\text{witness} & \text{classification} & R_{\mathrm{reopen}} & N_{\mathrm{debt}} & N_{\mathrm{boundary\ rewrite}}\\
\hline
I_0 & \text{invariant} & 0 & 0 & 0\\
I_1 & \text{invariant} & 0 & 0 & 0\\
I_2 & \text{invariant} & 0 & 0 & 0
\end{array}
\]
Thus the declared readout-invariance test has no reopen, no new debt, and no
boundary rewrite.
\end{proposition}

\begin{theorem}[GLOB2 declared-readout-invariance gate passed]
The witness family \(W_{\mathrm{inv}}\) discharges the declared
readout-invariance gate:
\[
    \mathrm{GLOB2}=1.
\]
After r17 the global-boundary ledger is
\[
\begin{array}{c|r}
\text{quantity} & \text{value}\\
\hline
\text{passed global gates} & 2\\
\text{open global gates} & 3\\
B_G & 1
\end{array}
\]
The boundary marker remains \(B_G=1\) because \(\mathrm{GLOB3}\) through
\(\mathrm{GLOB5}\) remain open.
\end{theorem}

\begin{remark}[Scope guard]
The word ``all'' in this r17 layer means all declared/admissible readout
transformations currently in the grammar.  It does not mean an unrestricted
theorem over every possible external readout or every possible physical model.
\end{remark}


% --- v4.5.0-r18 Execute W_tr Boundary Transport ---

\section{Boundary Transport Witness Execution}
\label{sec:v450-r18-execute-wtr-boundary-transport}

This section executes the third global-boundary witness family:
\[
    W_{\mathrm{tr}}=\{T_0,T_1,T_2\}.
\]
The purpose is to test whether the route-to-global boundary marker \(B_G\) can
be transported through the interface without creating new debt.  Transport is
not discharge: \(B_G\) remains an explicit boundary marker.

\begin{definition}[Boundary-transport witness family]
The family \(W_{\mathrm{tr}}\) consists of:
\[
\begin{array}{c|l}
T_0 & \text{boundary-carry witness}\\
T_1 & \text{debt-neutral transport witness}\\
T_2 & \text{OP5-return transport witness}
\end{array}
\]
It passes if \(B_G\) is carried as a boundary marker and no new scale, lens,
chart, selector, coverage, or OP5-return debt appears.
\end{definition}

\begin{proposition}[r18 boundary-transport register]
The transport register is
\[
\begin{array}{c|c|r|r}
\text{witness} & \text{classification} & N_{\mathrm{new\,debt}} & B_G\ \text{after}\\
\hline
T_0 & \text{transported boundary marker} & 0 & 1\\
T_1 & \text{debt neutral} & 0 & 1\\
T_2 & \text{return transport stable} & 0 & 1
\end{array}
\]
Thus boundary transport is stable and debt-free, but the boundary marker is
not erased.
\end{proposition}

\begin{theorem}[GLOB3 boundary-transport gate passed]
The witness family \(W_{\mathrm{tr}}\) discharges the boundary-transport gate:
\[
    \mathrm{GLOB3}=1.
\]
After r18 the global-boundary ledger is
\[
\begin{array}{c|r}
\text{quantity} & \text{value}\\
\hline
\text{passed global gates} & 3\\
\text{open global gates} & 2\\
B_G & 1
\end{array}
\]
The marker \(B_G=1\) remains because \(\mathrm{GLOB4}\) and
\(\mathrm{GLOB5}\) remain open.
\end{theorem}

\begin{remark}[No discharge by transport]
The r18 result says that the boundary can be carried without new debt.  It
does not say that the boundary is solved.  Therefore no global OP1/OP4/OP5
closure is claimed.
\end{remark}


% --- v4.5.0-r19 Execute W_debt No-New-Debt Extension ---

\section{No-New-Debt Extension Witness Execution}
\label{sec:v450-r19-execute-wdebt-no-new-debt-extension}

This section executes the fourth global-boundary witness family:
\[
    W_{\mathrm{debt}}=\{D_0,D_1,D_2,D_3,D_4,D_5\}.
\]
The purpose is to test whether the declared grammar-extension buckets create
new debt.  This is still not a global closure claim, because the functorial
comparison-stability gate remains open.

\begin{definition}[No-new-debt extension witness family]
The family \(W_{\mathrm{debt}}\) consists of:
\[
\begin{array}{c|l}
D_0 & \text{scale channel}\\
D_1 & \text{lens channel}\\
D_2 & \text{chart channel}\\
D_3 & \text{selector channel}\\
D_4 & \text{coverage channel}\\
D_5 & \text{OP5-return channel}
\end{array}
\]
It passes if none of these channels creates new debt, hidden leakage, or
reopening under the declared grammar extension.
\end{definition}

\begin{proposition}[r19 no-new-debt register]
The no-new-debt register is
\[
\begin{array}{c|c|r|r|r}
\text{witness} & \text{channel} & N_{\mathrm{new\,debt}} & R_{\mathrm{reopen}} & L_{\mathrm{leak}}\\
\hline
D_0 & \text{scale} & 0 & 0 & 0\\
D_1 & \text{lens} & 0 & 0 & 0\\
D_2 & \text{chart} & 0 & 0 & 0\\
D_3 & \text{selector} & 0 & 0 & 0\\
D_4 & \text{coverage} & 0 & 0 & 0\\
D_5 & \text{OP5-return} & 0 & 0 & 0
\end{array}
\]
Thus the declared grammar extension is debt-free for the tested channels.
\end{proposition}

\begin{theorem}[GLOB4 no-new-debt extension gate passed]
The witness family \(W_{\mathrm{debt}}\) discharges the no-new-debt extension
gate:
\[
    \mathrm{GLOB4}=1.
\]
After r19 the global-boundary ledger is
\[
\begin{array}{c|r}
\text{quantity} & \text{value}\\
\hline
\text{passed global gates} & 4\\
\text{open global gates} & 1\\
B_G & 1
\end{array}
\]
The marker \(B_G=1\) remains because \(\mathrm{GLOB5}\), functorial comparison
stability, remains open.
\end{theorem}

\begin{remark}[No global closure claim]
The r19 result says that the declared extension buckets are debt-free.  It
does not say that the comparison tensor is functorially stable under all needed
extension maps.  Therefore no global OP1/OP4/OP5 closure is claimed in r19.
\end{remark}


% --- v4.5.0-r20 Execute W_fun Functorial Comparison Stability ---

\section{Functorial Comparison Stability Witness Execution}
\label{sec:v450-r20-execute-wfun-functorial-comparison-stability}

This section executes the fifth global-boundary witness family:
\[
    W_{\mathrm{fun}}=\{F_0,F_1,F_2,F_3\}.
\]
The purpose is to test whether the comparison tensor \(\mathfrak C\) remains
stable under the declared extension maps.  Passing this gate makes the
boundary synthesis-ready, but it does not itself discharge \(B_G\).

\begin{definition}[Functorial comparison-stability witness family]
The family \(W_{\mathrm{fun}}\) consists of:
\[
\begin{array}{c|l}
F_0 & \mathfrak C\text{-naturalness}\\
F_1 & \text{composition stability}\\
F_2 & \text{identity stability}\\
F_3 & \text{variance preservation}
\end{array}
\]
It passes if the comparison tensor commutes with the declared extension maps,
respects composition and identity, preserves variance typing, and introduces
no raw same-variance comparison.
\end{definition}

\begin{proposition}[r20 functorial-stability register]
The functorial comparison register is
\[
\begin{array}{c|c|r}
\text{witness} & \text{condition} & \text{defect}\\
\hline
F_0 & \mathfrak C\text{-naturalness} & 0\\
F_1 & \text{composition} & 0\\
F_2 & \text{identity} & 0\\
F_3 & \text{variance preservation} & 0
\end{array}
\]
Thus the declared comparison tensor is functorially stable for the tested
extension maps.
\end{proposition}

\begin{theorem}[GLOB5 functorial-comparison-stability gate passed]
The witness family \(W_{\mathrm{fun}}\) discharges the functorial comparison
stability gate:
\[
    \mathrm{GLOB5}=1.
\]
After the functorial-comparison-stability gate, the global-boundary ledger is
\[
\begin{array}{c|r}
\text{quantity} & \text{value}\\
\hline
\text{passed global gates} & 5\\
\text{open global gates} & 0\\
B_G & 1\\
\text{synthesis-ready} & 1
\end{array}
\]
The marker \(B_G\) is not discharged by that individual gate.  It is carried into the
boundary-gate synthesis audit.
\end{theorem}

\begin{remark}[No premature global closure]
The functorial-stability result is the last individual gate execution.  It does not by itself
claim global OP1/OP4/OP5 closure.  The next layer must synthesize the five
passed gates and decide whether \(B_G\) can be discharged without overclaim.
\end{remark}


% --- v4.5.0-r21 Boundary-Gate Synthesis Audit ---

\section{Boundary-Gate Synthesis Audit}
\label{sec:v450-r21-boundary-gate-synthesis-audit}

This section synthesizes the five previously executed global-boundary gates.
It decides whether the retained boundary marker \(B_G\) may be discharged
inside the declared mathematical declared global-boundary synthesis grammar.  The result is not a
global physical closure theorem and not an unrestricted claim beyond the
declared extension maps.

\begin{definition}[Declared declared global-boundary synthesis grammar]
Let
\[
    \mathcal G_{\mathrm{global}}^{\mathrm{decl}}
\]
be the grammar obtained from the scoped compatibility grammar \(\mathcal G_C\)
by adjoining the five passed global-boundary gates:
\[
    \mathrm{GLOB1},\ldots,\mathrm{GLOB5}.
\]
This is a declared mathematical grammar.  It is not the class of all possible
external readouts or all possible physical models.
\end{definition}

\begin{proposition}[r21 synthesis ledger]
The r21 synthesis ledger is
\[
\begin{array}{c|r}
\text{quantity} & \text{value}\\
\hline
\text{global-boundary gates} & 5\\
\text{passed gates} & 5\\
\text{open gates} & 0\\
B_G\ \text{before synthesis} & 1\\
B_G\ \text{after synthesis} & 0
\end{array}
\]
All five global-boundary gates are passed, so the boundary marker is
discharged inside \(\mathcal G_{\mathrm{global}}^{\mathrm{decl}}\).
\end{proposition}

\begin{theorem}[Declared mathematical global-boundary closure]
Inside the declared mathematical declared global-boundary synthesis grammar
\(\mathcal G_{\mathrm{global}}^{\mathrm{decl}}\), the OP1/OP4/OP5 closure
surface is closed:
\[
    \mathrm{Closure}_{\mathcal G_{\mathrm{global}}^{\mathrm{decl}}}
    (\mathrm{OP1},\mathrm{OP4},\mathrm{OP5}).
\]
Equivalently,
\[
    B_G=0
\]
inside the declared global-boundary synthesis grammar.
\end{theorem}

\begin{proof}
The grammar-exhaustion gate discharges grammar exhaustion.  The readout-invariance gate discharges
declared readout invariance.  The debt-free transport gate transports the boundary without
new debt.  The no-new-debt extension gate proves no-new-debt for the declared extension
buckets.  The functorial-stability gate proves functorial stability of the comparison tensor
\(\mathfrak C\) under the declared extension maps.  Therefore all five
global-boundary gates are passed.  Since no open gate remains and no new debt,
reopening, leakage, variance defect, or comparison-instability is recorded, the
boundary marker \(B_G\) is discharged inside the declared mathematical
declared global-boundary synthesis grammar.
\end{proof}

\begin{remark}[Scope guard]
The declared global-boundary theorem is a mathematical grammar-closure theorem.
It is not a theorem over every possible external readout, not a global physical
closure theorem, not a QFT/GR closure theorem, and not release-final status.
Those claims remain outside the declared grammar scope.
\end{remark}



% --- v4.5.0-r22 Source-Native Active-Core Rebuild Checkpoint ---

\clearpage
\section{Compressed Editorial and Release-Provenance Layer}
\label{sec:rc16-compressed-editorial-release-provenance}

\noindent\textbf{Compression status.}
This section replaces a long sequence of internal editorial, proof-read, glossary-intake,
metadata, release-freeze, and reader-facing phrase-normalization worksheets from the
pre-v4.7 provenance layer.  The removed material is preserved verbatim in the RC16 source
packet as \texttt{RC14\_uncompressed\_reference\_source.tex}.  No theorem, definition,
operator, DOI, or claim-scope object is changed by this compression.

\begin{theoremlet}[Editorial-provenance compression guard]
Let \(\mathcal P_{\rm edit}\) denote the compressed editorial/release-provenance worksheet
chain.  Replacing \(\mathcal P_{\rm edit}\) by this summary is release-safe for the main
monolith provided that the uncompressed source is retained in the provenance packet and
that no mathematical label from \(\mathcal P_{\rm edit}\) is externally referenced.  In the
RC16 audit, the external-reference count for removed labels is zero.  Hence the replacement
is a presentation compression only.
\end{theoremlet}

\noindent\textbf{Compressed content.}
The compressed worksheet chain recorded: active-core rebuild checks, release-candidate
audit planning, reader-facing global-scope disclaimers, DOI synchronization, release-freeze
audits, frontmatter repair, line-level proofreading, terminology sweeps, glossary/index
intake, glossary appendix placement decisions, and compact glossary draft routing.  All of
these are editorial provenance and release-hygiene layers.  They do not alter the proof
spine, the local grammar, the bridge obstruction analysis, the OP1/OP4 status, the OP5/C5
classification, or the v4.7.0 guarded effective/local claim scope.

\noindent\textbf{Preserved guards.}
The intended v4.7.0 Version DOI remains
\href{https://doi.org/10.5281/zenodo.20232911}{10.5281/zenodo.20232911}.  The previous
baseline DOI remains \href{https://doi.org/10.5281/zenodo.20117361}{10.5281/zenodo.20117361}
only as provenance.  No public release is implied by this compression layer.

